\documentclass[11pt]{article}

\usepackage[a4paper,margin=1in]{geometry}
\usepackage{amsmath,amssymb,amsfonts,amsthm,mathtools}
\usepackage{bm}
\usepackage{booktabs}
\usepackage{longtable}
\usepackage{graphicx}
\usepackage[hidelinks]{hyperref}

\newtheorem{theorem}{Theorem}
\newtheorem{proposition}{Proposition}
\newtheorem{lemma}{Lemma}
\newtheorem{corollary}{Corollary}
\theoremstyle{definition}
\newtheorem{definition}{Definition}
\theoremstyle{remark}
\newtheorem{remark}{Remark}

\newcommand{\SU}{\mathrm{SU}}
\newcommand{\U}{\mathrm{U}}
\newcommand{\SO}{\mathrm{SO}}
\newcommand{\AdS}{\mathrm{AdS}}
\newcommand{\CP}{\mathbb{CP}}
\newcommand{\C}{\mathbb{C}}
\newcommand{\Z}{\mathbb{Z}}
\newcommand{\R}{\mathbb{R}}
\newcommand{\dd}{\mathrm d}
\newcommand{\half}{\tfrac12}
\newcommand{\cA}{\mathcal A}
\newcommand{\cK}{\mathcal K}
\newcommand{\cH}{\mathcal H}
\newcommand{\cO}{\mathcal O}
\newcommand{\Xp}{X_+}
\newcommand{\Xm}{X_-}
\newcommand{\sinsq}{\sin^2\theta_W}

\title{\bf A D10 Boundary-Current Route To The de Vries Electroweak Operator}
\author{Codex technical version}
\date{25 May 2026}

\begin{document}
\maketitle

% CONTROL: Proof gates for this manuscript are recorded in
% codexVersion-proof-gates.md. Keep gate lists, completion audits, and task
% bookkeeping outside the typeset text.

\begin{abstract}
\noindent
We formulate a ten-dimensional boundary-current route to the branch operator
\[
  \cA(J)=
  \begin{pmatrix}
  0&\sqrt J\\
  \sqrt J&-J
  \end{pmatrix},
  \qquad J=j(j+1),
\]
whose positive eigenvalues at \(j=\half,1\) reproduce the on-shell weak
mixing angle.  The constructive layer is Type IIA on
\(M_4\times K_6\), with \(K_6=\CP^2\times\CP^1\) and a flag/current
refinement.  The active construction is entirely ten-dimensional.  The D9
endpoint is the \(Q\)-boundary sector
\((\SU(3)_c\times\U(1)_Q)/\Z_3\).  Under the Stueckelberg Schur assumption
with unit current metric, the projected D10 fluctuation operator gives the de
Vries pole relation.  An almost-commutative finite layer supplies the scalar
dictionary: the Higgs expectation is the finite sheet length
\(\ell_H=\sqrt2/v_{\rm EW}\), the \(Q\)-neutral vacuum vector gives the
\(W/Z/\gamma\) Hessian, and modern scalar extensions add a neutral singlet
Hessian block.  The route emphasized here is the underused \(Q\)-boundary
mechanism: the D8-like boundary
\(\Gamma_{\rm EW}=M_4\times\CP^2\times S^1_Q\), the RR/current contribution
to \(L_{2,{\rm vert}}^2\), and the boundary auxiliary field that generates the
negative diagonal.
\end{abstract}

\tableofcontents

\section{D10 Boundary-Current Setup}

The constructive object is the Type IIA compactification
\[
  M_4\times K_6,
  \qquad
  K_6=\CP^2\times\CP^1,
\]
with retained Type IIA \(C_1,F_2\) data and with a boundary/current sector
carrying the electroweak operator.
The electroweak endpoint is carried by a different circle, the diagonal
weak-hypercharge frame circle
\[
  S^1_Q\subset P_{T_3}\times P_Y,
  \qquad Q=T_3+Y.
\]
This \(S^1_Q\) cycle is the circle used by the localized
\(\Gamma_{\rm EW}\) boundary action.

The manuscript develops the following conditional implication.  If the D10
boundary/current sector gives a projected operator
\[
  \Pi_j\cK_{D10,*}\Pi_j
  =
  \mu_*^2
  \begin{pmatrix}
  0&\sqrt{j(j+1)}\\
  \sqrt{j(j+1)}&-j(j+1)
  \end{pmatrix},
  \qquad
  j=0,\half,1,
\]
then
\[
  M_\gamma^2=0,\qquad
  M_W^2=\mu_*^2\Xp(3/4),\qquad
  M_Z^2=\mu_*^2\Xp(2),
\]
where
\[
  \Xp(J)=\frac{-J+\sqrt{J^2+4J}}{2}.
\]
The pole weak angle is then
\[
  \sin^2\theta_{\rm pole}
  =
  1-\frac{\Xp(3/4)}{\Xp(2)}
  =
  0.223101\ldots .
\]
The localized boundary matrix defines \(L_Q\), and \(L_Q\) fixes the absolute
electromagnetic normalization of a concrete realization.

\section{Weak Sectors And The Schur Operator}

The weak \(\CP^1\) factor carries the Borel--Weil sectors
\[
  H^0(\CP^1,\cO(2j))\simeq V_j,\qquad j=0,\half,1.
\]
The affine-current filter supplies the same finite set at level \(k=2\),
since integrable \(\SU(2)_k\) sectors obey
\[
  j=0,\half,\ldots,\frac{k}{2}.
\]
For this current-algebra filter the conformal weights are
\[
  \Delta_j=\frac{j(j+1)}{k+2},
  \qquad k=2,
\]
with conversion table
\[
\begin{array}{c|c|c|c}
j&J=j(j+1)&\Delta_j=J/4&4\Delta_j\\
\hline
0&0&0&0\\
\half&3/4&3/16&3/4\\
1&2&1/2&2
\end{array}
\]
The string/interface Regge role is the mass template
\[
  M^2(n,j,\sigma)
  =
  \frac{n-a_0}{\alpha'}
  +\mu^2\lambda_\sigma(j(j+1))
  +\Delta M^2_{\rm compactification}.
\]
The string oscillator term carries the tower.  The internal Schur term
carries the finite electroweak branch.
The target-space operator uses finite zero modes.  The conformal-weight
\(L_0\) route supplies the control normalization.  Let
\[
  A_j=\sum_{a=1}^3T_a^{(j)}\otimes e_a,
\]
with \(e_a\) an orthonormal current frame.  Then
\[
  A_j^\dagger A_j
  =
  \sum_aT_a^{(j)}T_a^{(j)}
  =
  j(j+1)\,\bm 1_{V_j}
  =
  J\,\bm 1_{V_j}.
\]
The same \(A_j\) admits an Atiyah--Manton-type holonomy reading
\cite{AtiyahMantonSchroers2011}: a connection determines a Wilson line, and
the first variation of the holonomy inserts the Lie-algebra generator in the
same current frame.  The candidate realization list retained for comparison is
\[
  {\rm Dolbeault/Dirac},\quad
  {\rm boundary\ supercharge},\quad
  {\rm current\ zero\ mode},\quad
  {\rm brane\ localized\ auxiliary}.
\]
The boundary-current branch below supplies the coefficient calculation through
the Payen endpoint Hilbert space and the localized \(\Gamma_{\rm EW}\) action.

The algebraic boundary Hessian on \(V_j\oplus{\rm Im}\,A_j\) is
\[
  \cH_j=
  \begin{pmatrix}
  0&A_j^\dagger\\
  A_j&-A_jA_j^\dagger
  \end{pmatrix}.
\]
Projection to a normalized image vector gives
\[
  \cH_j\longrightarrow
  \begin{pmatrix}
  0&\sqrt J\\
  \sqrt J&-J
  \end{pmatrix}.
\]
Consequently the secular equation is
\[
  X^2+JX-J=0.
\]

If the retained current metric is
\[
  \langle e_a,e_b\rangle_{\rm cur}
  =
  \kappa_{\rm cur}\delta_{ab},
  \qquad c=\sqrt{\kappa_{\rm cur}},
\]
the projected block becomes
\[
  \begin{pmatrix}
  0&c\sqrt J\\
  c\sqrt J&-c^2J
  \end{pmatrix},
  \qquad
  X^2+c^2JX-c^2J=0.
\]
Thus the physical normalization is the statement
\[
  \kappa_{\rm cur}=1.
\]

% CONTROL: The derivation of the D10 current metric is gate 3 in
% codexVersion-proof-gates.md.

\section{The Stueckelberg Normalization}

The coefficient \(c\) can be isolated without committing to a complete
compactification.  On the unit \(S^3\), the scalar Casimir sector gives
\[
  \Delta_0=4J,
\]
and the co-exact one-form sector has the Weitzenbock shift
\[
  \Delta_1|_{\rm co\text{-}exact}=\Delta_0+1=4J+1=(2j+1)^2.
\]
The de Vries operator uses the algebraic Casimir \(A_j^\dagger A_j=J\).
The shifted geometric operator supplies a concrete failure channel whenever
the diagonal and off-diagonal are built from independently shifted operators.
Writing the local quadratic terms as
\[
  \alpha_{\rm mix}\,{\rm Re}\,\langle a,A_j\phi\rangle
  -
  \beta_{\rm diag}\,
  \langle\phi,A_j^\dagger A_j\phi\rangle,
\]
and normalizing the off-diagonal to \(\sqrt J\), the coefficient in the
projected block is
\[
  c=\frac{\beta_{\rm diag}}{\alpha_{\rm mix}^2}.
\]
The single-coupling condition is the equality
\[
  \beta_{\rm diag}=\alpha_{\rm mix}^2.
\]

\begin{proposition}[Stueckelberg Schur normalization]
\label{prop:stueckelberg}
Suppose the order-parameter scalar \(\phi\) in the spin-\(j\) sector has
vanishing fluxless mass and couples to the background through a single
Stueckelberg coupling
\[
  a^\dagger A_j\phi+\hbox{h.c.},
\]
where the same \(A_j\) satisfies \(A_j^\dagger A_j=J\).  The reduced
fluctuation block is the single-operator Schur completion
\[
  \begin{pmatrix}
  0&A_j^\dagger\\
  A_j&-A_jA_j^\dagger
  \end{pmatrix},
\]
and the coefficient in the deformed block is
\[
  c=1.
\]
\end{proposition}

\begin{proof}
The fluxless diagonal term of \(\phi\) vanishes by hypothesis.  The single
coupling gives the mixing operator \(A_j\).  The rank-one self-energy generated
by the same coupling is \(-A_jA_j^\dagger\), so projection to
\({\rm Im}\,A_j\) gives the block
\[
  \begin{pmatrix}
  0&\sqrt J\\
  \sqrt J&-J
  \end{pmatrix}.
\]
The ratio of the diagonal coefficient to the square of the off-diagonal
coefficient is therefore \(1\).
\end{proof}

The same one-coupling pattern has a standard higher-dimensional home in
gauge--Higgs unification \cite{Manton1979,Hosotani1983}.  Splitting a
higher-dimensional weak gauge field into \(A_\mu\) and internal components
\(A_a\), the internal components are Wilson-line or Hosotani scalar variables.
The cross term \(F_{\mu a}\) supplies the Stueckelberg coupling, and on the
\(\CP^1\) line-bundle modes the weak current zero mode acts as
\[
  J_0^a|_{V_j}=T_a^{(j)}.
\]
Thus the gauge--Higgs coupling uses the algebraic current operator
\[
  A_j=\sum_aT_a^{(j)}\otimes e_a,\qquad A_j^\dagger A_j=J.
\]
The shifted curl eigenvalue belongs to the cotangent-bundle control operator;
the gauge--Higgs coupling supplies the bare-Casimir normalization needed by the
Schur block.  The Schur completion itself is then the D10 boundary/current
statement.

The coefficient test is local.  Let
\[
  U_H(\xi)
  =
  {\cal P}\exp\left(
  i\mu_\Gamma\int_{\gamma_H}\xi^a e_aT_a^{(j)}
  \right),
\]
where \(e_a\) is the same current frame used in \(A_j\).  Then
\[
  \frac{\partial U_H}{\partial\xi^a}\bigg|_0
  =
  i\mu_\Gamma T_a^{(j)},
  \qquad
  \frac{\partial^2U_H}{\partial\xi^a\partial\xi^b}\bigg|_0
  =
  -\mu_\Gamma^2T_a^{(j)}T_b^{(j)}.
\]
Thus the Wilson-line expansion gives
\[
  \alpha_{\rm mix}=\mu_\Gamma,\qquad
  \beta_{\rm diag}=\mu_\Gamma^2,\qquad
  \frac{\beta_{\rm diag}}{\alpha_{\rm mix}^2}=1,
\]
after the common physical factor has been extracted.  For the stabilizer
cycle,
\[
  U_Q(\theta)=\exp(i\theta Q),\qquad Q\eta_0=0,
\]
the first two insertions on \(\eta_0\) vanish.  The Hosotani cycle must
therefore come from the retained transverse weak/current or Chan--Paton
boundary-complex sector.

The proposition separates two issues.  The ratio \(c=1\) follows from the
single Stueckelberg Schur structure.  The de Vries number additionally uses the
bare-Casimir argument \(J\) in \(\Xp(J)\).  The curvature-shifted alternatives
provide numerical controls:
\[
  J=\frac34,2
  \quad\leadsto\quad
  \sinsq=0.2231\ldots ,
\]
The independent shifted assignment gives
\[
  c_W=\frac{4(3/4)+1}{4(3/4)}=\frac43,\qquad
  c_Z=\frac{4(2)+1}{4(2)}=\frac98,
\]
and hence
\[
  \sin^2\theta_W|_{\rm shifted}=0.267\ldots .
\]

% CONTROL: The single-coupling hypothesis itself remains assigned to the
% auxiliary derivation list. Hosotani target: identify phi with a Wilson-line zero
% mode theta_H=int_gamma A_int, derive alpha_mix A_j and beta_diag A_j^dagger A_j
% from one localized D10 expansion, verify beta_diag=alpha_mix^2, holonomy
% periodicity, Q-neutral vacuum component, absence of an independent shifted
% operator, finite Wilson-line potential, and Connes sheet compatibility.
% Calculation sequence: select gamma_H; compute f_H^2 and mu_Gamma from the
% localized Gamma_EW action; expand U_H=P exp(i mu_Gamma int_gamma_H xi^a e_a
% T_a^(j)); canonically normalize alpha_mix=mu_Gamma/f_H and
% beta_diag=mu_Gamma^2/f_H^2; derive V_H(theta_H)=sum_m c_m cos(m w_H theta_H);
% check KO6/Poincare-duality finite-triple compatibility and survival of Q on
% eta_0. Pure S1_Q supplies electromagnetic stabilizer data and gives zero
% Schur insertion on eta_0.
% Hosotani local data in auxiliary notes:
% (gamma_H,omega_H,n_H^a,(f_H^2)_{ab},w_H), int_gamma_H omega_H=1,
% current-frame direction n_H^a e_a, localized kinetic matrix from DBI/YM plus
% endpoint terms, charge-lattice period theta_H~theta_H+2pi/w_H, stationary map
% |f_H theta_H,*|=v_EW/sqrt(2), and full doublet support before Q-neutral
% unitary-gauge projection to eta_0.

\section{The \texorpdfstring{\(Q\)}{Q}-Protected Boundary Cycle}

The weak doublet section is
\[
  \eta_0\in H^0(\CP^1,\cO(1)),
  \qquad
  T_3\eta_0=-\frac12\eta_0,
  \qquad
  Y\eta_0=\frac12\eta_0.
\]
Therefore
\[
  Q\eta_0=(T_3+Y)\eta_0=0,
\]
and the photon direction is protected by the same line sector that carries
the \(j=\half\) slot.

Let \(P_{T_3}\) and \(P_Y\) be the unit-frame bundles for the weak weight line
and the hypercharge line.  The neutral character is
\[
  \chi_N(u_T,u_Y)=u_T^{-1}u_Y.
\]
The \(Q\)-circle sector is
\[
  P_Q=\ker\chi_N\subset P_{T_3}\times P_Y.
\]
For the selected line data,
\[
  c_1(L_Y)=6h+7\eta,\qquad c_1(L_{T_3})=\eta,
\]
so the neutral line has
\[
  c_1(L_N)=c_1(L_Y)-c_1(L_{T_3})
  =
  6h+6\eta
  =
  6(h+\eta).
\]
The diagonal six-unit line class also decomposes the selected class as
\[
  6h+7\eta=6(h+\eta)+\eta,
\]
namely a diagonal six-unit line class plus one weak \(\cO(1)\) unit.

Using integer hypercharge \(y=6Y\), the ultraviolet diagonal element is
\[
  \zeta_6=(\omega_3,-1,e^{i\pi/3}),
  \qquad
  \omega_3=e^{2\pi i/3}.
\]
On a surviving state with color triality \(t\), weak weight \(T_3=m\), and
charge \(Q=m+Y\), its action becomes
\[
  \omega_3^t\exp(i2\pi m)\exp(i\pi y/3)
  =
  \omega_3^t\exp(i2\pi Q).
\]
Writing \(q_3=3Q\), the residual generator is
\[
  \zeta_3=(\omega_3,e^{2\pi i/3}),
\]
and the endpoint group is
\[
  G_{D9}=
  \frac{\SU(3)_c\times\U(1)_Q}{\Z_3}.
\]

% CONTROL: Detailed finite-triple audits live in
% connes-kk-higgs-vacuum-bridge.md,
% connes-finite-hessian-schur-comparison.md, and
% connes-determinant-one-full-lift-audit.md.

\section{Finite-Higgs Sheet And Scalar Vacuum}

The same \(Q\)-neutral weak line admits an almost-commutative
finite-geometric reading.  In the Connes--Lott electroweak construction and in
the later distance calculations for scalar fluctuations, the Higgs field is a
finite-direction connection and its expectation value is a metric datum
\cite{ConnesLott1990,Connes1989,MartinettiWulkenhaar2002}.  For a two-point
finite direction,
\[
  D_F(\Phi)=
  \begin{pmatrix}
  0&\Phi\\
  \Phi^\dagger&0
  \end{pmatrix},
\]
and for \(a=(\lambda_L,\lambda_R)\),
\[
  [D_F(\Phi),a]=
  \begin{pmatrix}
  0&(\lambda_R-\lambda_L)\Phi\\
  (\lambda_L-\lambda_R)\Phi^\dagger&0
  \end{pmatrix}.
\]
Consequently
\[
  \|[D_F(\Phi),a]\|
  =
  |\lambda_L-\lambda_R|\,\|\Phi\|,
  \qquad
  d_{\rm sheet}=\frac{1}{\|\Phi\|}.
\]
For the project scalar,
\[
  \Phi_0=v_F\eta_0,\qquad Q\eta_0=0,
\]
this gives
\[
  \ell_H=\frac{1}{|v_F|}
  =
  \frac{\sqrt2}{v_{\rm EW}}.
\]
Thus the electroweak vacuum amplitude is simultaneously a finite sheet length
and a D10 scalar coordinate.

The finite spectral datum underlying this calculation is
\[
  \cA_F^{\rm EW}=\C_r\oplus{\mathbb H}_w,\qquad
  \cH_F^{\rm EW}=H_r\oplus H_w,\qquad
  \Phi:H_r\longrightarrow H_w .
\]
The represented edge matched to the project is
\[
  H_r\xrightarrow{\ \Phi_0=v_F\eta_0\ }H_w,
  \qquad
  Y(H_r)=0,\qquad
  Y(H_w)=\half,\qquad
  Q\eta_0=0 .
\]
Inner fluctuation of \(D_F\) gives the finite connection \(\Phi\).  The
spectral distance gives \(\ell_H\).  The electroweak covariant derivative gives
the \(W/Z/\gamma\) Hessian.  The D10 boundary-current action compares this
finite Hessian with the Schur image Hessian through the common vector
\(\eta_0\).

The finite layer used here is the electroweak Glashow--Weinberg--Salam sector
of Connes--Lott, together with the scalar-fluctuation distance formula.  It
supplies the local \(\U(1)\times\SU(2)\) gauge data, a finite Higgs doublet,
and the auxiliary-field vocabulary.  The full Standard-Model finite triple
supplies the KO6, Poincare, orientability, and unimodularity controls used
below.

The local finite potential supplies the radial Higgs curvature.  With
\[
  V_F(\Phi)=\lambda_F
  \left(\langle\Phi,\Phi\rangle-v_F^2\right)^2,
  \qquad
  \Phi=v_F\eta_0+\delta\Phi,
\]
and
\[
  h=\sqrt2\,{\rm Re}\,\langle\eta_0,\delta\Phi\rangle,
\]
the quadratic term is
\[
  V_F^{(2)}
  =
  2\lambda_Fv_F^2h^2
  =
  \frac12 m_H^2h^2,
  \qquad
  m_H^2=4\lambda_Fv_F^2.
\]
Equivalently,
\[
  m_H^2=2\lambda_Fv_{\rm EW}^2,
  \qquad
  \lambda_F=\frac{m_H^2}{2v_{\rm EW}^2}.
\]
The value \(v_F\) fixes the finite sheet scale; the potential shape requires
the finite triple, spectral-action coefficients, Yukawa data, cutoff
normalization, and running.  The KK side supplies a separate effective
potential for compactification variables.  The common object is therefore
\[
  V_{\rm eff}(\rho,\lambda,v_F,\theta_H,\sigma,\ldots),
\]
with stationarity equations for the finite amplitude, the weak/base shape,
the overall compactification scale, and possible Wilson-line or singlet
directions.  A D10 electroweak phase additionally requires the threshold
window
\[
  \Lambda_{9\to10}(\rho_*,\lambda_*)
  <
  E_{\rm probe}
  <
  \Lambda_{{\rm UV},10}(\rho_*,\lambda_*) .
\]
Equivalently, writing the resolved weak-projective scale as
\[
  R_{\rm EW}(\ell)=\ell\,\widehat R_{\rm EW},
\]
the stationary vacuum must satisfy
\[
  R_{\rm EW}(\ell_*)>\Lambda_{{\rm UV},10}^{-1}.
\]
This inequality gives the D10 probe regime.  Collapse of the weak/projective
direction gives the D9 endpoint.
For the elementary quartic finite-sector model with \(\lambda_F>0\),
\[
  V_F=\lambda_F\left(\langle\Phi,\Phi\rangle-v_F^2\right)^2
  \quad\Longrightarrow\quad
  V_F\to+\infty\quad {\rm as}\quad |\Phi|\to\infty .
\]
In that phase diagram \(v_F=0\) can lose stability as the quadratic
coefficient changes sign, a finite \(v_0\) becomes the electroweak minimum,
and the \(v_F\to\infty\) boundary has infinite energy.  A D9 endpoint in the
present model therefore requires a coupled KK direction whose large-field
asymptotic flattens or turns the effective potential along a
compactification-collapse coordinate.
A factorized potential,
\[
  V_{\rm eff}=V_{\rm KK}(\rho,\lambda,\tau)+V_F(v_F,\sigma,\ldots),
\]
leaves the finite electroweak scale and the dimension window as separate
stationarity problems.  The D10 electroweak phase therefore requires mixed
Hessian entries such as
\[
  K_{\rho v}
  =
  \left.\frac{\partial^2V_{\rm eff}}{\partial\rho\,\partial v_F}\right|_*,
  \qquad
  K_{\lambda v}
  =
  \left.\frac{\partial^2V_{\rm eff}}{\partial\lambda\,\partial v_F}\right|_*,
\]
or a shared boundary/current scale
\[
  \mu_*^2=\mu_*^2(\rho_*,\lambda_*,v_F,\theta_H,\ldots).
\]
The reduced radial coefficient skeleton retained for the one-scale pass is
\[
  V_{\rm rad}(\ell)
  =
  -B_R\ell^{-2}
  +B_H\ell^{-6}
  +B_0\ell^6
  +B_2\ell^2
  +B_4\ell^{-2}
  +B_6\ell^{-6}
  +V_{\rm int}.
\]
The stationary calculation must pair this HKT scaling block with the finite
quartic \(V_F\) and the localized boundary term \(V_\Gamma\).
% CONTROL: Coefficient extraction and the stationarity proof remain in
% d10-one-scale-potential-gate.md and project-status-gates-audit.md.
The gauge Hessian uses the same vacuum vector:
\[
  K_{AB}^{\rm gauge}
  =
  v_F^2\,g_Ag_B\,
  \langle T_A\eta_0,T_B\eta_0\rangle.
\]
Together with \(T_3\eta_0=-\half\eta_0\), \(Y\eta_0=\half\eta_0\), this gives
\[
  M_W^2=\frac{g_2^2v_{\rm EW}^2}{4},
  \qquad
  M_Z^2=\frac{(g_2^2+g_Y^2)v_{\rm EW}^2}{4},
  \qquad
  M_\gamma^2=0.
\]
Using Weinberg's current-length normalization, let
\(C=2\pi\sqrt{16\pi G}\).  The KK current-length dictionary is
\[
  g_2=\frac{C}{L_2},\qquad
  g_Y=\frac{C}{L_Y},\qquad
  e=\frac{C}{L_Q},\qquad
  L_Q^2=L_2^2+L_Y^2 .
\]
Hence
\[
  \sin^2\theta_W=\frac{L_2^2}{L_Q^2},
  \qquad
  \tan^2\theta_W=\frac{L_2^2}{L_Y^2}.
\]
The same dictionary rewrites the vector masses as
\[
  M_W=\frac{ev_{\rm EW}}{2}\frac{L_Q}{L_2}
  =
  \frac{ev_{\rm EW}}{2\sin\theta_W},
  \qquad
  M_Z=\frac{ev_{\rm EW}}{2}\frac{L_Q^2}{L_2L_Y}
  =
  \frac{ev_{\rm EW}}{2\sin\theta_W\cos\theta_W}.
\]
Equivalently, with
\[
  G_{\rm EW}(\rho,\lambda,\ldots)
  =
  \frac{1}{L_2(\rho,\lambda,\ldots)^2}
  +
  \frac{1}{L_Y(\rho,\lambda,\ldots)^2},
\]
one has
\[
  M_Z^2=\frac{C^2}{4}v_{\rm EW}^2G_{\rm EW},
  \qquad
  \mu_*^2=
  \frac{C^2}{4\Xp(2)}v_{\rm EW}^2G_{\rm EW}.
\]
Thus the current-length dependence produces mixed scale sensitivity,
\[
  \frac{\partial^2M_Z^2}{\partial\rho\,\partial v_{\rm EW}}
  =
  \frac{C^2}{2}v_{\rm EW}\,\partial_\rho G_{\rm EW},
  \qquad
  \frac{\partial^2M_Z^2}{\partial\lambda\,\partial v_{\rm EW}}
  =
  \frac{C^2}{2}v_{\rm EW}\,\partial_\lambda G_{\rm EW}.
\]
The localized source term with the physical electroweak amplitude uses the
canonically normalized finite mode.  Let
\[
  I_0=
  \int_{\Gamma_{\rm EW,int}}\sqrt{\gamma_\Gamma}\,|\eta_0|^2,
  \qquad
  I_{AB}=
  \int_{\Gamma_{\rm EW,int}}\sqrt{\gamma_\Gamma}\,
  \langle T_A\eta_0,T_B\eta_0\rangle .
\]
For an unnormalized boundary-profile amplitude \(v_b\), the common
DBI/open-boundary coefficient \(Z_0\) gives
\[
  v_{\rm EW}^2=Z_0I_0v_b^2.
\]
The physical source term is therefore
\[
  V_{\Gamma}
  =
  \frac12v_{\rm EW}^2
  \Sigma_\Gamma^{AB}{\cal G}_{AB}(\rho,\lambda,\ldots),
  \qquad
  \Sigma_\Gamma^{AB}=\frac{I_{AB}}{I_0},
  \qquad Q\eta_0=0.
\]
For the selected normalized \(O(1)\) component,
\[
  T_3\eta_0=-\half\eta_0,\qquad
  Y\eta_0=\half\eta_0,
\]
so
\[
  \Sigma_{33}=\frac14,\qquad
  \Sigma_{YY}=\frac14,\qquad
  \Sigma_{3Y}=-\frac14.
\]
Hence
\[
  \Sigma_\Gamma(Q,Q)=0,\qquad Q=T_3+Y,
\]
and the photon direction is protected at the source tensor level.  The
charged weak entries satisfy
\[
  \|T_1\eta_0\|_\Gamma^2
  =
  \|T_2\eta_0\|_\Gamma^2
  =
  \frac14,
  \qquad
  {\rm Re}\,\langle T_1\eta_0,T_2\eta_0\rangle_\Gamma=0.
\]
With \(c_2=C/L_2\) and \(c_Y=C/L_Y\), the neutral source block is
\[
  K_{\rm neutral}
  =
  \frac{v_{\rm EW}^2}{4}
  \begin{pmatrix}
    c_2^2&-c_2c_Y\\
    -c_2c_Y&c_Y^2
  \end{pmatrix}
  =
  \frac{C^2v_{\rm EW}^2}{4}
  \begin{pmatrix}
    L_2^{-2}&-(L_2L_Y)^{-1}\\
    -(L_2L_Y)^{-1}&L_Y^{-2}
  \end{pmatrix}.
\]
Its null eigenvector is proportional to \((L_2,L_Y)\), giving
\[
  A_\mu=\frac{L_2}{L_Q}W^3_\mu+\frac{L_Y}{L_Q}B_\mu,
\]
and its nonzero eigenvalue is \(M_Z^2=(C^2v_{\rm EW}^2/4)G_{\rm EW}\).
% CONTROL: Boundary-frame neutral matrix propagation is audited in
% codexVersion-proof-gates.md and progressReport10.tex.
The localized boundary frame admits a general neutral length matrix
\[
  {\cal L}_{\Gamma}
  =
  \begin{pmatrix}
    L_2^2&\epsilon_\Gamma\\
    \epsilon_\Gamma&L_Y^2
  \end{pmatrix},
  \qquad
  \det{\cal L}_{\Gamma}=L_2^2L_Y^2-\epsilon_\Gamma^2 .
\]
The \(Q\)-circle length is
\[
  L_Q^2
  =
  (1,1){\cal L}_{\Gamma}\binom{1}{1}
  =
  L_2^2+L_Y^2+2\epsilon_\Gamma,
  \qquad
  \alpha_\Gamma=\frac{C^2}{4\pi L_Q^2}.
\]
The neutral source covector is \(s=(-1,1)\), hence the generalized neutral
coefficient is
\[
  G_Z^\Gamma
  =
  s{\cal L}_{\Gamma}^{-1}s^T
  =
  \frac{L_2^2+L_Y^2+2\epsilon_\Gamma}
       {L_2^2L_Y^2-\epsilon_\Gamma^2}
  =
  \frac{L_Q^2}{\det{\cal L}_{\Gamma}},
\]
and
\[
  M_Z^2=\frac{C^2v_{\rm EW}^2}{4}G_Z^\Gamma .
\]
The diagonal source metric appears at \(\epsilon_\Gamma=0\), recovering
\[
  L_Q^2=L_2^2+L_Y^2,\qquad
  G_Z^\Gamma=L_2^{-2}+L_Y^{-2}.
\]
Let
\[
  R_{\rm dV}=\frac{\Xp(3/4)}{\Xp(2)}.
\]
Since
\[
  M_W^2=\frac{C^2v_{\rm EW}^2}{4L_2^2},
\]
the boundary length matrix gives
\[
  \frac{M_W^2}{M_Z^2}
  =
  \frac{\det{\cal L}_{\Gamma}}{L_2^2L_Q^2}
  =
  \frac{L_2^2L_Y^2-\epsilon_\Gamma^2}
       {L_2^2(L_2^2+L_Y^2+2\epsilon_\Gamma)}.
\]
The de Vries pass equation for the neutral boundary matrix is
\[
  \epsilon_\Gamma^2
  +2R_{\rm dV}L_2^2\epsilon_\Gamma
  +R_{\rm dV}L_2^4
  +(R_{\rm dV}-1)L_2^2L_Y^2
  =
  0.
\]
With
\[
  u=\frac{L_2^2}{L_Y^2},
  \qquad
  \delta=\frac{\epsilon_\Gamma}{L_Y^2},
\]
this becomes
\[
  R_{\rm dV}
  =
  \frac{u-\delta^2}{u(u+1+2\delta)},
  \qquad
  \delta_\pm(u)
  =
  -R_{\rm dV}u
  \pm
  \sqrt{u(1-R_{\rm dV})(1-R_{\rm dV}u)}.
\]
The diagonal source metric sits at
\[
  \delta_+(u_{\rm dV})=0,
  \qquad
  u_{\rm dV}=\frac{1-R_{\rm dV}}{R_{\rm dV}}
  =
  \tan^2\theta_{\rm pole}.
\]
Since \({\cal L}_\Gamma\) is a kinetic matrix, the allowed region is
\[
  u>0,\qquad
  |\delta|<\sqrt u,\qquad
  u+1+2\delta>0.
\]
The de Vries branches are real for
\[
  0<u\le R_{\rm dV}^{-1}=1.2871691363429836\ldots .
\]
At \(u=u_{\rm dV}\) the second branch gives
\[
  \delta_-(u_{\rm dV})
  =
  -2(1-R_{\rm dV})
  =
  -0.4462026446017324\ldots .
\]
It also satisfies the kinetic inequalities.  The localized
\(\Gamma_{\rm EW}\) action is therefore characterized by the neutral branch
data, including \(L_Q\), together with the charged current length.
Equivalently, define the normalized neutral current correlation
\[
  \chi_\Gamma
  =
  \frac{\epsilon_\Gamma}{L_2L_Y}
  =
  \frac{\delta}{\sqrt u}.
\]
Then \(-1<\chi_\Gamma<1\), and the pass equation reads
\[
  R_{\rm dV}
  =
  \frac{1-\chi_\Gamma^2}
       {u+1+2\chi_\Gamma\sqrt u}.
\]
Solving for the correlation gives
\[
  \chi_\pm(u)
  =
  -R_{\rm dV}\sqrt u
  \pm
  \sqrt{(1-R_{\rm dV})(1-R_{\rm dV}u)}.
\]
At the diagonal source point,
\[
  \chi_+(u_{\rm dV})=0,
  \qquad
  \chi_-(u_{\rm dV})
  =
  -2\sqrt{R_{\rm dV}(1-R_{\rm dV})}
  =
  -0.8326514812056032\ldots .
\]
The neutral branch question is therefore a bounded current-correlation
calculation in the localized boundary action.
The same data factorize the matrix as
\[
  {\cal L}_\Gamma
  =
  L_Y^2
  \begin{pmatrix}
    u&\chi_\Gamma\sqrt u\\
    \chi_\Gamma\sqrt u&1
  \end{pmatrix}.
\]
Consequently
\[
  \det{\cal L}_\Gamma=L_Y^4u(1-\chi_\Gamma^2),
  \qquad
  L_Q^2=L_Y^2(u+1+2\chi_\Gamma\sqrt u),
\]
and
\[
  G_Z^\Gamma
  =
  \frac{u+1+2\chi_\Gamma\sqrt u}
       {L_Y^2u(1-\chi_\Gamma^2)},
  \qquad
  \alpha_\Gamma
  =
  \frac{C^2}
       {4\pi L_Y^2(u+1+2\chi_\Gamma\sqrt u)}.
\]
The localized action target is the triple
\[
  S_\Gamma
  \quad\leadsto\quad
  \left(L_Y^2,u,\chi_\Gamma\right),
\]
with
\[
  u=
  \frac{\int d\nu_\Gamma\,\|\theta_\Gamma^3\|^2}
       {\int d\nu_\Gamma\,\|\theta_\Gamma^Y\|^2},
  \qquad
  \chi_\Gamma=
  \frac{\int d\nu_\Gamma\,
  \gamma_\Gamma^{-1}(\theta_\Gamma^3,\theta_\Gamma^Y)}
  {L_2L_Y}.
\]
Writing \(y=\sqrt u\), the pass equation is equivalently
\[
  R_{\rm dV}y^2
  +2R_{\rm dV}\chi_\Gamma y
  +R_{\rm dV}+\chi_\Gamma^2-1
  =
  0,
\]
so
\[
  y_\pm(\chi_\Gamma)
  =
  -\chi_\Gamma
  \pm
  \sqrt{\frac{(1-R_{\rm dV})(1-\chi_\Gamma^2)}{R_{\rm dV}}},
  \qquad
  u_\pm(\chi_\Gamma)=y_\pm(\chi_\Gamma)^2 .
\]
At \(\chi_\Gamma=0\), the positive root is
\[
  u_+(0)=\frac{1-R_{\rm dV}}{R_{\rm dV}}.
\]
At \(u=u_{\rm dV}\), the correlated branch has
\[
  \chi_\Gamma=-2\sqrt{R_{\rm dV}(1-R_{\rm dV})}.
\]
Define
\[
  D_\Gamma=u+1+2\chi_\Gamma\sqrt u,
  \qquad
  \Lambda_\Gamma=\frac{C^2}{L_Y^2}.
\]
Then the factorized matrix gives
\[
  \alpha_\Gamma
  =
  \frac{\Lambda_\Gamma}{4\pi D_\Gamma},
  \qquad
  M_W^2
  =
  \frac{\Lambda_\Gamma v_{\rm EW}^2}{4u},
  \qquad
  M_Z^2
  =
  \frac{\Lambda_\Gamma v_{\rm EW}^2D_\Gamma}
       {4u(1-\chi_\Gamma^2)}.
\]
Eliminating \(\Lambda_\Gamma\) gives
\[
  \frac{M_Z^2}{\pi\alpha_\Gamma v_{\rm EW}^2}
  =
  \frac{D_\Gamma^2}{u(1-\chi_\Gamma^2)},
  \qquad
  \frac{M_W^2}{\pi\alpha_\Gamma v_{\rm EW}^2}
  =
  \frac{D_\Gamma}{u}.
\]
The diagonal branch reduces this to
\[
  \frac{M_Z^2}{\pi\alpha_\Gamma v_{\rm EW}^2}
  =
  \frac{1}{R_{\rm dV}(1-R_{\rm dV})},
  \qquad
  \frac{M_W^2}{\pi\alpha_\Gamma v_{\rm EW}^2}
  =
  \frac{1}{1-R_{\rm dV}}.
\]
The electromagnetic normalization is the localized computation of
\(\Lambda_\Gamma=C^2/L_Y^2\).
Given \((u,\chi_\Gamma)\) from the neutral current matrix, the three
observable channels infer
\[
  \Lambda_\Gamma^{(\alpha)}
  =
  4\pi\alpha_\Gamma D_\Gamma,
  \qquad
  \Lambda_\Gamma^{(W)}
  =
  \frac{4uM_W^2}{v_{\rm EW}^2},
  \qquad
  \Lambda_\Gamma^{(Z)}
  =
  \frac{4u(1-\chi_\Gamma^2)M_Z^2}
       {v_{\rm EW}^2D_\Gamma}.
\]
The D10 observable closure is
\[
  \Lambda_\Gamma^{(\alpha)}
  =
  \Lambda_\Gamma^{(W)}
  =
  \Lambda_\Gamma^{(Z)}.
\]
Equivalently,
\[
  4\pi\alpha_\Gamma D_\Gamma
  =
  \frac{4uM_W^2}{v_{\rm EW}^2}
  =
  \frac{4u(1-\chi_\Gamma^2)M_Z^2}
       {v_{\rm EW}^2D_\Gamma}.
\]
Define
\[
  S_W=
  \frac{M_W^2}{\pi\alpha_\Gamma v_{\rm EW}^2},
  \qquad
  S_Z=
  \frac{M_Z^2}{\pi\alpha_\Gamma v_{\rm EW}^2},
  \qquad
  R_{\rm obs}=\frac{S_W}{S_Z}.
\]
Then
\[
  D_\Gamma=S_Wu,
  \qquad
  1-\chi_\Gamma^2=R_{\rm obs}S_Wu.
\]
With \(\varsigma_\Gamma={\rm sgn}(\chi_\Gamma)\),
\[
  \chi_\Gamma
  =
  \varsigma_\Gamma
  \sqrt{1-R_{\rm obs}S_Wu}.
\]
The inverse electroweak reconstruction equation is
\[
  S_Wu
  =
  u+1+
  2\varsigma_\Gamma
  \sqrt{u(1-R_{\rm obs}S_Wu)}.
\]
For the exact de Vries pass, \(R_{\rm obs}=R_{\rm dV}\).  The common scale is
\[
  \Lambda_\Gamma=4\pi\alpha_\Gamma S_Wu.
\]
For canonical finite-mode normalization and frame-frozen
\(\Sigma_\Gamma^{AB}\),
\[
  \partial_\rho\log G_{\rm EW}=-\frac1\rho,
  \qquad
  \partial_\lambda\log G_{\rm EW}
  =
  -\frac{0.118746651783551}{\lambda}.
\]
Therefore
\[
  \frac{\partial^2M_Z^2}{\partial\rho\,\partial v_{\rm EW}}
  =
  -\frac{2M_Z^2}{v_{\rm EW}\rho},
  \qquad
  \frac{\partial^2M_Z^2}{\partial\lambda\,\partial v_{\rm EW}}
  =
  -0.237493303567102\,
  \frac{M_Z^2}{v_{\rm EW}\lambda}.
\]
The same dimensionless coefficients apply to
\(\mu_*^2=M_Z^2/\Xp(2)\).
For the charged weak block,
\[
  \frac{\partial^2M_W^2}{\partial\rho\,\partial v_{\rm EW}}
  =
  -\frac{2M_W^2}{v_{\rm EW}\rho},
  \qquad
  \frac{\partial^2M_W^2}{\partial\lambda\,\partial v_{\rm EW}}
  =
  -0.838811351483078\,
  \frac{M_W^2}{v_{\rm EW}\lambda}.
\]
The general mixed block is
\[
  K_{\rho v}^{\Gamma}
  =
  v_{\rm EW,*}
  \left[
  \Sigma_\Gamma^{AB}\partial_\rho{\cal G}_{AB}
  +(\partial_\rho\Sigma_\Gamma^{AB}){\cal G}_{AB}
  \right]_*,
\]
with the analogous \(\lambda\)-derivative.
In the one-scale window, \(L_A(\ell,\lambda)=\ell\widehat L_A(\lambda)\)
and \(\rho=c_\rho(\lambda)\ell^2\).  At fixed \(\lambda\),
\[
  {\cal G}^{AB}(\rho)
  =
  \frac{c_\rho}{\rho}\widehat{\cal G}^{AB},
  \qquad
  \partial_\rho{\cal G}^{AB}
  =
  -\frac{1}{\rho}{\cal G}^{AB}.
\]
For frame-frozen \(\Sigma_\Gamma^{AB}\), a positive coefficient gives
\[
  \partial_\rho V_\Gamma=-\frac{1}{\rho}V_\Gamma,
  \qquad
  K_{\rho v}^{\Gamma}
  =
  -\frac{1}{\rho_*}
  v_{\rm EW,*}\Sigma_\Gamma^{AB}{\cal G}_{AB}\big|_* .
\]
For an unnormalized boundary-profile amplitude, the DBI-scaled finite-section
coefficient carries the D8 source scaling
\[
  Z_\Gamma\propto\rho\,\tau^{-3}.
\]
At fixed \(\tau\), this cancels the radial scaling of
\({\cal G}^{AB}\).  With
\[
  \tau=\tau_*\left(\frac{\ell}{\ell_*}\right)^s,
  \qquad
  \rho=c_\rho\ell^2,
\]
one obtains
\[
  \partial_\rho\log(Z_\Gamma{\cal G}^{AB})
  =
  -\frac{3s}{2\rho}.
\]
In the canonical \(v_{\rm EW}\) convention, the common DBI factor is absorbed
into \(v_{\rm EW}\), leaving the normalized matrix \(\Sigma_\Gamma^{AB}\).
At fixed \(\rho\), the selected source metric also gives local
\(\lambda\)-derivatives.  For
\[
  u(\beta)=\frac{L_2^2}{L_Y^2}
  =
  3\frac{n^2q^2}{r^2}\frac{1+\beta}{\beta},
  \qquad
  \lambda(\beta)
  =
  \frac34\frac{p^2}{q^2}\beta(3-4\beta)^2 ,
\]
one obtains
\[
  \left.\frac{\partial\log u}{\partial\log\lambda}\right|_*
  =
  1.347631602696264 .
\]
Combining this with the weak length coefficient gives
\[
  \partial_\lambda\log L_2^2
  =
  \frac{0.419405675741539}{\lambda},
  \qquad
  \partial_\lambda\log L_Y^2
  =
  -\frac{0.928225926954725}{\lambda}.
\]
Consequently,
\[
  \partial_\lambda\log L_Q^2
  =
  -\frac{0.627566902996737}{\lambda},
  \qquad
  \partial_\lambda\log G_{\rm EW}
  =
  -\frac{0.118746651783551}{\lambda}.
\]
The total \(\lambda\)-block for the canonical finite-section source adds
\(\partial_\lambda\Sigma_\Gamma^{AB}\).  For the unnormalized
boundary-profile control pass, at fixed \(\rho\) and fixed \(\tau\), write
\[
  Z_\Gamma=Z_{\Gamma,0}z_\Gamma(\lambda).
\]
For a constant-profile boundary-cycle normalization on
\(\CP^2\times S^1_Q\),
\[
  z_{\Gamma,{\rm cp}}(\lambda)
  \propto
  c_\rho(\lambda)^{-5/2}\widehat L_Q(\lambda),
\]
and the selected source metric gives
\[
  \partial_\lambda\log z_{\Gamma,{\rm cp}}
  =
  -\frac{0.980450118165035}{\lambda}.
\]
Consequently,
\[
  \partial_\lambda\log(z_{\Gamma,{\rm cp}}G_{\rm EW})
  =
  -\frac{1.099196769948586}{\lambda}.
\]
A canonically normalized four-dimensional finite mode can absorb the
cycle-volume factor into \(v_{\rm EW}\).  In that convention
\[
  z_{\Gamma,{\rm unit}}(\lambda)=1,
  \qquad
  \partial_\lambda\log(z_{\Gamma,{\rm unit}}G_{\rm EW})
  =
  -\frac{0.118746651783551}{\lambda}.
\]
% CONTROL: derive Z_Gamma and the direct localized-current tensor signs in
% mixed-current-vacuum-hessian-gate.md.
Standard KK scaling gives a second D9 boundary at finite electroweak vacuum
value.  Since \(g_2\propto L_2^{-1}\),
\[
  L_2\to0,\qquad v_{\rm EW}=v_0<\infty
  \quad\Longrightarrow\quad
  M_W,M_Z\to\infty .
\]
With finite \(L_Y\),
\[
  e=\frac{g_2g_Y}{\sqrt{g_2^2+g_Y^2}}\to g_Y,
  \qquad
  \sin^2\theta_W\to0,
\]
so the surviving photon is the finite \(U(1)_Q\) current at the D9 boundary.
Indeed, with
\[
  A_\mu=\sin\theta_W\,W^3_\mu+\cos\theta_W\,B_\mu,
  \qquad
  Q=T_3+Y,
\]
the field direction tends to \(B_\mu\), and the charge generator remains
\(Q\) because
\[
  g_2\sin\theta_W=e,
  \qquad
  g_Y\cos\theta_W=e .
\]
The branch with \(L_2,L_Y\to0\) has \(e\to\infty\) and gives a singular
endpoint.  This finite-\(v\) route is controlled by the current metric,
and the large-\(v_F\) route is controlled by the scalar asymptotic.  As a
vacuum endpoint, finite \(v_{\rm EW}\) with \(M_W,M_Z\to\infty\) is a
pathological decoupling limit; its value here is diagnostic for the
\(Q\)-projection.

The de Vries operator enters through the boundary-current image channel.  For
\[
  A_j=\sum_aT_a^{(j)}\otimes e_a,\qquad A_j^\dagger A_j=J\bm1,
\]
the source/image Hessian is
\[
  \cH_j=
  \begin{pmatrix}
  0&\sqrt J\\
  \sqrt J&-J
  \end{pmatrix},
  \qquad
  \det\cH_j=-J.
\]
For \(j=\half\), the same neutral doublet gives \(J=3/4\).  The finite
potential, the gauge Hessian, and the Schur image Hessian therefore share the
vacuum vector \(\eta_0\) and occupy different quadratic sectors of the D10
interface.
The corresponding Schur mass scale obeys
\[
  \mu_*^2=\frac{M_Z^2}{X_+(2)}
  =
  \frac{g_2^2+g_Y^2}{4X_+(2)}v_{\rm EW}^2.
\]
The canonical electroweak source tensor fixes the physical \(j=1\) neutral
eigenvalue, so the dimensionful Schur scale is anchored by
\[
  \mu_*^2X_+(2)=
  M_Z^2
  =
  \frac{C^2v_{\rm EW}^2}{4}G_{\rm EW}.
\]
The coefficient split is
\[
  \hbox{scale:}\quad
  \mu_*^2=
  \frac{C^2v_{\rm EW}^2}{4X_+(2)}G_{\rm EW},
  \qquad
  \hbox{ratio:}\quad
  a=1,\quad \kappa_{\rm cur}=1.
\]
Anchoring the \(j=1\) channel converts the charged channel into the
weak-angle length equation.  The source block gives
\[
  \frac{M_W^2}{M_Z^2}
  =
  \frac{L_Y^2}{L_2^2+L_Y^2}
  =
  \frac{1}{1+u},
  \qquad
  u=\frac{L_2^2}{L_Y^2}.
\]
The Schur block gives
\[
  \frac{M_W^2}{M_Z^2}
  =
  \frac{X_+(3/4)}{X_+(2)}.
\]
Source/Schur agreement is therefore equivalent to
\[
  u=
  \frac{X_+(2)}{X_+(3/4)}-1
  =
  0.2871691363429836.
\]
Thus the source tensor supplies the common mass scale, and the Payen
Wilson/current-frame calculation supplies the Schur ratio.
% CONTROL: Local Gamma_EW action derivation of the common mu_*^2 coefficient
% and M_perp^2 gap lives in gamma-ew-boundary-quadratic-action-audit.md.
Modern spectral-action variants can add a real singlet scalar \(\sigma\)
coupled to the finite Higgs sector
\cite{ChamseddineConnes2012,ChamseddineConnesVanSuijlekom2013a,
ChamseddineConnesVanSuijlekom2013b,Martinetti2015,
FilaciMartinettiPesco2020}.  A local model is
\[
  V_{h\sigma}
  =
  \lambda_F(\langle\Phi,\Phi\rangle-v_F^2)^2
  +\frac12M_{\sigma,0}^2\sigma^2
  +a_\sigma\sigma(\langle\Phi,\Phi\rangle-v_F^2).
\]
In the canonical basis \((h,\sigma)\), the Hessian is
\[
  K_{h\sigma}=
  \begin{pmatrix}
  2\lambda_Fv_{\rm EW}^2&a_\sigma v_{\rm EW}\\
  a_\sigma v_{\rm EW}&M_{\sigma,0}^2
  \end{pmatrix}.
\]
Writing \(M_{h,0}^2=2\lambda_Fv_{\rm EW}^2\), the scalar eigenvalues are
\[
  m_\pm^2=
  \frac12\left(
  M_{h,0}^2+M_{\sigma,0}^2
  \pm
  \sqrt{(M_{h,0}^2-M_{\sigma,0}^2)^2
  +4a_\sigma^2v_{\rm EW}^2}
  \right),
\]
with mixing angle
\[
  \tan 2\alpha
  =
  \frac{2a_\sigma v_{\rm EW}}
       {M_{\sigma,0}^2-M_{h,0}^2}.
\]
Since \(Q\sigma=0\), this mixing leaves the photon-protection equation
unchanged and enlarges the scalar vacuum system.

The finite-triple control layer uses the Standard-Model KO6 structure
\cite{ChamseddineConnesMarcolli2007,Stephan2009,DAndreaDabrowski2016}.  A
local four-summand lift is
\[
  A_{\rm lift}=
  M_3(\C)_c\oplus{\mathbb H}_w\oplus\C_r\oplus\C_q,
\]
with project edge
\[
  H_{rq}^-\longrightarrow H_{wq}^+.
\]
The first-order factors are
\[
  (a_w-a_r)(b_q-b_q)=0,
  \qquad
  (a_q-a_q)(b_w-b_r)=0,
\]
and the local charge equations are
\[
  Y(H_{rq})=0,\qquad Y(H_{wq})=\half.
\]
With \(y_w=0\), this fixes
\[
  y_r=y_q=-\half.
\]
The project edge contribution to the four-summand intersection form has
\[
  \det(\cap_\Phi)=0.
\]
The full represented Krajewski matrix supplies the Standard-Model Poincare
form, determinant-one hypercharge lift, and fermion Hilbert-space data.
Modern finite-geometry checks add Hodge-duality and spectral-torsion
constraints involving Yukawa and Majorana data
\cite{DabrowskiSitarz2019,DabrowskiMukhopadhyayPozar2025}.

\section{The D8-Like Boundary And Schur Auxiliary}

The common boundary interface for \(Q\) protection and the Schur auxiliary
channel is
\[
  \Gamma_{\rm EW}=M_4\times\CP^2\times S^1_Q.
\]
The localized action on this interface is split as
\[
  S_\Gamma
  =
  S_{\rm DBI}^{\Gamma}
  +S_{\rm CS}^{\Gamma}
  +S_{\rm Payen}^{\Gamma}
  +S_{\rm aux}^{\Gamma}.
\]
The D8-like DBI term is
\[
  S_{\rm DBI}^{\Gamma}
  =
  -T_8\int_{\Gamma_{\rm EW}}d^9\xi\,
  e^{-\phi}
  \sqrt{-\det\!\left(P[g+B]+2\pi\alpha'F_Q\right)} .
\]
Its quadratic finite/current density is
\[
  d\nu_\Gamma
  =
  T_8(2\pi\alpha')^2e^{-\phi}\sqrt{\gamma_\Gamma}\,
  {\cal W}_\Gamma\,d^5\sigma ,
\]
where \({\cal W}_\Gamma\) collects the orientifold-image, frame, and
reduction factors.
It is D8-like in Type IIA.  The universal localized-source scaling gives
\[
  V_{Dp}\propto \rho^{(p-6)/2}\tau^{-3},
\]
so for \(p=8\),
\[
  V_{D8}\propto \rho\,\tau^{-3}.
\]
With one-scale variables
\[
  \rho=c_\rho\ell^2,\qquad
  \tau=\tau_*\left(\frac{\ell}{\ell_*}\right)^s,
\]
this contribution scales as
\[
  V_{D8}\propto x^{2-3s},
  \qquad x=\frac{\ell}{\ell_*}.
\]
The Romans-mass jump takes the distributional form
\[
  \dd F_0=N_8\,\delta_{\Gamma_{\rm EW}}.
\]

The Wess--Zumino skeleton is
\[
  S_{\rm WZ}
  =
  \mu_8\int_{\Gamma_{\rm EW}}
  C\wedge{\rm ch}(E_Q)
  \sqrt{\frac{\widehat A(T\Gamma_{\rm EW})}
             {\widehat A(N\Gamma_{\rm EW})}}.
\]
On the \(\CP^2\) boundary the spin\(^c\) shift gives
\[
  \frac{F_Q}{2\pi}=6h_Q+\frac{x}{2},
  \qquad x=(2m+1)h_Q.
\]
For the minimal lift,
\[
  \frac{F_Q}{2\pi}=\frac{13}{2}h_Q,
  \qquad
  \frac{2F_Q}{2\pi}=13h_Q.
\]
The boundary therefore carries the determinant signal \(P+Q=13\).  In a
realization with weak period \(Q=7\), the internal consistency check reads
\[
  r=Q(P+Q)=7\cdot13=91.
\]
This relation is used inside a specified realization as an internal
consistency condition.

The source-backed boundary placement is Payen's Chan--Paton action for a
group-valued endpoint field \(g(\tau)\) \cite{Payen2007}:
\[
  S_{\rm CP}
  =
  \int_{\partial\Sigma}d\tau\,
  \kappa\!\left(
  -ig^{-1}(\partial_\tau+A_\tau)g
  \right).
\]
Canonical quantization gives a finite-dimensional representation \(R\) of the
endpoint group, and the boundary path integral gives the Wilson loop
\[
  {\rm Tr}\,{\cal T}_{\partial\Sigma}
  \exp\left(
  i\int_{\partial\Sigma}d\tau\,R(A_\tau)
  \right).
\]
For \(G=\SU(2)\), \(R=V_j\), and
\[
  R_j(A_\tau)=A_\tau^aT_a^{(j)},
\]
the Wilson expansion supplies the linear charge insertion \(T_a^{(j)}\) and
the bilinear insertion \(T_a^{(j)}T_b^{(j)}\).
With a boundary fluctuation
\[
  R_j(A_\tau)=\mu\,\xi^aT_a^{(j)},
\]
the expansion has one tensor normalization:
\[
  \delta W_j|_{\xi=0}\sim i\mu\,T_a^{(j)}\delta\xi^a,
  \qquad
  \delta^2W_j|_{\xi=0}\sim
  -\mu^2T_a^{(j)}T_b^{(j)}\delta\xi^a\delta\xi^b.
\]
The representation product is therefore fixed once the boundary fluctuation
scale \(\mu\) is fixed.  Equivalently, a rescaling \(A_j\mapsto aA_j\)
sends the image bilinear to \(a^2A_jA_j^\dagger\).  The boundary exponential
therefore gives the tensor relation \(b=a^2\).  After the common physical
scale \(\mu_*^2\) is factored out, the orthonormal current frame gives the
canonical dimensionless coefficient \(a=1\).
The coefficient \(\mu\) belongs to the boundary fluctuation and the retained
current frame.  Representation dependence enters through \(T_a^{(j)}\) and
its Casimir \(j(j+1)\).  Thus a single \(\Gamma_{\rm EW}\) boundary sector
uses one worldvolume scale \(\mu_\Gamma^2\) for \(j=\half\) and \(j=1\).

The same boundary can carry the auxiliary field \(y\in V_j\).  Before the
\(Q\) projection, the weak Chan--Paton line sectors are \(\cO(0)\) and
\(\cO(2j)\).  Their mixed boundary-changing zero-mode space is
\[
  H^0(\CP^1,{\rm Hom}(\cO(0),\cO(2j)))
  =
  H^0(\CP^1,\cO(2j))
  =
  V_j.
\]
Thus the representation space for the auxiliary is the same Borel--Weil
sector used by the current map.
The D-brane charge formula uses a Chan--Paton bundle with connection on the
wrapped worldvolume \cite{MinasianMoore1997}, and Payen's boundary action
quantizes the endpoint degree into the representation \(V_j\)
\cite{Payen2007}.  For the two-line weak interface set
\[
  E_0=\cO(0),\qquad E_j=\cO(2j),\qquad
  \mathcal E_j={\rm Hom}(E_0,E_j)\simeq \cO(2j).
\]
The Hermitian metrics on \(E_0\) and \(E_j\) induce one Hermitian metric on
\(\mathcal E_j\), hence on
\[
  V_j=H^0(\CP^1,\mathcal E_j).
\]
Thus
\[
  y,\quad t,\quad A_j^\dagger n
  \in V_j
\]
are measured by the same boundary-complex norm.  The physical coefficient is
the common worldvolume normalization multiplying this norm.
For the neutral anchor \(j=1\), this boundary-complex metric is tied to the
finite weak doublet metric by the line-square map.  Let \(L=\cO(1)\).  Then
\[
  V_{\half}=H^0(\CP^1,L),
  \qquad
  V_1=H^0(\CP^1,L^{\otimes2})
  \simeq {\rm Sym}^2 V_{\half}.
\]
If the \(Q\)-neutral section \(\eta_0\in V_{\half}\) has unit norm, the
induced tensor-square metric gives
\[
  \|\eta_0\odot\eta_0\|_{V_1}
  =
  \|\eta_0\|_{V_{\half}}^2
  =
  1.
\]
Thus the \(j=1\) Schur scale anchor can use the same finite-section
normalization that defines \(v_{\rm EW}^2=Z_0I_0v_b^2\).
Writing \(Z_{\mathcal E,1}\) for the localized line-square boundary-complex
norm, the source-matched coefficient is
\[
  \mu_\Gamma^2
  =
  \frac{C^2}{4X_+(2)}
  Z_{\mathcal E,1}v_b^2\,G_{\rm EW},
  \qquad
  Z_{\mathcal E,1}=Z_0I_0.
\]
In a standard spin basis, the \(Q\)-neutral choice
\[
  \eta_0=|{\textstyle -\frac12}\rangle\in V_{\half}
\]
has tensor square
\[
  \eta_0\odot\eta_0=|-1\rangle\in V_1.
\]
Its charges are
\[
  T_3(\eta_0\odot\eta_0)=-(\eta_0\odot\eta_0),
  \qquad
  Y(\eta_0\odot\eta_0)=\eta_0\odot\eta_0,
\]
so
\[
  Q(\eta_0\odot\eta_0)=0.
\]
Consequently \(\eta_0\odot\eta_0\) is basic for the \(S^1_Q\) boundary action
and descends to the \(\Gamma_{\rm EW}\) boundary.
The spin-one Casimir then gives
\[
  \sum_aT_a^{(1)}T_a^{(1)}|-1\rangle=2|-1\rangle,
\]
so the neutral Schur anchor carries \(J=2\) from the symmetric square of the
same finite section that supplies the electroweak source tensor.

The scale equality is now a zero-mode normalization test.  In the Hilbert
normalization of the boundary zero modes,
\[
  \|\widehat\eta_0\|_{V_{\half}}=1,
  \qquad
  \|\widehat\eta_0\odot\widehat\eta_0\|_{{\rm Sym}^2V_{\half}}=1,
\]
so the induced neutral-anchor coefficient is
\[
  Z_{\mathcal E,1}^{\rm zm}=Z_0I_0,
  \qquad
  \chi_\Gamma^{\rm zm}=1.
\]
A pointwise tensor-profile reading gives a distinct diagnostic.  With
\[
  I_2=
  \int_{\Gamma_{\rm EW,int}}
  \sqrt{\gamma_\Gamma}\,|\eta_0|^4 ,
\]
the same local density gives
\[
  Z_{\mathcal E,1}^{\rm pt}=Z_0I_2,
  \qquad
  \chi_\Gamma^{\rm pt}=\frac{I_2}{I_0}.
\]
Thus source matching selects the zero-mode Hilbert tensor metric, or a
localized profile identity \(I_2=I_0\).
The Payen boundary construction supports the zero-mode branch: the endpoint
variable quantizes to a finite-dimensional Hilbert space carrying \(R\), and
the boundary path integral is the Wilson trace \cite{Payen2007}.  With
\[
  R=V_j=H^0(\CP^1,\cO(2j)),
\]
the line-square anchor is a tensor inside the endpoint Hilbert representation.
The pointwise \(I_2/I_0\) branch is a separate DBI profile assumption for a
composite line-square field.
Payen's non-oriented extension adds an orientifold compatibility check:
the endpoint representation must be equivalent to its dual
\cite{Payen2007}.  For \(\SU(2)\),
\[
  V_j^*\simeq V_j,
\]
with integer \(j\) real and half-integer \(j\) pseudoreal.  Hence the retained
\(j=\half\) doublet and \(j=1\) triplet both pass the D8/O8 endpoint reality
test.
Endpoint reality is a survival test.  The finite retained set
\[
  j=0,\half,1
\]
still comes from the line-bundle or level-two selection.

\begin{proposition}[Zero-mode scale pass]
Assume that the localized \(\Gamma_{\rm EW}\) boundary action uses the Payen
endpoint Hilbert representation \(R=V_j\), tensors the \(Q\)-neutral line
inside \(V_{\half}\), and multiplies the finite source and neutral Schur
anchor by the same D8/open-boundary density.  Then
\[
  Z_{\mathcal E,1}=Z_0I_0,
  \qquad
  \chi_\Gamma=1,
\]
and
\[
  \mu_\Gamma^2
  =
  \frac{C^2}{4\Xp(2)}Z_0I_0v_b^2G_{\rm EW}
  =
  \mu_*^2 .
\]
\end{proposition}

\begin{proof}
The zero-mode tensor metric gives
\(\|\widehat\eta_0\odot\widehat\eta_0\|=1\).  The shared
D8/open-boundary density gives \(Z_{\mathcal E,1}=Z_0I_0\).  Substitution in
\(\chi_\Gamma=Z_{\mathcal E,1}/(Z_0I_0)\) gives \(\chi_\Gamma=1\), and the
definition of \(\mu_\Gamma^2\) gives the stated scale.
\end{proof}

The shared-density form can be written explicitly as
\[
  d\nu_\Gamma
  =
  T_8(2\pi\alpha')^2e^{-\phi}\sqrt{\gamma_\Gamma}\,
  {\cal W}_\Gamma\,d^5\sigma ,
\]
where \({\cal W}_\Gamma\) collects the frame, orientifold-image, and
reduction factors of the localized action.  The finite source coefficient is
\[
  Z_0I_0
  =
  \int_{\Gamma_{\rm EW,int}}d\nu_\Gamma\,|\eta_0|^2 .
\]
For the endpoint Hilbert embedding,
\[
  Z_{\mathcal E,1}^{\rm zm}
  =
  \int_{\Gamma_{\rm EW,int}}d\nu_\Gamma\,|\eta_0|^2
  \|\widehat\eta_0\odot\widehat\eta_0\|^2
  =
  Z_0I_0 .
\]
Thus the coefficient datum is the same \(d\nu_\Gamma\) in the finite source
term and in the Payen endpoint Hilbert action.
Consequently, for any modulus \(m\) whose dependence in both coefficients is
carried by the same \(d\nu_\Gamma\),
\[
  \partial_m\log Z_{\mathcal E,1}^{\rm zm}
  =
  \partial_m\log(Z_0I_0),
  \qquad
  \partial_m\log\chi_\Gamma^{\rm zm}=0.
\]
Frame, orientifold-image, and reduction factors affect the common scale
\(Z_0I_0\) and cancel from the scale-pass ratio.
At fixed shape and finite profile, the D8-like density carries the source
scaling
\[
  Z_0I_0\propto \rho\,\tau^{-3},
  \qquad
  G_{\rm EW}\propto \rho^{-1}.
\]
Hence the Schur scale has
\[
  \mu_\Gamma^2\propto Z_0I_0G_{\rm EW}\propto \tau^{-3},
  \qquad
  \partial_\rho\log(Z_0I_0G_{\rm EW})\big|_{\tau,\lambda}=0.
\]
Along \(\rho=c_\rho\ell^2\) and
\(\tau=\tau_*(\ell/\ell_*)^s\),
\[
  \partial_{\log\ell}\log\mu_\Gamma^2=-3s.
\]

For \(J=j(j+1)>0\), the image projector is fixed by the Casimir identity:
\[
  A_j^\dagger A_j=J\,\bm1_{V_j},
  \qquad
  P_{\rm im}^{(j)}=\frac{1}{J}A_jA_j^\dagger .
\]
The current codomain decomposes under the diagonal \(\SU(2)\) action as
\[
  V_j\otimes V_1=V_{j+1}\oplus V_j\oplus V_{j-1},
\]
with the absent negative-spin summand omitted.  The map \(A_j\) is an
equivariant embedding of \(V_j\) into the middle summand.  Thus
\[
  n_{\rm im}=P_{\rm im}^{(j)}n,
  \qquad
  n_\perp=(1-P_{\rm im}^{(j)})n
\]
gives a canonical image/orthogonal split.  For \(j=0\), \(A_0=0\) gives the
photon slot.
Equivalently, with
\[
  {\cal B}_j=\sum_aT_a^{(j)}\otimes S_a^{(1)}
\]
on \(V_j\otimes V_1\), the diagonal Casimir gives eigenvalues
\[
  j,\quad -1,\quad -(j+1)
\]
on \(V_{j+1}\), \(V_j\), and \(V_{j-1}\), respectively.  Hence
\[
  P_{\rm im}^{(j)}
  =
  -\frac{({\cal B}_j-j)({\cal B}_j+j+1)}{j(j+1)}.
\]
The complementary projector is
\[
  P_{\perp}^{(j)}
  =
  1-P_{\rm im}^{(j)}
  =
  \frac{{\cal B}_j({\cal B}_j+1)}{j(j+1)}.
\]
It has eigenvalue zero on the image summand and eigenvalue one on the
orthogonal summands.  Thus \(M_\perp^2\langle n,P_\perp^{(j)}n\rangle\) is a
canonical local gap term for \(M_\perp^2>0\).
Equivalently, let \(\Delta_{\Gamma,j}\) be the self-adjoint Laplacian of the
localized boundary complex whose zero modes are the image summand.  Compactness
of the internal boundary image gives
\[
  \lambda_{\perp,j}
  =
  \inf_{\substack{n\in P_\perp^{(j)}(V_j\otimes V_1)\\ \|n\|=1}}
  \langle n,\Delta_{\Gamma,j}n\rangle
  >
  0.
\]
The complement gap can then be written as \(M_\perp^2=\lambda_{\perp,j}\)
up to the common boundary kinetic normalization.
% CONTROL: Complement-kernel and density audit lives in
% codexVersion-proof-gates.md.
With the shared density, the finite projector acts fibrewise:
\[
  {\cal H}_{\Gamma,j}
  =
  L^2(\Gamma_{\rm EW,int},d\nu_\Gamma)
  \otimes (V_j\otimes V_1),
\]
\[
  \langle n,P_\perp^{(j)}n\rangle_\Gamma
  =
  \int_{\Gamma_{\rm EW,int}}d\nu_\Gamma\,
  \langle n(\sigma),P_\perp^{(j)}n(\sigma)\rangle_{V_j\otimes V_1}.
\]
All frame, orientifold-image, and reduction factors stay in \(d\nu_\Gamma\).
If
\[
  \ker\Delta_{\Gamma,j}={\rm Im}\,A_j,
\]
then
\[
  M_{\perp,j}^2
  =
  \inf_{\substack{
  n\in P_\perp^{(j)}{\cal H}_{\Gamma,j}\\
  \|n\|_\Gamma=1}}
  \langle n,\Delta_{\Gamma,j}n\rangle_\Gamma
  >
  0.
\]
The residual dimensionless number is
\[
  r_{\perp,j}=\frac{M_{\perp,j}^2}{\mu_\Gamma^2}.
\]
It leaves the image-block de Vries eigenvalue unchanged and assigns the
complement datum to the localized boundary complex.

Let
\[
  t\in V_j,\qquad n=n_{\rm im}+n_\perp,\qquad
  n_{\rm im}\in{\rm Im}\,A_j,\qquad y\in V_j.
\]
The dimensionful local auxiliary form is
\[
  S_{\Gamma,j}^{(2)}
  =
  \mu_\Gamma^2
  \left[
  \|t\|^2+\|y\|^2+
  \langle y,t-A_j^\dagger n\rangle+
  \langle t-A_j^\dagger n,y\rangle
  \right]
  +M_\perp^2\langle n,P_\perp^{(j)}n\rangle
\]
eliminates \(y\) and leaves
\[
  \mu_\Gamma^2
  \left[
  \langle t,A_j^\dagger n\rangle+
  \langle n,A_jt\rangle-
  \langle n,A_jA_j^\dagger n\rangle
  \right]
  +M_\perp^2\langle n,P_\perp^{(j)}n\rangle .
\]
Equivalently, setting \(s=t-A_j^\dagger n\), the \(y\)-dependent part is
\[
  \mu_\Gamma^2
  \left[
  \|t\|^2+\|y\|^2+\langle y,s\rangle+\langle s,y\rangle
  \right]
  =
  \mu_\Gamma^2
  \left[
  \|y+s\|^2+\|t\|^2-\|s\|^2
  \right].
\]
Thus the equal-coefficient test is local: the D8 boundary complex provides
the same Hermitian metric for \(y\), for \(t-A_j^\dagger n\), and for
the endpoint field \(t\).
At the selected stationary point, the source-to-Schur scale equation is
\[
  \mu_\Gamma^2=\mu_*^2
  =
  \frac{C^2v_{\rm EW}^2}{4X_+(2)}G_{\rm EW}.
\]
Since \(X_+(2)=\sqrt3-1\), equivalently
\[
  \mu_\Gamma^2
  =
  \frac{\sqrt3+1}{2}\,M_Z^2.
\]
In boundary-profile variables, \(v_{\rm EW}^2=Z_0I_0v_b^2\) gives
\[
  \mu_\Gamma^2
  =
  \frac{C^2}{4X_+(2)}
  Z_0I_0v_b^2\,G_{\rm EW}.
\]
% CONTROL: Auxiliary coefficient pass scalar:
% chi_Gamma = mu_Gamma^2 X_+(2)/M_Z^2
% = 4 X_+(2) mu_Gamma^2 /(C^2 Z_0 I_0 v_b^2 G_EW).
This boundary Schur term supplies the negative diagonal in the projected
block on \(V_j\oplus{\rm Im}\,A_j\).  The orthogonal image complement is
gapped by \(M_\perp^2\).

% CONTROL: The concrete derivation of y, its sign, and its coefficient lives
% in codexVersion-proof-gates.md and the D10 boundary notes.

\section{RR/Current Boundary Length}

The weak coupling normalization follows Weinberg's rms-circumference rule,
\[
  g^2=\frac{4\pi^2\kappa^2}{L^2},
  \qquad
  \kappa^2=16\pi G.
\]
The D10 current-frame normalization is the retained weak length from the
localized current frame on \(\Gamma_{\rm EW}\).  Write
\[
  \Pi_{\rm EW}\iota_\Gamma^*C_1
  =
  q_\Gamma\,\theta_\Gamma^3,
  \qquad
  \Pi_{\rm EW}\iota_\Gamma^*F_2
  =
  q_\Gamma\,d_\Gamma\theta_\Gamma^3 .
\]
The \(\SU(2)\)-equivariant completion gives the triplet
\(\theta_\Gamma^a\).  The boundary-current length matrix is
\[
  K_{ab}^\Gamma
  =
  \int_{\Gamma_{\rm EW,int}}d\nu_\Gamma\,
  \gamma_\Gamma^{-1}(\theta_\Gamma^a,\theta_\Gamma^b).
\]
The unit current metric condition is the D10 equation
\[
  K_{ab}^\Gamma=L_{2,\Gamma}^2\delta_{ab}.
\]
A useful split is
\[
  L_{2,\Gamma}^2=L_{2,{\rm horiz}}^2+L_{2,{\rm vert}}^2,
\]
with the horizontal term supplied by the \(\CP^1\) weak orbit and the vertical
term supplied by the localized \(C_1,F_2\) current channel.  On a round
\(\CP^1\simeq S^2\) weak factor of radius \(R_w\),
\[
  K_a=x\times e_a,
  \qquad
  |K_a|^2=R_w^2-x_a^2,
\]
and the sphere average gives
\[
  \frac{1}{{\rm Vol}(\CP^1)}
  \int_{\CP^1}\sqrt g\,g(K_a,K_b)
  =
  \frac{2}{3}R_w^2\delta_{ab}.
\]
Thus
\[
  L_{2,{\rm horiz}}^2=\frac23\ell^2\lambda R_0^2.
\]
The vertical term is the localized RR/current integral
\[
  L_{2,{\rm vert}}^2\delta_{ab}
  =
  \int_{\Gamma_{\rm EW,int}}d\nu_\Gamma\,
  \gamma_\Gamma^{-1}(\theta_{\Gamma,{\rm vert}}^a,
  \theta_{\Gamma,{\rm vert}}^b).
\]
Equivalently, let \(\widetilde K_a\) be the retained boundary-current lift
whose horizontal part is the \(\CP^1\) rotation field \(K_a\), with vertical
part supplied by the localized \(C_1,F_2\) current channel.  The averaged
metric condition is
\[
  {\rm Av}_\Gamma\,g(\widetilde K_a,\widetilde K_b)
  =
  \left(L_{2,{\rm horiz}}^2+L_{2,{\rm vert}}^2\right)\delta_{ab}
  =
  L_{2,\Gamma}^2\delta_{ab}.
\]
Thus the current frame used in the Schur operator is
\[
  e_a=\frac{\widetilde K_a}{L_{2,\Gamma}},
  \qquad
  A_j=\sum_aT_a^{(j)}\otimes e_a,
  \qquad
  A_j^\dagger A_j=j(j+1)\bm1 .
\]
The same density \(d\nu_\Gamma\) appears in the finite source coefficient
and in the Payen endpoint Hilbert coefficient, so the vertical length and the
Schur scale are computed inside one boundary action.

\subsection{Neutral Matrix Extraction}

Set
\[
  \Theta_\Gamma
  =
  \binom{\theta_\Gamma^3}{\theta_\Gamma^Y},
  \qquad
  I_{AB}^\Gamma
  =
  \int_{\Gamma_{\rm EW,int}}d\nu_\Gamma\,
  \gamma_\Gamma^{-1}(\theta_\Gamma^A,\theta_\Gamma^B),
  \quad A,B\in\{3,Y\}.
\]
The localized action extracts the neutral length matrix as
\[
  {\cal L}_\Gamma
  =
  \begin{pmatrix}
  I_{33}^\Gamma&I_{3Y}^\Gamma\\
  I_{3Y}^\Gamma&I_{YY}^\Gamma
  \end{pmatrix},
\]
so
\[
  L_2^2=I_{33}^\Gamma,\qquad
  L_Y^2=I_{YY}^\Gamma,\qquad
  \epsilon_\Gamma=I_{3Y}^\Gamma.
\]
Equivalently,
\[
  u=\frac{I_{33}^\Gamma}{I_{YY}^\Gamma},
  \qquad
  \chi_\Gamma=
  \frac{I_{3Y}^\Gamma}{\sqrt{I_{33}^\Gamma I_{YY}^\Gamma}},
  \qquad
  D_\Gamma=u+1+2\chi_\Gamma\sqrt u .
\]
The common current scale is
\[
  \Lambda_\Gamma=\frac{C^2}{I_{YY}^\Gamma}.
\]
The local boundary calculation then gives
\[
  \alpha_\Gamma
  =
  \frac{\Lambda_\Gamma}{4\pi D_\Gamma},
  \qquad
  M_W^2
  =
  \frac{\Lambda_\Gamma v_{\rm EW}^2}{4u},
  \qquad
  M_Z^2
  =
  \frac{\Lambda_\Gamma v_{\rm EW}^2D_\Gamma}
       {4u(1-\chi_\Gamma^2)}.
\]
The de Vries pass condition becomes the integral constraint
\[
  R_{\rm dV}
  =
  \frac{I_{33}^\Gamma I_{YY}^\Gamma-(I_{3Y}^\Gamma)^2}
       {I_{33}^\Gamma
       (I_{33}^\Gamma+I_{YY}^\Gamma+2I_{3Y}^\Gamma)}.
\]
The D10 task has therefore been reduced to constructing
\(\theta_\Gamma^3,\theta_\Gamma^Y,d\nu_\Gamma\) from the localized
DBI/CS boundary system and evaluating these three current integrals.

A rank-one specialization gives a viability test.  If
\[
  \theta_\Gamma^Y=c\,\theta_\Gamma^3
\]
in the retained current Hilbert space, then
\[
  I_{YY}^\Gamma=c^2I_{33}^\Gamma,
  \qquad
  I_{3Y}^\Gamma=cI_{33}^\Gamma,
  \qquad
  \det{\cal L}_\Gamma=0.
\]
The neutral coefficient \(G_Z^\Gamma=L_Q^2/\det{\cal L}_\Gamma\) lacks a
finite value.  Thus the \(S^1_Q\) factor supplies photon protection, and a
viable neutral matrix requires a second current direction from the
\(\CP^2\), Chan--Paton, or boundary-complex sector.

A rank-two normal form makes this requirement explicit.  Let \(e_Q,e_N\) be
orthonormal retained current directions, with \(e_Q\) the \(Q\)-protected
photon direction, and write
\[
  \theta_\Gamma^3=a_\Gamma e_Q+b_\Gamma e_N,
  \qquad
  \theta_\Gamma^Y=c_\Gamma e_Q-b_\Gamma e_N.
\]
Then
\[
  \theta_\Gamma^Q=\theta_\Gamma^3+\theta_\Gamma^Y
  =
  (a_\Gamma+c_\Gamma)e_Q,
\]
so the massive-neutral direction cancels from the photon current.  The
integrals are
\[
  I_{33}^\Gamma=a_\Gamma^2+b_\Gamma^2,
  \qquad
  I_{YY}^\Gamma=c_\Gamma^2+b_\Gamma^2,
  \qquad
  I_{3Y}^\Gamma=a_\Gamma c_\Gamma-b_\Gamma^2.
\]
Consequently
\[
  L_Q^2=(a_\Gamma+c_\Gamma)^2,
  \qquad
  \det{\cal L}_\Gamma=b_\Gamma^2(a_\Gamma+c_\Gamma)^2,
\]
and
\[
  G_Z^\Gamma=\frac{1}{b_\Gamma^2},
  \qquad
  G_W^\Gamma=\frac{1}{a_\Gamma^2+b_\Gamma^2}.
\]
The de Vries pass condition is therefore the current-splitting condition
\[
  R_{\rm dV}
  =
  \frac{M_W^2}{M_Z^2}
  =
  \frac{b_\Gamma^2}{a_\Gamma^2+b_\Gamma^2},
\]
or
\[
  b_\Gamma^2=R_{\rm dV}L_2^2,
  \qquad
  a_\Gamma^2=(1-R_{\rm dV})L_2^2,
  \qquad
  \frac{|b_\Gamma|}{|a_\Gamma|}
  =
  \sqrt{\frac{R_{\rm dV}}{1-R_{\rm dV}}}
  =
  1.866083698\ldots .
\]
The neutral datum is \(L_Q=|a_\Gamma+c_\Gamma|\), equivalently
\[
  c_\Gamma=\sigma_Q L_Q-a_\Gamma,\qquad \sigma_Q=\pm1.
\]

The second direction has a canonical boundary-line source.  The combined
weak-hypercharge frame bundle carries the neutral character
\[
  \chi_N(u_T,u_Y)=u_T^{-1}u_Y,
  \qquad
  L_N=L_{T_3}^{-1}\otimes L_Y,
\]
whose connection is represented by
\[
  \theta_\Gamma^N=\theta_\Gamma^Y-\theta_\Gamma^3.
\]
For the neutral finite section,
\[
  T_3\eta_0=-\frac12\eta_0,
  \qquad
  Y\eta_0=\frac12\eta_0,
\]
and therefore
\[
  (T_3\theta_\Gamma^3+Y\theta_\Gamma^Y)\eta_0
  =
  \frac12\theta_\Gamma^N\eta_0.
\]
The massive-neutral current is the current-metric projection of this neutral
line away from the photon direction:
\[
  P_{Q^\perp}\omega
  =
  \omega
  -
  \frac{\langle \omega,\theta_\Gamma^Q\rangle_\Gamma}
       {\langle \theta_\Gamma^Q,\theta_\Gamma^Q\rangle_\Gamma}
  \theta_\Gamma^Q ,
  \qquad
  e_N
  =
  -\frac{P_{Q^\perp}\theta_\Gamma^N}
        {\|P_{Q^\perp}\theta_\Gamma^N\|_\Gamma}.
\]
The projection has a matrix expression:
\[
  \|P_{Q^\perp}\theta_\Gamma^N\|_\Gamma^2
  =
  \|\theta_\Gamma^N\|_\Gamma^2
  -
  \frac{\langle\theta_\Gamma^N,\theta_\Gamma^Q\rangle_\Gamma^2}
       {\|\theta_\Gamma^Q\|_\Gamma^2}
  =
  \frac{4\det{\cal L}_\Gamma}{L_Q^2}.
\]
The rank-two normal form gives
\[
  P_{Q^\perp}\theta_\Gamma^N=-2b_\Gamma e_N,
  \qquad
  \|P_{Q^\perp}\theta_\Gamma^N\|_\Gamma^2=4b_\Gamma^2.
\]
Thus the de Vries pass becomes the projected neutral-line norm condition
\[
  \|P_{Q^\perp}\theta_\Gamma^N\|_\Gamma^2
  =
  4R_{\rm dV}L_2^2.
\]
Equivalently,
\[
  \det{\cal L}_\Gamma=R_{\rm dV}L_2^2L_Q^2.
\]
This replaces the abstract second-current condition by concrete D10 data:
\(L_N\), its projected boundary norm using \(d\nu_\Gamma\), and the
comparison with the weak-current norm.

The associated line-lattice data are
\[
  c_1(L_Y)=6h+7\eta,
  \qquad
  c_1(L_{T_3})=\eta,
\]
so
\[
  c_1(L_N)
  =
  c_1(L_Y)-c_1(L_{T_3})
  =
  6(h+\eta).
\]
The topological class fixes the neutral line; the physical condition is the
projected norm
\[
  {\cal N}_N^\Gamma
  =
  \int_{\Gamma_{\rm EW,int}}d\nu_\Gamma\,
  \gamma_\Gamma^{-1}
  \left(
  P_{Q^\perp}\theta_\Gamma^N,
  P_{Q^\perp}\theta_\Gamma^N
  \right),
\]
with
\[
  \frac{{\cal N}_N^\Gamma}{4L_2^2}
  =
  R_{\rm dV}
  =
  0.776898678\ldots .
\]
Thus the diagonal six-unit line class supplies the neutral line class, and the
localized DBI/CS action supplies the current metric whose projected norm gives
the de Vries ratio.

Equivalently, let \(\omega_h,\omega_\eta\) be the localized boundary
representatives of \(h,\eta\) and define
\[
  H_{ij}^\Gamma
  =
  \int_{\Gamma_{\rm EW,int}}d\nu_\Gamma\,
  \gamma_\Gamma^{-1}(\omega_i,\omega_j),
  \qquad i,j\in\{h,\eta\}.
\]
For
\[
  \omega_T=\omega_\eta,
  \qquad
  \omega_Y=6\omega_h+7\omega_\eta,
\]
the neutral and photon representatives are
\[
  \omega_N=6(\omega_h+\omega_\eta),
  \qquad
  \omega_Q=6\omega_h+8\omega_\eta.
\]
Hence
\[
  {\cal N}_N^\Gamma
  =
  \langle \omega_N,\omega_N\rangle_\Gamma
  -
  \frac{\langle \omega_N,\omega_Q\rangle_\Gamma^2}
       {\langle \omega_Q,\omega_Q\rangle_\Gamma},
\]
with
\[
  \langle \omega_N,\omega_N\rangle_\Gamma
  =
  36\left(H_{hh}^\Gamma+2H_{h\eta}^\Gamma+H_{\eta\eta}^\Gamma\right),
\]
\[
  \langle \omega_N,\omega_Q\rangle_\Gamma
  =
  36H_{hh}^\Gamma+84H_{h\eta}^\Gamma+48H_{\eta\eta}^\Gamma,
\]
and
\[
  \langle \omega_Q,\omega_Q\rangle_\Gamma
  =
  36H_{hh}^\Gamma+96H_{h\eta}^\Gamma+64H_{\eta\eta}^\Gamma.
\]
The neutral-line pass becomes a boundary Hodge-matrix equation:
\[
  {\cal N}_N^\Gamma=4R_{\rm dV}L_2^2.
\]
Writing
\[
  A=H_{hh}^\Gamma,\qquad B=H_{h\eta}^\Gamma,\qquad C=H_{\eta\eta}^\Gamma,
\]
the projection simplifies to
\[
  {\cal N}_N^\Gamma
  =
  \frac{36(AC-B^2)}{9A+24B+16C}.
\]
Thus the localized pass equation is
\[
  \frac{9(AC-B^2)}{9A+24B+16C}
  =
  R_{\rm dV}L_2^2.
\]
The weak-current representative \(\omega_T=\omega_\eta\) gives a useful
specialization:
\[
  L_2^2=C=H_{\eta\eta}^\Gamma.
\]
With
\[
  r_H=\frac{A}{C},\qquad s_H=\frac{B}{C},
\]
the pass condition is
\[
  \frac{9(r_H-s_H^2)}{9r_H+24s_H+16}
  =
  R_{\rm dV},
\]
so
\[
  r_H(s_H)
  =
  \frac{9s_H^2+24R_{\rm dV}s_H+16R_{\rm dV}}
       {9(1-R_{\rm dV})}.
\]
Along this curve,
\[
  r_H-s_H^2
  =
  \frac{R_{\rm dV}(3s_H+4)^2}{9(1-R_{\rm dV})},
  \qquad
  9r_H+24s_H+16
  =
  \frac{(3s_H+4)^2}{1-R_{\rm dV}}.
\]
The diagonal Hodge case gives the concrete target
\[
  s_H=0,\qquad
  \frac{H_{hh}^\Gamma}{H_{\eta\eta}^\Gamma}
  =
  \frac{16R_{\rm dV}}{9(1-R_{\rm dV})}
  =
  6.190\ldots .
\]
For comparison, the bulk KLT harmonic representatives
\[
  h=x+y,\qquad
  \eta_{\rm harm}=-\frac12x+\frac12y+z
\]
with \(x=e^{12}\), \(y=e^{34}\), \(z=e^{56}\), have weights
\[
  \|x\|^2=\|y\|^2\propto\lambda^{1/2},
  \qquad
  \|z\|^2\propto\lambda^{-3/2}.
\]
They give
\[
  A_{\rm bulk}=2\lambda^{1/2},\qquad
  B_{\rm bulk}=0,\qquad
  C_{\rm bulk}=\frac12\lambda^{1/2}+\lambda^{-3/2},
\]
so
\[
  r_H^{\rm bulk}(\lambda)
  =
  \frac{4\lambda^2}{\lambda^2+2}.
\]
On \(2/5\leq\lambda\leq2\), this ranges from \(0.296\ldots\) to
\(2.666\ldots\).  The boundary-current pass therefore requires either a
localized boundary Hodge measure or a weak-current frame factor beyond the
unlocalized KLT harmonic metric.  If
\[
  L_2^2=\kappa_H C_{\rm bulk},
\]
then
\[
  \kappa_H(\lambda)
  =
  \frac{9\lambda^2}{R_{\rm dV}(13\lambda^2+8)}.
\]
At \(\lambda=0.514554893751227\), one obtains \(\kappa_H=0.268\ldots\).

A localized representative can also shift the horizontal component:
\[
  \eta(t)=t(x+y)+y+z,\qquad q=\lambda^2.
\]
For the same KLT weights,
\[
  A=2q,\qquad
  B=(2t+1)q,\qquad
  C=(2t^2+2t+1)q+1,
\]
so
\[
  AC-B^2=q(q+2),
\]
and the weak-current identification gives
\[
  9q(q+2)
  =
  2R_{\rm dV}
  (16qt^2+40qt+29q+8)(q(2t^2+2t+1)+1).
\]
For example, \(\lambda=1.915\) has roots
\[
  t=-0.811\ldots,\qquad (r_H,s_H)=(2.070\ldots,-0.644\ldots),
\]
and
\[
  t=-0.939\ldots,\qquad (r_H,s_H)=(1.727\ldots,-0.758\ldots).
\]
The quartic can be centered by setting \(s=t+7/8\).  It becomes
\[
  F(q,s)=F_0(q)-2R_{\rm dV}q(7q+32)s^2-64R_{\rm dV}q^2s^4,
\]
where
\[
  F_0(q)
  =
  \frac{576q^2+1152q-R_{\rm dV}(625q^2+1600q+1024)}{64}.
\]
Consequently real representatives occur for \(F_0(q)\ge0\).  The positive
threshold is
\[
  q_c=
  \frac{32(25R_{\rm dV}+18\sqrt{1-R_{\rm dV}}-18)}
       {576-625R_{\rm dV}}
  =
  3.511614004347\ldots ,
  \qquad
  \lambda_c=1.873930095907\ldots .
\]
For \(q\ge q_c\), the roots are
\[
  t_\pm=-\frac78\pm
  \sqrt{
  \frac{-(7q+32)+\sqrt{(7q+32)^2+64F_0(q)/R_{\rm dV}}}
       {64q}} .
\]
At \(\lambda=2\), this gives
\[
  t_-=-0.983371\ldots,\qquad t_+=-0.766629\ldots .
\]
The localized representative \(t\) is therefore a computable boundary datum
near \(t=-7/8\).
With a weak-current frame factor \(L_2^2=\kappa_H C\), the same calculation
gives
\[
  \kappa_H(q,t)
  =
  \frac{9q(q+2)}
       {2R_{\rm dV}(16qt^2+40qt+29q+8)(q(2t^2+2t+1)+1)}.
\]
The maximum over this representative family occurs at \(t=-7/8\):
\[
  \kappa_H^{\max}(q)
  =
  \frac{576q(q+2)}{R_{\rm dV}(25q+32)^2}.
\]
Thus \(\kappa_H^{\max}=1\) gives the same threshold \(q_c\).  Below
\(\lambda_c\), the neutral boundary condition requires a reduced
weak-current frame factor.
The center also has a boundary-orthogonality interpretation.  For horizontal
weights
\[
  \langle x,x\rangle_\partial=w_x,\qquad
  \langle y,y\rangle_\partial=w_y,\qquad
  \langle x,y\rangle_\partial=0,
\]
the condition \(\langle\eta(t),h\rangle_\partial=0\) gives
\[
  t=-\frac{w_y}{w_x+w_y}.
\]
Hence \(t=-7/8\) corresponds to a localized horizontal split
\[
  w_y:w_x=7:1.
\]
For the weak-period value \(Q=7\), this split is the special case of
\[
  w_y:w_x=Q:1,\qquad t_Q=-\frac{Q}{Q+1}.
\]
Together with the primitive line-lattice relation \(P=Q-1\), the retained
class
\[
  \zeta=P h+Q\eta(t_Q)
\]
has the form
\[
  \zeta
  =
  Q(y+z)-\frac{h}{Q+1}.
\]
Thus for \(Q=7\),
\[
  \zeta=7(y+z)-\frac18h,\qquad 8\zeta=56(y+z)-h.
\]
The centered representative therefore converts the weak period and the single
weak \(\cO(1)\) unit into an eightfold boundary lift.
The D8/O8 bookkeeping uses
\[
  N_8=2n,\qquad
  q_8^{\rm net}=N_8-16=2(n-8),\qquad
  k=n-8.
\]
Identifying the lift denominator with the Chan--Paton half-rank,
\[
  n=Q+1,
\]
gives, for \(Q=7\),
\[
  n=8,\qquad N_8=16,\qquad q_8^{\rm net}=0,\qquad k=0.
\]
The same integer that centers the representative therefore selects the locally
neutral D8/O8 endpoint.
Borel--Weil supplies a concrete source for the half-rank relation.  The
weak-period line has
\[
  L_Q|_{\CP^1}\simeq\cO(Q),
\]
and hence
\[
  H^0(\CP^1,\cO(Q))\simeq{\rm Sym}^Q(\mathbb C^2)^*,
  \qquad
  \dim H^0(\CP^1,\cO(Q))=Q+1.
\]
Taking the Chan--Paton half-stack to be this zero-mode space gives
\[
  n=\dim H^0(\CP^1,\cO(Q))=Q+1.
\]
For \(Q=7\),
\[
  N_8=2\,\dim H^0(\CP^1,\cO(7))=16.
\]
The same line gives the horizontal weight split through point evaluation.
Fix \(p\in\CP^1\).  Since the ideal sheaf of \(p\) is \(\cO(-1)\), tensoring
by \(\cO(Q)\) gives
\[
  0\longrightarrow \cO(Q-1)
  \longrightarrow \cO(Q)
  \stackrel{\mathrm{ev}_p}{\longrightarrow}
  \cO(Q)|_p
  \longrightarrow 0 .
\]
For \(Q\ge1\), global sections form the exact sequence
\[
  0\longrightarrow H^0(\CP^1,\cO(Q-1))
  \longrightarrow H^0(\CP^1,\cO(Q))
  \stackrel{\mathrm{ev}_p}{\longrightarrow}
  \cO(Q)|_p
  \longrightarrow 0 ,
\]
with
\[
  \dim\ker{\rm ev}_p=Q,\qquad
  \dim{\rm im}\,{\rm ev}_p=1 .
\]
In coordinates \([u:v]\) with \(p=[1:0]\),
\[
  H^0(\CP^1,\cO(Q))
  =
  {\rm span}\{u^Q,u^{Q-1}v,\ldots,v^Q\},
\]
and \({\rm ev}_p\) selects the coefficient of \(u^Q\).  If \(P_p\) projects
onto \(u^Q\) and \(K_p=1-P_p\), the trace metric gives
\[
  {\rm Tr}\,P_p=1,\qquad {\rm Tr}\,K_p=Q.
\]
The same projector gives the local current metric on the transverse weak
plane.  Let \(J_Q=Q/2\), let \(|0\rangle=|J_Q,J_Q\rangle\), and define
\[
  G_{ab}^{(p)}
  =
  {\rm Re}\,\langle0|T_aK_pT_b|0\rangle .
\]
With
\[
  T_1=\frac12(T_++T_-),\qquad
  T_2=\frac{1}{2i}(T_+-T_-),\qquad
  T_3|0\rangle=J_Q|0\rangle,
\]
and \(T_-|0\rangle=\sqrt Q\,|J_Q,J_Q-1\rangle\), one obtains
\[
  G_{ab}^{(p)}
  =
  \frac Q4
  \begin{pmatrix}
  1&0&0\\
  0&1&0\\
  0&0&0
  \end{pmatrix}.
\]
With common localized density factor \(f_\Gamma^2\), the transverse Hosotani
normalization is
\[
  f_{H,\perp}^2=f_\Gamma^2\frac Q4.
\]
Thus a holonomy scale \(\mu_\Gamma\) gives
\[
  \alpha_{\rm mix}=\frac{2\mu_\Gamma}{\sqrt Q\,f_\Gamma},
  \qquad
  \beta_{\rm diag}=\frac{4\mu_\Gamma^2}{Qf_\Gamma^2},
  \qquad
  \frac{\beta_{\rm diag}}{\alpha_{\rm mix}^2}=1 .
\]
The endpoint line also carries the neutral weak-hypercharge character.  Write
\(Q_{\rm w}\) for the weak-period integer here and set
\[
  J_{\rm w}=\frac{Q_{\rm w}}2,\qquad |0\rangle=|J_{\rm w},J_{\rm w}\rangle,
  \qquad
  Y|0\rangle=-J_{\rm w}|0\rangle .
\]
Then \(Q_{\rm em}=T_3+Y\) satisfies \(Q_{\rm em}|0\rangle=0\).  For
\((T_3,Y)\), the endpoint charge vector is \(q=(J_{\rm w},-J_{\rm w})\), hence
\[
  N_{AB}^{(p)}
  =
  \langle0|q_Aq_B|0\rangle
  =
  \frac{Q_{\rm w}^2}{4}
  \begin{pmatrix}
  1&-1\\
  -1&1
  \end{pmatrix}.
\]
The null vector is \((1,1)\), the \(Q_{\rm em}\) direction.  The vector
\((1,-1)\) has eigenvalue \(Q_{\rm w}^2/2\).  This endpoint-line block is the
rank-one neutral mass block; the gauge-kinetic neutral lengths are supplied by
the localized current matrix.
With gauge fields \((W^3_\mu,B_\mu)\), this block becomes
\[
  M_N^2
  =
  v_\Gamma^2J_{\rm w}^2
  \begin{pmatrix}
  g_2^2&-g_2g_Y\\
  -g_2g_Y&g_Y^2
  \end{pmatrix}.
\]
Thus
\[
  A_\mu=\frac{g_YW^3_\mu+g_2B_\mu}{\sqrt{g_2^2+g_Y^2}},
  \qquad
  Z_\mu=\frac{g_2W^3_\mu-g_YB_\mu}{\sqrt{g_2^2+g_Y^2}},
\]
and
\[
  M_\gamma^2=0,\qquad
  M_Z^2=v_\Gamma^2J_{\rm w}^2(g_2^2+g_Y^2).
\]
The transverse \(K_p\) metric gives
\[
  M_W^2=v_\Gamma^2\frac{Q_{\rm w}}4g_2^2.
\]
For \(Q_{\rm w}=1\), the \(\mathcal O(1)\) doublet gives
\[
  M_W^2=\frac{g_2^2v_\Gamma^2}{4},
  \qquad
  M_Z^2=\frac{(g_2^2+g_Y^2)v_\Gamma^2}{4}.
\]
The Connes finite-line amplitude satisfies
\[
  |v_F|=\frac{v_{\rm EW}}{\sqrt2},
\]
so the endpoint gauge-mass convention is
\[
  v_\Gamma=v_{\rm EW}=\sqrt2\,|v_F|.
\]
Equivalently,
\[
  M_W^2=\frac{g_2^2|v_F|^2}{2},
  \qquad
  M_Z^2=\frac{(g_2^2+g_Y^2)|v_F|^2}{2}.
\]
The \(Q_{\rm w}=7\) period stack supplies the boundary filtration, half-rank,
and endpoint metric factors; electroweak mass normalization uses the
\(\mathcal O(1)\) doublet together with the localized current lengths.
The endpoint convention then gives the localized Schur scale target
\[
  \mu_{\Gamma,{\rm Schur}}^2
  =
  \mu_*^2
  =
  \frac{M_Z^2}{\Xp(2)}
  =
  \frac{g_2^2+g_Y^2}{4\Xp(2)}v_\Gamma^2
  =
  \frac{g_2^2+g_Y^2}{2\Xp(2)}|v_F|^2 .
\]
For a stationary Hosotani scalar with
\[
  f_H\theta_{H,*}=|v_F|=\frac{v_{\rm EW}}{\sqrt2},
  \qquad
  v_\Gamma=\sqrt2\,f_H\theta_{H,*},
\]
this becomes
\[
  \mu_{\Gamma,{\rm Schur}}^2
  =
  \frac{g_2^2+g_Y^2}{2\Xp(2)}f_H^2\theta_{H,*}^2 .
\]
Thus the D10 matching datum is the dimensionless scale-ratio coordinate
\[
  \chi_H
  =
  \frac{2\Xp(2)\mu_{\Gamma,{\rm Schur}}^2}
       {(g_2^2+g_Y^2)f_H^2\theta_{H,*}^2}
  =
  1.
\]
With the point-projector normalization \(f_H^2=Qf_\Gamma^2/4\), the same
target is
\[
  \mu_{\Gamma,{\rm Schur}}^2
  =
  \frac{Q(g_2^2+g_Y^2)}{8\Xp(2)}
  f_\Gamma^2\theta_{H,*}^2 .
\]
Writing \(x=w_H\theta_H\), the minimal localized Wilson potential with an
interior stationary value is
\[
  V_H(\theta_H)=c_1\cos x+c_2\cos(2x).
\]
Its stationary equation is
\[
  \partial_{\theta_H}V_H
  =
  -w_H\sin x\,[c_1+4c_2\cos x],
\]
so the interior branch satisfies
\[
  \cos x_*=-\frac{c_1}{4c_2},
  \qquad
  c_2>0,\qquad |c_1|<4c_2.
\]
Matching the finite-sheet value
\[
  x_*=\frac{w_Hv_{\rm EW}}{\sqrt2\,f_H}
\]
fixes the Fourier ratio
\[
  \frac{c_1}{c_2}=-4\cos\!\left(
  \frac{w_Hv_{\rm EW}}{\sqrt2\,f_H}
  \right).
\]
The canonical scalar curvature from this Wilson sector is
\[
  m_\theta^2
  =
  \frac{1}{f_H^2}\partial_{\theta_H}^2V_H\big|_{\theta_{H,*}}
  =
  \frac{4w_H^2c_2}{f_H^2}\sin^2x_*.
\]
The full-curvature specialization \(m_\theta^2=M_H^2\) gives
\[
  c_2=
  \frac{M_H^2f_H^2}{4w_H^2\sin^2x_*},
  \qquad
  c_1=
  -\frac{M_H^2f_H^2}{w_H^2}
  \frac{\cos x_*}{\sin^2x_*}.
\]
With \(M_H^2=2\lambda_Fv_{\rm EW}^2\), the same coefficients are
\[
  c_2=
  \frac{\lambda_Fv_{\rm EW}^2f_H^2}{2w_H^2\sin^2x_*},
  \qquad
  c_1=
  -\frac{2\lambda_Fv_{\rm EW}^2f_H^2}{w_H^2}
  \frac{\cos x_*}{\sin^2x_*}.
\]
Eliminating \(f_H\) by
\[
  f_H=\frac{w_Hv_{\rm EW}}{\sqrt2\,x_*}
\]
gives the dimensionless coefficient target
\[
  \frac{c_2}{v_{\rm EW}^4}
  =
  \frac{\lambda_F}{4x_*^2\sin^2x_*},
  \qquad
  \frac{c_1}{v_{\rm EW}^4}
  =
  -\frac{\lambda_F\cos x_*}{x_*^2\sin^2x_*}.
\]
A localized spectrum calculation of \(c_1/c_2\) fixes
\[
  x_*=\arccos\!\left(-\frac{c_1}{4c_2}\right),
  \qquad
  f_H=\frac{w_Hv_{\rm EW}}{\sqrt2\,x_*},
\]
on the chosen \(0<x_*<\pi\) minimum branch.
For the standard four-dimensional Coleman--Weinberg integral over one
circle, a massless charge-one tower contributes
\[
  c_m^{(1)}=\frac{A_1}{m^5}.
\]
Such a tower gives \(c_1/c_2=32\).  The interior minimum condition
\[
  \left|\frac{c_1}{c_2}\right|<4
\]
therefore demands additional charged sectors, massive thresholds, or
localized boundary terms.  With charge-one amplitude \(A_1\) and charge-two
amplitude \(A_2\) at the fundamental period,
\[
  c_1=A_1,\qquad c_2=\frac{A_1}{32}+A_2,
\]
and the stationary electroweak angle fixes
\[
  \frac{A_2}{A_1}
  =
  -\frac{1}{4\cos x_*}-\frac{1}{32}.
\]
The local minimum condition \(c_2>0\) gives two sign branches:
\[
  A_1>0:\quad
  \frac{A_2}{A_1}>\frac{7}{32},
  \qquad \frac{\pi}{2}<x_*<\pi,
\]
and
\[
  A_1<0:\quad
  \frac{A_2}{A_1}<-\frac{9}{32},
  \qquad 0<x_*<\frac{\pi}{2}.
\]
More generally, let \(A_r\) be the total signed amplitude of massless towers
with charge \(r\) under \(x=w_H\theta_H\).  Then
\[
  V_H(x)
  =
  \sum_{r\ge1}A_r
  \sum_{m\ge1}\frac{\cos(mrx)}{m^5}
  =
  \sum_{n\ge1}c_n\cos(nx),
\]
with divisor-sum coefficients
\[
  c_n=\frac{1}{n^5}\sum_{r|n}r^5A_r .
\]
Thus
\[
  c_1=A_1,\qquad
  c_2=\frac{A_1}{32}+A_2,\qquad
  c_3=\frac{A_1}{243}+A_3,\qquad
  c_4=\frac{A_1}{1024}+\frac{A_2}{32}+A_4 .
\]
The two-harmonic approximation requires the higher harmonics to be small
against \(|c_1|+|c_2|\).  Massive thresholds and localized boundary terms
enter through the same harmonic bookkeeping with threshold-adjusted
coefficients.
In the Borel--Weil convention \(H^0(\CP^1,\cO(n))\simeq V_{n/2}\), use
integer current charge \(q=2m\).  The charge set is
\[
  q=-n,-n+2,\ldots,n-2,n.
\]
Thus \(\cO(1)\) supplies \(|q|=1\), \(\cO(2)\) supplies \(|q|=2\), and
\(\cO(7)\) supplies \(|q|=1,3,5,7\).  The physical doublet gives the
charge-one tower.  The active \(\cO(1)\oplus\cO(2)\) charged set has
\(\gcd(1,2)=1\), and the \(\cO(7)\) stack has \(\gcd(1,3,5,7)=1\).  Thus
\[
  w_H=1
\]
in the integer current-charge convention.  A global-form lift rescales the
endpoint charge lattice by an explicit factor.  The weak-period stack gives
zero \(A_2\) and higher odd harmonics.  The minimal Borel--Weil
second-harmonic source is
\[
  \cO(2)\simeq {\rm Sym}^2\cO(1),
\]
with
\[
  A_1>0:\quad A_2>\frac{7}{32}A_1,
  \qquad
  A_1<0:\quad A_2<-\frac{9}{32}A_1.
\]
The line-square sector used for the \(j=1\) Schur anchor therefore supplies
the Hosotani second harmonic in the minimal localized spectrum.
Let the localized towers coupled to \(\gamma_H\) have integral charges
\(q_\alpha\), one-loop signs \(\sigma_\alpha\), degeneracies \(d_\alpha\),
and positive threshold weights \(\tau_\alpha\).  With common loop factor
\({\cal C}_H>0\),
\[
  A_r=
  {\cal C}_H
  \sum_{\alpha:\ |q_\alpha|=r}
  \sigma_\alpha d_\alpha\tau_\alpha .
\]
For \(\cO(1)\oplus\cO(2)\), write the signed charged weights as
\[
  S_1=
  \sum_{\alpha\in\cO(1),\ |q_\alpha|=1}
  \sigma_\alpha d_\alpha\tau_\alpha,
  \qquad
  S_2=
  \sum_{\alpha\in\cO(2),\ |q_\alpha|=2}
  \sigma_\alpha d_\alpha\tau_\alpha .
\]
Then \(A_2/A_1=S_2/S_1\), and the finite Hosotani minimum requires
\[
  S_1>0:\quad \frac{S_2}{S_1}>\frac{7}{32},
  \qquad
  S_1<0:\quad \frac{S_2}{S_1}<-\frac{9}{32}.
\]
A same-sign charged-weight pass has \(S_2=S_1\).  It gives
\[
  \frac{c_1}{c_2}=\frac{32}{33},
  \qquad
  \cos x_*=-\frac{8}{33},
  \qquad
  x_*=1.815660181000518\ldots ,
\]
and hence
\[
  \frac{f_H}{v_{\rm EW}}
  =
  \frac{w_H}{\sqrt2\,x_*}
  =
  0.389448856446748\ldots\, w_H .
\]
For \(m_\theta^2=M_H^2=2\lambda_Fv_{\rm EW}^2\),
\[
  \frac{c_2}{v_{\rm EW}^4}
  =
  0.080570282164285\ldots\,\lambda_F,
  \qquad
  \frac{c_1}{v_{\rm EW}^4}
  =
  0.078128758462337\ldots\,\lambda_F .
\]
The Borel--Weil weights occur in charge-conjugate pairs.  If the localized
real structure pairs \(+r\) with \(-r\) and preserves the threshold weight,
then \(S_r=2\sigma_rd_r\tau_r\).  The \(\cO(1)\) sector has \(q=\pm1\); the
\(\cO(2)\) sector has \(q=\pm2\), with a \(\theta_H\)-independent middle
state at \(q=0\).  Hence
\[
  \frac{S_2}{S_1}
  =
  \frac{\sigma_2d_2\tau_2}{\sigma_1d_1\tau_1}
  =:\chi_{12}.
\]
The equal-weight pass is \(\chi_{12}=1\).  More generally,
\[
  \sigma_1d_1\tau_1>0:\quad \chi_{12}>\frac{7}{32},
  \qquad
  \sigma_1d_1\tau_1<0:\quad \chi_{12}<-\frac{9}{32}.
\]
For harmonic computations, split the threshold weight by wrapping number.
For a flat local circle pass with charged mass \(M_r\), set
\[
  z_r=M_rL_H,
  \qquad
  F(z)=e^{-z}\left(1+z+\frac{z^2}{3}\right).
\]
The \(m\)-th wrapping of the charge-\(r\) pair carries
\(\tau_{r,m}=F(mz_r)\).  With signed base pair weights
\[
  D_r={\cal C}_H\,2\sigma_rd_r,
  \qquad
  \chi_{12}^{(0)}=\frac{D_2}{D_1},
\]
the first two Fourier coefficients are
\[
  c_1=D_1F(z_1),
  \qquad
  c_2=\frac{D_1}{32}F(2z_1)+D_2F(z_2).
\]
Hence
\[
  \frac{c_2}{c_1}
  =
  \frac{F(2z_1)}{32F(z_1)}
  +
  \chi_{12}^{(0)}\frac{F(z_2)}{F(z_1)}.
\]
For a chosen stationary angle,
\[
  \chi_{12}^{(0)}
  =
  \frac{F(z_1)}{F(z_2)}
  \left[
  -\frac{1}{4\cos x_*}
  -
  \frac{F(2z_1)}{32F(z_1)}
  \right].
\]
The equal signed-base pass gives
\[
  \cos x_*=
  -\frac{1}{4\left(
  F(2z_1)/(32F(z_1))+F(z_2)/F(z_1)
  \right)}.
\]
At \(z_1=z_2=0\), the result is \(\cos x_*=-8/33\).
Since
\[
  F'(z)=-\frac13 e^{-z}z(1+z)\le0,
\]
the equal signed-base pass with \(D_1>0\) has a simple threshold domain:
\[
  F(z_2)>\frac14F(z_1)-\frac1{32}F(2z_1).
\]
For \(z_1=0\), this gives
\[
  F(z_2)>\frac{7}{32},
  \qquad
  z_2<3.773046018411012\ldots .
\]
For common thresholds \(z_2=z_1\), the pass condition is automatic because
\[
  1+\frac{F(2z_1)}{32F(z_1)}>\frac14 .
\]
The Borel--Weil sectors themselves give a distinguished zero-mode pass.  For
\(n\ge0\),
\[
  H^0(\CP^1,\cO(n))=\ker\bar\partial_{\cO(n)}.
\]
Thus the charge pair \(q=\pm1\) in \(\cO(1)\) and the charge pair \(q=\pm2\)
in \(\cO(2)\) lie in holomorphic zero-eigenvalue sectors.  The localized
determinant then has
\[
  z_1=z_2=0,
  \qquad F(z_1)=F(z_2)=F(2z_1)=1.
\]
The \(q=0\) middle state of \(\cO(2)\) is independent of \(\theta_H\).  With
common signed oscillator data,
\[
  \sigma_2d_2=\sigma_1d_1,
\]
one obtains
\[
  c_1=D_1,\qquad c_2=\frac{33}{32}D_1,
  \qquad
  \frac{c_1}{c_2}=\frac{32}{33}.
\]
Therefore
\[
  \cos x_*=-\frac{8}{33},
  \qquad
  \frac{f_H}{v_{\rm EW}}=
  0.389448856446748\ldots\,w_H .
\]
In homogeneous coordinates \([u:v]\) on \(\CP^1\), the charged bases are
\[
  H^0(\CP^1,\cO(1)):
  \quad u\ (q=+1),\quad v\ (q=-1),
\]
and
\[
  H^0(\CP^1,\cO(2)):
  \quad u^2\ (q=+2),\quad uv\ (q=0),\quad v^2\ (q=-2).
\]
Thus each charged pair has one positive-charge and one negative-charge complex
mode, and the \(uv\) mode is neutral for the Wilson angle.  For a common
localized boundary complex \({\cal B}_\Gamma\) whose oscillator operator acts
with the same sign and degeneracy on \((u,v)\) and \((u^2,v^2)\),
\[
  \sigma_2=\sigma_1,\qquad d_2=d_1,\qquad D_2=D_1.
\]
The holomorphic zero-mode pass then yields the coefficient ratio \(32/33\).
Equivalently, define charged projectors by
\[
  P_1^\chi:\ H^0(\CP^1,\cO(1))\to{\rm span}\{u,v\},
  \qquad
  P_2^\chi:\ H^0(\CP^1,\cO(2))\to{\rm span}\{u^2,v^2\}.
\]
Then
\[
  {\rm Tr}\,P_1^\chi={\rm Tr}\,P_2^\chi=2,
  \qquad
  P_2^\chi=1-|uv\rangle\langle uv|.
\]
For a common localized charged oscillator block
\[
  {\cal B}_\Gamma^\chi=B_0\otimes P_r^\chi,
\]
the signed base weights are
\[
  D_r={\cal C}_H\,\sigma\,{\rm Tr}_{P_r^\chi}{\bf 1}
  =
  2{\cal C}_H\sigma .
\]
Thus
\[
  D_2=D_1,
  \qquad
  V_H(x)=D_1\cos x+\frac{33}{32}D_1\cos(2x).
\]
The same equality can be expressed as an equivariant index calculation.  For
the \(S^1\) action generated by \(x=w_H\theta_H\),
\[
  {\rm Ind}_{S^1}\bar\partial_{\cO(n)}(x)
  =
  \sum_{k=0}^{n}\exp\!\big(i(n-2k)x\big),
  \qquad n\ge0,
\]
because \(H^1(\CP^1,\cO(n))=0\).  Consequently
\[
  I_1(x)=2\cos x,
  \qquad
  I_2(x)=1+2\cos(2x).
\]
The charged endpoint trace is
\[
  I_\chi(x)=I_1(x)+I_2(x)-1
  =
  2\cos x+2\cos(2x).
\]
The \(1\) in \(I_2\) is the Wilson-neutral \(uv\) mode.  The charged
character coefficients for the \(\cO(1)\) and \(\cO(2)\) sectors are equal.
Payen's boundary path integral fixes the trace convention
\cite{Payen2007}: the discretized endpoint integral uses
\[
  \int N_R\,Dg\,R(g)v_\kappa v_\kappa^\dagger R^\dagger(g)
  =
  {\bf 1}_{V_R},
\]
and the external Yang--Mills coupling reduces to
\[
  {\rm Tr}_{V_R}\,
  {\cal T}_{\partial\Sigma}
  \exp\left(
    i\int_{\partial\Sigma}d\tau\,R(A_\tau(\tau))
  \right).
\]
The endpoint determinant therefore uses the ordinary finite representation
trace with no dimension divisor.  In the charged Borel--Weil subspaces,
\[
  {\rm Tr}_{P_1^\chi}{\bf 1}
  =
  {\rm Tr}_{P_2^\chi}{\bf 1}
  =
  2,
\]
so the source-backed endpoint trace gives \(d_1=d_2=d_\chi\) before oscillator
thresholds.
An endpoint determinant with common sign and localized density on these
Dolbeault zero-mode sectors gives \(A_1=A_2\), hence
\[
  c_1=A_1,\qquad
  c_2=\frac{A_1}{32}+A_2=\frac{33}{32}A_1,
  \qquad
  \frac{c_1}{c_2}=\frac{32}{33}.
\]
In pair-normalized form this contribution is
\[
  V_{\rm BW}(x)
  =
  D\sum_{m\ge1}\frac{\cos(mx)+\cos(2mx)}{m^5},
\]
so that
\[
  c_1=D,\qquad
  c_2=\frac{33}{32}D .
\]
At the two-harmonic level this gives \(\cos x_*=-8/33\).  The full massless
character has stationary equation
\[
  \operatorname{Im}{\rm Li}_4(e^{ix})
  +
  2\operatorname{Im}{\rm Li}_4(e^{2ix})
  =0,
\]
with interior solution
\[
  x_{\rm full}=1.837672606049940\ldots,
  \qquad
  \frac{f_H}{v_{\rm EW}}
  =
  \frac{w_H}{\sqrt2\,x_{\rm full}}
  =
  0.384783872197163\ldots\,w_H .
\]
At this point
\[
  K_{\rm full}
  =
  -{\rm Re}\,{\rm Li}_3(e^{ix_{\rm full}})
  -
  4{\rm Re}\,{\rm Li}_3(e^{2ix_{\rm full}})
  =
  3.557941701817797\ldots ,
\]
and matching \(m_\theta^2=M_H^2=2\lambda_Fv_{\rm EW}^2\) gives
\[
  \frac{D}{v_{\rm EW}^4}
  =
  \frac{\lambda_F}{x_{\rm full}^2K_{\rm full}}
  =
  0.083227124394645\ldots\,\lambda_F .
\]
With a common threshold \(z\), the same character gives
\[
  V_z(x)
  =
  D\sum_{m\ge1}
  F(mz)\frac{\cos(mx)+\cos(2mx)}{m^5},
\]
and the stationary equation
\[
  \sum_{m\ge1}F(mz)
  \frac{\sin(mx)+2\sin(2mx)}{m^4}
  =0 .
\]
The root moves from \(x_0=1.837672606\ldots\) to
\[
  x_\infty=\arccos(-1/4)=1.8234765819\ldots
\]
as higher wrappings are suppressed.  Correspondingly
\[
  \frac{f_H}{w_Hv_{\rm EW}}
  =
  0.384783872\ldots\quad(z=0),
  \qquad
  0.387779469\ldots\quad(z=\infty).
\]
The corresponding D10 coefficient target is the determinant amplitude.  Write
\[
  {\cal N}_\Gamma^\chi
  =
  {\rm Str}_{\chi}\,\zeta_\alpha
\]
for the signed charged endpoint supertrace after canonical CP2 and boundary
zero-mode normalization.  The weights \(\zeta_\alpha\) are dimensionless and
may include threshold factors.  For a tower with quadratic density \(Z_\alpha\),
\[
  \frac12{\rm Tr}\log[Z_\alpha\Delta_\alpha(x)]
  =
  \frac12{\rm Tr}\log\Delta_\alpha(x)
  +
  \frac12{\rm Tr}\log Z_\alpha .
\]
The last term is angle-independent and is absorbed into the vacuum constant.
Thus the angle-dependent determinant uses the canonical signed supertrace.
The DBI/open-boundary density enters \(f_H^2\) and the threshold parameters.
Then
\[
  D=
  \frac{3}{4\pi^2}\,
  \frac{{\cal N}_\Gamma^\chi}{L_H^4}.
\]
The coefficient follows from the standard circle determinant
\[
  \frac12\sum_{n\in\mathbb Z}
  \int\frac{d^4p_E}{(2\pi)^4}
  \log\!\left[
    p_E^2+\left(\frac{2\pi n+x}{L_H}\right)^2
  \right]
\]
after Schwinger parameterization and Poisson resummation; the tower sign is
included in \({\cal N}_\Gamma^\chi\).
For the full massless character, the Higgs-curvature match is
\[
  \frac{3}{4\pi^2}
  \frac{{\cal N}_\Gamma^\chi}{L_H^4v_{\rm EW}^4}
  =
  0.083227124394645\ldots\,\lambda_F .
\]
Equivalently,
\[
  \frac{{\cal N}_\Gamma^\chi}{L_H^4v_{\rm EW}^4}
  =
  1.095225057620535\ldots\,\lambda_F .
\]
The same loop length enters the scalar kinetic term.  With arclength \(s\) on
\(\gamma_H\),
\[
  \omega_H=\frac{ds}{L_H},\qquad
  \int_{\gamma_H}\omega_H=1,\qquad
  \int_{\gamma_H}|\omega_H|^2ds=\frac1{L_H}.
\]
Writing \(Z_{\Gamma,H}\) for the localized density gives
\[
  f_H^2=\frac{Z_{\Gamma,H}}{L_H}.
\]
Since \(f_H\theta_{H,*}=v_{\rm EW}/\sqrt2\) and
\(x_*=w_H\theta_{H,*}\),
\[
  L_H=
  \frac{2x_*^2Z_{\Gamma,H}}{w_H^2v_{\rm EW}^2}.
\]
For the full massless-character root, the amplitude target becomes
\[
  \frac{{\cal N}_\Gamma^\chi w_H^8 v_{\rm EW}^4}
       {16x_*^8 Z_{\Gamma,H}^4}
  =
  1.095225057620535\ldots\,\lambda_F .
\]
The point-projector current metric in the transverse weak channel gives
\[
  f_H^2=\frac{Q}{4}f_\Gamma^2 .
\]
Combining this metric with \(x_{\rm full}=1.837672606049940\ldots\) gives
\[
  \frac{f_H}{v_{\rm EW}}
  =
  \frac{w_H}{\sqrt2\,x_{\rm full}}
  =
  0.384783872197163\ldots\,w_H,
\]
and hence
\[
  \frac{f_\Gamma}{v_{\rm EW}}
  =
  \frac{\sqrt2}{\sqrt Q\,x_{\rm full}}\,w_H .
\]
For the \(Q=7\) boundary stack,
\[
  \frac{f_\Gamma}{v_{\rm EW}}
  =
  0.290869266954901\ldots\,w_H .
\]
The loop and point-projector normalizations meet through
\[
  \frac{Z_{\Gamma,H}}{L_H}=\frac{Q}{4}f_\Gamma^2 .
\]
Eliminating \(L_H\) gives the localized density target
\[
  \frac{{\cal N}_\Gamma^\chi}{Z_{\Gamma,H}^4}
  =
  \frac{2279.126786313921\ldots}{w_H^8v_{\rm EW}^4}\,
  \lambda_F
\]
at the full massless-character root.
The kinetic density in this equation comes from the localized Yang--Mills
term
\[
  S_{\rm YM}^{\Gamma}
  =
  -\frac{T_8(2\pi\alpha')^2}{4}
  \int_{\Gamma_{\rm EW}}
  e^{-\phi}\sqrt{\gamma_\Gamma}\,
  {\rm tr}_{\rm end}(F_{MN}F^{MN}) .
\]
For the Wilson ansatz
\[
  A_s(x,y)=\theta_H(x)\omega_{H,s}(y)n_H^aT_a ,
  \qquad
  F_{\mu s}=(\partial_\mu\theta_H)\omega_{H,s}n_H^aT_a ,
\]
one obtains
\[
  f_H^2
  =
  T_8(2\pi\alpha')^2
  \int_{\Gamma_{\rm EW,int}}
  e^{-\phi}\sqrt{\gamma_\Gamma}\,
  g^{ss}\omega_{H,s}\omega_{H,s}\,
  {\cal T}_H ,
\]
with
\[
  {\cal T}_H=
  {\rm tr}_{\rm end}(n_H^aT_a\,n_H^bT_b)\Pi_{ab}^{\rm Schur}.
\]
For \(\omega_H=ds/L_H\) and uniform transverse density,
\[
  Z_{\Gamma,H}
  =
  T_8(2\pi\alpha')^2
  \int_{\CP^2}
  e^{-\phi}\sqrt{\gamma_\perp}\,
  {\cal T}_H .
\]
The determinant factor is the canonical charged-kernel supertrace
\[
  {\cal N}_\Gamma^\chi
  =
  \sum_{\alpha\in\ker{\cal B}_\Gamma^\chi}
  \sigma_\alpha d_\alpha\tau_\alpha .
\]
In the equal-amplitude Borel--Weil pass,
\[
  {\cal N}_{\Gamma,1}^\chi
  =
  {\cal N}_{\Gamma,2}^\chi
  =
  {\cal N}_\Gamma^\chi ,
\]
so the total charged supertrace over
\(\cO(1)\oplus\cO(2)\) is
\[
  {\cal N}_{\Gamma,{\rm tot}}^\chi
  =
  2{\cal N}_\Gamma^\chi .
\]
The curvature target uses the per-pair amplitude
\({\cal N}_\Gamma^\chi\), multiplying
\(\cos(mx)+\cos(2mx)\) with a common coefficient.
For the zero-threshold endpoint trial, the charged kernels are
\[
  H^0(\CP^1,\cO(1))^\chi={\rm span}\{u,v\},
  \qquad
  H^0(\CP^1,\cO(2))^\chi={\rm span}\{u^2,v^2\}.
\]
The \(uv\) state is neutral under the Wilson angle.  With common oscillator
multiplicity \(d_\chi\) and Dolbeault index sign \(\sigma_\chi\),
\[
  {\cal N}_\Gamma^\chi=2\sigma_\chi d_\chi .
\]
The Higgs-curvature match has \(D>0\) for \(\lambda_F>0\), hence
\[
  \sigma_\chi d_\chi>0 .
\]
For the minimal even endpoint determinant,
\[
  d_\chi=1,\qquad
  \sigma_\chi=+1,\qquad
  {\cal N}_\Gamma^\chi=2 .
\]
With threshold parameters \(z_r=M_rL_H\), the charged-pair amplitudes become
\[
  {\cal N}_{\Gamma,1}^\chi
  =
  2\sigma_\chi d_1F(z_1),
  \qquad
  {\cal N}_{\Gamma,2}^\chi
  =
  2\sigma_\chi d_2F(z_2),
\]
where
\[
  F(z)=e^{-z}\left(1+z+\frac{z^2}{3}\right).
\]
The equal-amplitude Borel--Weil shape requires
\[
  d_1F(z_1)=d_2F(z_2).
\]
A sufficient local source for this equality is a product Dolbeault boundary
complex
\[
  {\cal D}_\Gamma^\chi
  =
  \left(
  \bar\partial_{\cO(1)}^\chi
  \oplus
  \bar\partial_{\cO(2)}^\chi
  \right)\otimes{\bf 1}_{E_\chi}
  +
  {\bf 1}\otimes{\cal D}_{E,\chi}.
\]
On the product zero-mode reduction,
\[
  ({\cal D}_\Gamma^\chi)^\dagger{\cal D}_\Gamma^\chi
  =
  \Delta_{\bar\partial,\chi}\otimes{\bf 1}_{E_\chi}
  +
  {\bf 1}\otimes\Delta_{E,\chi}.
\]
% CONTROL: Bochner split test. A rough CP1 kinetic determinant gives
% alpha_B=kappa_B, alpha_C=0 and splits the O(1), O(2) thresholds.
% A pure Casimir threshold gives alpha_B=0, alpha_C=kappa_C and also
% splits them. The Hosotani pass row uses a cohomological determinant
% alpha_B=alpha_C=0 or a sourced mixed relation alpha_C=-4alpha_B/5.
For \(n\geq0\),
\[
  H^0(\CP^1,\cO(n))=\ker\bar\partial_{\cO(n)},
  \qquad
  H^1(\CP^1,\cO(n))=0 .
\]
Thus \(\Delta_{\bar\partial,\chi}\) vanishes on the charged projectors
\[
  P_1^\chi={\rm span}\{u,v\},
  \qquad
  P_2^\chi={\rm span}\{u^2,v^2\}.
\]
The charged-pair threshold spectrum is then supplied by the common oscillator
block \(\Delta_{E,\chi}\), so a single eigenvalue \(M^2\) gives
\[
  z_1=z_2=ML_H,\qquad d_1=d_2=d_\chi .
\]
Degree-dependent rough-Laplacian, Casimir, or curvature terms give the
splitting parameters in the general equation above.
The minimal common-threshold pass has
\[
  d_1=d_2=1,\qquad z_1=z_2=z,
\]
and gives
\[
  {\cal N}_\Gamma^\chi=2F(z)
\]
for positive endpoint parity.
The point-projector endpoint trace evaluates the Schur-channel trace factor:
\[
  G_{ab}^{(p)}
  =
  {\rm Re}\,\langle0|T_aK_pT_b|0\rangle
  =
  \frac{Q}{4}{\rm diag}(1,1,0)_{ab}.
\]
Writing
\[
  Z_{\Gamma,0}
  =
  T_8(2\pi\alpha')^2
  \int_{\CP^2}
  e^{-\phi}\sqrt{\gamma_\perp}\,{\cal Z}_{\rm end},
\]
one gets
\[
  Z_{\Gamma,H}=\frac{Q}{4}Z_{\Gamma,0},
  \qquad
  f_\Gamma^2=\frac{Z_{\Gamma,0}}{L_H}.
\]
The same localized measure defines the charged-kernel inner product.  For an
orthonormal determinant basis
\[
  \psi_\alpha\in
  L^2(\Gamma_{\rm EW,int},d\nu_\Gamma)\otimes E_\chi,
  \qquad
  \int_{\Gamma_{\rm EW,int}}d\nu_\Gamma\,|\psi_\alpha|^2=1 ,
\]
the angle-dependent Wilson determinant has the finite coefficient
\[
  {\cal N}_\Gamma^\chi
  =
  \sum_{\alpha\in\ker{\cal B}_\Gamma^\chi}
  \sigma_\alpha d_\alpha\tau_\alpha .
\]
Consequently the local density fixes the Hilbert normalization and the kinetic
scale \(Z_{\Gamma,0}\).  The Wilson amplitude is the finite supertrace over
the normalized kernel.
For \(Q=7\), the same determinant target becomes
\[
  \frac{{\cal N}_\Gamma^\chi}{Z_{\Gamma,0}^4}
  =
  \frac{21375.716460702053\ldots}
       {w_H^8v_{\rm EW}^4}\,\lambda_F .
\]
In the minimal even endpoint trial the equivalent form is
\[
  \frac{Z_{\Gamma,0}}{v_{\rm EW}}
  =
  \left(
  \frac{2}{21375.716460702053\ldots\,\lambda_F}
  \right)^{1/4}w_H^2 .
\]
With a common threshold \(z\), this becomes
\[
  \frac{Z_{\Gamma,0}}{v_{\rm EW}}
  =
  \left(
  \frac{2F(z)}{21375.716460702053\ldots\,\lambda_F}
  \right)^{1/4}w_H^2 .
\]
Since
\[
  F'(z)=-\frac13e^{-z}z(1+z),
\]
one has \(0<F(z)\leq1\) for \(z\geq0\).  Define
\[
  C_0=
  \left(
  \frac{2}{21375.716460702053\ldots}
  \right)^{1/4}
  =
  0.0983506713707568\ldots .
\]
Then every minimal positive endpoint realization satisfies the density window
\[
  0<
  \frac{Z_{\Gamma,0}}{w_H^2v_{\rm EW}}
  \leq
  C_0\lambda_F^{-1/4}.
\]
Equivalently, a localized D8/O8 calculation passes this row exactly when
\[
  0<
  \frac{21375.716460702053\ldots\,\lambda_F}{2w_H^8}
  \left(\frac{Z_{\Gamma,0}}{v_{\rm EW}}\right)^4
  \leq1,
\]
with \(z\) then fixed by the monotone equation for \(F(z)\).
It is useful to name the left-hand side
\[
  \Xi_\Gamma
  =
  \frac{21375.716460702053\ldots\,\lambda_F}{2w_H^8}
  \left(\frac{Z_{\Gamma,0}}{v_{\rm EW}}\right)^4 .
\]
The value \(\Xi_\Gamma=1\) gives the massless endpoint threshold \(z=0\).
Values \(0<\Xi_\Gamma<1\) give a unique \(z>0\) through
\[
  e^{-z}\left(1+z+\frac{z^2}{3}\right)=\Xi_\Gamma .
\]
Values \(\Xi_\Gamma>1\) reject the minimal positive endpoint row and require
additional localized spectrum data.
Equivalently,
\[
  R_\Gamma
  =
  \Xi_\Gamma^{1/4}
  =
  \frac{Z_{\Gamma,0}}
       {C_0\lambda_F^{-1/4}w_H^2v_{\rm EW}}
  =
  F(z)^{1/4}.
\]
Thus \(R_\Gamma\) is the direct threshold-suppression coordinate for this
localized row.
% CONTROL: The auxiliary ratio chi_H tracks the local Hosotani scale match.
% It belongs to codexVersion-proof-gates.md, the Hosotani auxiliary note,
% and connes-kk-scale-normalization-bridge.md.
An endpoint trace metric therefore gives
\[
  w_y:w_x=Q:1,
\]
and for \(Q=7\) the center is \(t_Q=-7/8\).  The same filtration gives
\[
  \dim H^0(\CP^1,\cO(Q))=Q+1,
\]
the Chan--Paton half-rank used above.
Substitution of \(t_Q\) in the neutral pass equation yields
\[
  9q(q+2)
  =
  2R_{\rm dV}
  \left(
  q\frac{5Q^2+18Q+29}{(Q+1)^2}+8
  \right)
  \left(
  q\frac{Q^2+1}{(Q+1)^2}+1
  \right).
\]
For \(Q=7\), the equation becomes
\[
  576q(q+2)=R_{\rm dV}(25q+32)^2,
\]
with positive root
\[
  q_c=3.511614004347\ldots,\qquad
  \lambda_c=1.873930095907\ldots .
\]
The same point-evaluation specialization has an explicit photon length.
With the primitive relation \(P=Q-1\),
\[
  \omega_T=\eta(t_Q),
  \qquad
  \omega_Y=(Q-1)h+Q\eta(t_Q),
  \qquad
  \omega_Q=(Q-1)h+(Q+1)\eta(t_Q).
\]
For \(t_Q=-Q/(Q+1)\),
\[
  A=2q,\qquad
  B=-\frac{Q-1}{Q+1}q,\qquad
  C=\frac{Q^2+1}{(Q+1)^2}q+1.
\]
In the weak-current specialization \(L_2^2=C\), one obtains
\[
  L_Q^2
  =
  (Q-1)^2A+2(Q-1)(Q+1)B+(Q+1)^2C
  =
  (Q+1)^2L_2^2.
\]
For \(Q=7\), this gives \(L_Q^2=64L_2^2\).  The neutral determinant is
\[
  \det{\cal L}_\Gamma=(Q-1)^2q(q+2),
\]
so the equation
\[
  \det{\cal L}_\Gamma=R_{\rm dV}L_2^2L_Q^2
\]
recovers the reduced \(Q\)-row above.
On this neutral pass surface,
\[
  G_Z^\Gamma
  =
  \frac{L_Q^2}{\det{\cal L}_\Gamma}
  =
  \frac{1}{R_{\rm dV}L_2^2}.
\]
Writing \(C_{\rm W}\) for the Weinberg length-normalization constant in this
local calculation, the charged and neutral source masses read
\[
  M_W^2=\frac{C_{\rm W}^2v_{\rm EW}^2}{4L_2^2},
  \qquad
  M_Z^2=\frac{C_{\rm W}^2v_{\rm EW}^2}{4R_{\rm dV}L_2^2}.
\]
Since \(R_{\rm dV}=\Xp(3/4)/\Xp(2)\), the Schur coefficient becomes
\[
  \mu_\Gamma^2
  =
  \frac{M_Z^2}{\Xp(2)}
  =
  \frac{M_W^2}{\Xp(3/4)}
  =
  \frac{C_{\rm W}^2v_{\rm EW}^2}{4\Xp(3/4)L_2^2}.
\]
The electromagnetic normalization in the same specialization is
\[
  \alpha_\Gamma
  =
  \frac{C_{\rm W}^2}{4\pi L_Q^2}
  =
  \frac{C_{\rm W}^2}{4\pi(Q+1)^2L_2^2}
  =
  \frac{M_W^2}{\pi(Q+1)^2v_{\rm EW}^2}.
\]
For \(Q=7\), this specialization gives
\[
  \alpha_\Gamma=\frac{M_W^2}{64\pi v_{\rm EW}^2}.
\]
Using \(M_W=80.377\,{\rm GeV}\) and
\(v_{\rm EW}=246.21965\,{\rm GeV}\) as representative electroweak inputs,
\[
  \alpha_\Gamma^{-1}=1886.738366864\ldots .
\]
The physical electromagnetic input \(\alpha^{-1}=137.035999084\ldots\)
corresponds to
\[
  k_\alpha=
  \frac{M_W^2}{\pi\alpha v_{\rm EW}^2}
  =
  4.648394337766\ldots,
  \qquad
  \frac{64}{k_\alpha}=13.768195069001\ldots .
\]
Thus the point-evaluation Hodge row is a neutral-matrix filter.  Full
electromagnetic normalization requires an additional localized contribution
to the \(Q\)-current norm.
Equivalently, set \(L_2^2=1\), \(k=k_\alpha\), and write
\[
  \widehat{\cal L}_\Gamma=
  \begin{pmatrix}
    1&\varepsilon\\
    \varepsilon&y
  \end{pmatrix}.
\]
The alpha-closed neutral matrix obeys
\[
  1+y+2\varepsilon=k,
  \qquad
  y-\varepsilon^2=R_{\rm dV}k,
\]
hence
\[
  \varepsilon_\pm=-1\pm\sqrt{k(1-R_{\rm dV})},
  \qquad
  y_\pm=k+1\mp2\sqrt{k(1-R_{\rm dV})}.
\]
For the numerical inputs above,
\[
\begin{array}{c|c|c|c}
{\rm branch}&\varepsilon&y&\varepsilon/\sqrt y\\ \hline
+&0.018362864274\ldots&3.611668609219\ldots&0.009662432587\ldots\\
-&-2.018362864274\ldots&7.685120066313\ldots&-0.728071311237\ldots .
\end{array}
\]
At the point-evaluation neutral root, the corresponding normalized matrix is
\[
  \widehat{\cal L}_{\rm pt}
  =
  \begin{pmatrix}
  1&2.778688215142\ldots\\
  2.778688215142\ldots&57.442623569715\ldots
  \end{pmatrix}.
\]
A sharper decomposition writes each alpha-closed branch \(T_\pm\) as
\[
  T_\pm=\frac{\widehat{\cal L}_{\rm pt}-D_\pm}{s_\pm},
  \qquad
  D_\pm\succeq0,\qquad {\rm rank}\,D_\pm=1 .
\]
The two decompositions are
\[
\begin{array}{c|c|c|c}
{\rm branch}&s&{\rm nonzero\ eigenvalue\ of\ }D&
{\rm slope\ of\ the\ image\ line}\\ \hline
+&0.859516413258733\ldots&
54.478818707582\ldots&
19.667102015222\ldots\\
-&0.672682628439277\ldots&
52.600294175197\ldots&
12.637293988103\ldots .
\end{array}
\]
Thus the localized action data are a rank-one subtraction in the
\(Q\)-current plane together with a common current-frame rescaling.  The local
coefficient row is the extraction of
\(H_{hh}^\Gamma,H_{h\eta}^\Gamma,H_{\eta\eta}^\Gamma\) and \(L_2^2\) from the
same \(d\nu_\Gamma\), endpoint image data, and DBI/CS pullback forms.
A single boundary auxiliary gives the required algebra.  For \(J=(J_3,J_Y)\),
write
\[
  S_{\rm img}^{(2)}
  =
  \frac12J^T\widehat{\cal L}_{\rm pt}J
  +r_{\rm img}\ell_\pm^TJ
  +\frac12m_\pm r_{\rm img}^2 .
\]
Eliminating \(r_{\rm img}\) gives
\[
  T_\pm
  =
  s_\pm^{-1}
  \left(
  \widehat{\cal L}_{\rm pt}
  -\frac{\ell_\pm\ell_\pm^T}{m_\pm}
  \right),
  \qquad
  \frac{\ell_\pm\ell_\pm^T}{m_\pm}
  =
  \delta_\pm
  \binom{1}{\rho_\pm}
  \begin{pmatrix}1&\rho_\pm\end{pmatrix}.
\]
The coefficient extraction target is therefore
\[
\begin{aligned}
+:\quad&
s_+=0.859516413258733\ldots,
\quad
s_+^{-1}=1.163444914575446\ldots,\\
&
\delta_+=0.140483586741\ldots,
\quad
\rho_+=19.667102015222\ldots,\\
-:\quad&
s_-=0.672682628439277\ldots,
\quad
s_-^{-1}=1.486585140930646\ldots,\\
&
\delta_-=0.327317371561\ldots,
\quad
\rho_-=12.637293988103\ldots .
\end{aligned}
\]
Equivalently, for \(c_\rho=(1,\rho)^T\) and
\[
  M_\rho=\widehat{\cal L}_{\rm pt}-\delta(\rho)c_\rho c_\rho^T ,
\]
the equations
\[
  (1,1)M_\rho\binom11=k_\alpha(M_\rho)_{33},
  \qquad
  \det M_\rho=R_{\rm dV}k_\alpha(M_\rho)_{33}^2
\]
give
\[
\begin{aligned}
0=P_\rho(\rho)
={}&
\rho^4
-29.874189325032\ldots\,\rho^3
+170.407256725780\ldots\,\rho^2\\
&+591.897671281551\ldots\,\rho
+93.118987613750\ldots .
\end{aligned}
\]
Positive Schur subtraction selects the roots
\[
  \rho_-=12.637293988103\ldots,
  \qquad
  \rho_+=19.667102015221\ldots .
\]
For positive image mass \(m>0\), the coefficient extraction reads
\[
  \frac{\ell_3^2}{m}=\delta,\qquad
  \frac{\ell_Y}{\ell_3}=\rho,\qquad
  s=(M_\rho)_{33}=1-\delta .
\]
Thus a real localized image auxiliary realizes either
\[
  \delta_+=0.140483586741\ldots,\quad
  \rho_+=19.667102015221\ldots,
\]
or
\[
  \delta_-=0.327317371561\ldots,\quad
  \rho_-=12.637293988103\ldots .
\]
In local variables, choose a normalized image eigenmode \(\psi\) with
\[
  \int_{\Gamma_{\rm EW,int}}d\nu_\Gamma|\psi|^2=1,\qquad
  m_\psi=\langle\psi,{\cal M}_{\rm img}\psi\rangle_\Gamma .
\]
For image current profile \(\Theta_{\rm img}\), define
\[
  {\cal J}_A=\gamma_\Gamma^{-1}(\theta_\Gamma^A,\Theta_{\rm img}),
  \qquad
  \ell_A[\psi]=
  \int_{\Gamma_{\rm EW,int}}d\nu_\Gamma\,\psi{\cal J}_A,
  \qquad A=3,Y .
\]
The extraction invariants are
\[
  \delta_\psi=\frac{\ell_3[\psi]^2}{m_\psi L_2^2},
  \qquad
  \rho_\psi=\frac{\ell_Y[\psi]}{\ell_3[\psi]},
  \qquad
  s_\psi=1-\delta_\psi .
\]
Equivalently, the full image sector gives the Green-kernel correction
\[
  {\cal D}_{AB}
  =
  \frac{1}{L_2^2}
  \left\langle
  {\cal J}_A,
  {\cal M}_{\rm img}^{-1}{\cal J}_B
  \right\rangle_\Gamma,
  \qquad A,B=3,Y .
\]
The one-field Schur row is the rank-one case
\[
  {\cal D}_{33}{\cal D}_{YY}-{\cal D}_{3Y}^2=0,
  \qquad
  {\cal D}_{33}>0 .
\]
A spectral resolution makes this condition computable from the localized
image spectrum.  For
\[
  {\cal M}_{\rm img}\psi_n=m_n\psi_n,
  \qquad m_n>0,
\]
set
\[
  \ell_A^{(n)}
  =
  \int_{\Gamma_{\rm EW,int}}d\nu_\Gamma\,\psi_n{\cal J}_A .
\]
Then
\[
  {\cal D}_{AB}
  =
  \sum_n
  \frac{\ell_A^{(n)}\ell_B^{(n)}}{m_nL_2^2}.
\]
Writing
\[
  u_n
  =
  \frac{1}{L_2\sqrt{m_n}}
  \binom{\ell_3^{(n)}}{\ell_Y^{(n)}},
\]
one has
\[
  {\cal D}=\sum_n u_nu_n^T,
  \qquad
  \Delta_{\rm rank}
  =
  {\cal D}_{33}{\cal D}_{YY}-{\cal D}_{3Y}^2
  =
  \sum_{n<m}
  \left(u_n^3u_m^Y-u_n^Yu_m^3\right)^2 .
\]
Thus rank one is the collinearity condition for the sourced overlap vectors.
For \({\cal D}_{33}>0\), put
\[
  \delta_D={\cal D}_{33},
  \qquad
  \rho_D=\frac{{\cal D}_{3Y}}{{\cal D}_{33}} .
\]
The two Schur branches are selected by
\[
  \Delta_{\rm rank}=0,\qquad
  (\delta_D,\rho_D)
  \in
  \{(\delta_+,\rho_+),(\delta_-,\rho_-)\}.
\]
Several localized image modes can contribute to one branch when the
source-slope ratio is common and the strengths add to \(\delta_\pm\).
A Green-metric form records the same condition directly on the source
profiles.  Define
\[
  (F,G)_G
  =
  \frac{1}{L_2^2}
  \langle F,{\cal M}_{\rm img}^{-1}G\rangle_\Gamma .
\]
Then \({\cal D}_{AB}=({\cal J}_A,{\cal J}_B)_G\).  For
\(\sigma\in\{+,-\}\), the target branch is
\[
  {\cal D}^{(\sigma)}
  =
  \delta_\sigma
  \binom{1}{\rho_\sigma}
  \begin{pmatrix}1&\rho_\sigma\end{pmatrix},
\]
equivalently
\[
  \|{\cal J}_Y-\rho_\sigma{\cal J}_3\|_G^2=0,
  \qquad
  \|{\cal J}_3\|_G^2=\delta_\sigma .
\]
For a positive image Green metric, the local calculation reduces to source
proportionality in the sourced image quotient plus one norm calculation.
If the image source is the pure weak-period line
\[
  {\cal J}_Y=\rho_Q{\cal J}_3,\qquad \rho_Q=2Q+1,
\]
then
\[
  {\cal D}
  =
  \delta_Q
  \binom{1}{\rho_Q}
  \begin{pmatrix}1&\rho_Q\end{pmatrix}.
\]
For \(Q=7\), \(\rho_Q=15\).  The Schur branch equations require
\[
  \rho_D
  \in
  \{12.637293988103\ldots,\;19.667102015221\ldots\}.
\]
The pure weak-period line is therefore excluded as the final image source.
With the alpha-length value,
\[
  \delta_Q
  =
  \frac{64-k_\alpha}{(1+\rho_Q)^2-k_\alpha}
  =
  0.236129805122\ldots,
  \qquad
  {\cal R}(15)=0.999999862211\ldots .
\]
Thus the localized action must deform the source line to one Schur slope or
project the weak-period source away from the Schur image channel.
Using \(\delta_Q\) as the alpha-normalized reference strength, a coframe
deformation in the sourced quotient may be parameterized by
\[
  {\cal J}_3^{(\sigma)}=a_\sigma{\cal J}_3^{(Q)},
  \qquad
  {\cal J}_Y^{(\sigma)}
  =
  a_\sigma(\rho_Q+\zeta_\sigma){\cal J}_3^{(Q)} .
\]
Branch matching gives
\[
  \zeta_\sigma=\rho_\sigma-\rho_Q,
  \qquad
  a_\sigma^2=\frac{\delta_\sigma}{\delta_Q}.
\]
For \(Q=7\),
\[
\begin{array}{c|c|c|c}
{\rm branch}&\zeta_\sigma&a_\sigma^2&a_\sigma\\ \hline
+&4.667102015221\ldots&0.594942204219\ldots&0.771324966677\ldots\\
-&-2.362706011897\ldots&1.386175588431\ldots&1.177359583318\ldots .
\end{array}
\]
These two rows are the local coefficient targets for a sourced D8/O8 coframe
deformation.
Equivalently, let
\[
  {\cal C}_Q=
  \binom{{\cal J}_3^{(Q)}}{{\cal J}_Y^{(Q)}},
  \qquad
  {\cal J}_Y^{(Q)}=\rho_Q{\cal J}_3^{(Q)} .
\]
The deformation can be placed in a lower-triangular source-coframe map
\[
  {\cal C}_\sigma=R_\sigma{\cal C}_Q,
  \qquad
  R_\sigma
  =
  a_\sigma
  \begin{pmatrix}
  1&0\\
  \zeta_\sigma&1
  \end{pmatrix}.
\]
It satisfies
\[
  R_\sigma\binom{1}{\rho_Q}
  =
  a_\sigma\binom{1}{\rho_\sigma},
  \qquad
  R_\sigma{\cal D}_QR_\sigma^T
  =
  {\cal D}_\sigma .
\]
For \(Q=7\),
\[
R_+
=
\begin{pmatrix}
0.771324966677\ldots&0\\
3.599852306368\ldots&0.771324966677\ldots
\end{pmatrix}
\]
and
\[
R_-
=
\begin{pmatrix}
1.177359583318\ldots&0\\
-2.781754565670\ldots&1.177359583318\ldots
\end{pmatrix}.
\]
The diagonal entry is the Green-norm rescaling, and the lower-left entry is
the localized \(T_3\)-to-\(Y\) source mixing required from the D8/O8 pullback.
The equivalent pullback condition on localized source coframes is
\[
  \theta_{\Gamma,\sigma}^{3}
  =
  a_\sigma\theta_{\Gamma,Q}^{3},
  \qquad
  \theta_{\Gamma,\sigma}^{Y}
  =
  a_\sigma
  \left(
  \theta_{\Gamma,Q}^{Y}
  +\zeta_\sigma\theta_{\Gamma,Q}^{3}
  \right)
\]
inside the image source quotient.  Thus the localized action returns
\[
  Z_{\rm img}^{(\sigma)}=a_\sigma^2Z_{\rm img}^{(Q)},
  \qquad
  \beta_{Y3}^{(\sigma)}=\zeta_\sigma .
\]
For \(Q=7\),
\[
\begin{array}{c|c|c}
{\rm branch}&Z_{\rm img}^{(\sigma)}/Z_{\rm img}^{(Q)}
&\beta_{Y3}^{(\sigma)}\\ \hline
+&0.594942204219\ldots&4.667102015221\ldots\\
-&1.386175588431\ldots&-2.362706011897\ldots .
\end{array}
\]
For a computed localized action, define
\[
  \widehat\delta_G
  =
  ({\cal J}_3,{\cal J}_3)_G={\cal D}_{33},
  \qquad
  \widehat\rho_G
  =
  \frac{({\cal J}_3,{\cal J}_Y)_G}
       {({\cal J}_3,{\cal J}_3)_G}
  =
  \frac{{\cal D}_{3Y}}{{\cal D}_{33}},
\]
and
\[
  \widehat\epsilon_G^2
  =
  \frac{\|{\cal J}_Y-\widehat\rho_G{\cal J}_3\|_G^2}
       {\widehat\delta_G}
  =
  \frac{{\cal D}_{33}{\cal D}_{YY}-{\cal D}_{3Y}^2}
       {{\cal D}_{33}^2}.
\]
The Schur-branch extraction condition is
\[
  (\widehat\delta_G,\widehat\rho_G,\widehat\epsilon_G^2)
  =
  (\delta_\sigma,\rho_\sigma,0),
  \qquad \sigma\in\{+,-\}.
\]
The deformation coefficients are then
\[
  a_\sigma^2=\frac{\widehat\delta_G}{\delta_Q},
  \qquad
  \zeta_\sigma=\widehat\rho_G-\rho_Q .
\]
In the pullback notation the direct readout is
\[
  \widehat Z_{\rm img}= {\cal D}_{33},
  \qquad
  \widehat\beta_{Y3}=\frac{{\cal D}_{3Y}}{{\cal D}_{33}}-\rho_Q,
\]
\[
  \widehat E_\perp
  =
  {\cal D}_{YY}-\frac{{\cal D}_{3Y}^2}{{\cal D}_{33}} .
\]
The branch extraction condition becomes
\[
  \frac{\widehat Z_{\rm img}}{Z_{\rm img}^{(Q)}}=a_\sigma^2,
  \qquad
  \widehat\beta_{Y3}=\zeta_\sigma,
  \qquad
  \widehat E_\perp=0 .
\]
A projected-image realization uses a \(G\)-self-adjoint idempotent
\(\Pi_\sigma\) satisfying
\[
  \Pi_\sigma^2=\Pi_\sigma,
  \qquad
  (\Pi_\sigma F,G)_G=(F,\Pi_\sigma G)_G,
\]
\[
  \Pi_\sigma{\cal J}_Y
  =
  \rho_\sigma\Pi_\sigma{\cal J}_3,
  \qquad
  \|\Pi_\sigma{\cal J}_3\|_G^2=\delta_\sigma .
\]
The coframe map also defines a symmetric source-index block
\[
  {\cal H}_\sigma=R_\sigma R_\sigma^T
  =
  a_\sigma^2
  \begin{pmatrix}
  1&\zeta_\sigma\\
  \zeta_\sigma&1+\zeta_\sigma^2
  \end{pmatrix}.
\]
It has
\[
  \det{\cal H}_\sigma=a_\sigma^4,
  \qquad
  \chi_\sigma
  =
  \frac{\zeta_\sigma}{\sqrt{1+\zeta_\sigma^2}} .
\]
For \(Q=7\),
\[
{\cal H}_+
=
\begin{pmatrix}
0.594942204219\ldots&2.776655960251\ldots\\
2.776655960251\ldots&13.553878831880\ldots
\end{pmatrix},
\qquad
\chi_+=0.977806418243\ldots ,
\]
and
\[
{\cal H}_-
=
\begin{pmatrix}
1.386175588431\ldots&-3.275125396331\ldots\\
-3.275125396331\ldots&9.124334052058\ldots
\end{pmatrix},
\qquad
\chi_-=-0.920912188652\ldots .
\]
The same rank-one image subtraction can be written with one real localized
field \(r_\sigma\):
\[
  S_{{\rm img},\sigma}^{(2)}
  =
  \frac12m_\sigma r_\sigma^2
  +L_2\sqrt{m_\sigma}\,
  r_\sigma\,g_\sigma^TJ,
  \qquad
  J=\binom{J_3}{J_Y}.
\]
With
\[
  g_\sigma=\sqrt{\delta_\sigma}\binom{1}{\rho_\sigma},
\]
eliminating \(r_\sigma\) gives
\[
  -\frac12L_2^2J^Tg_\sigma g_\sigma^TJ,
  \qquad
  {\cal D}_\sigma=g_\sigma g_\sigma^T .
\]
For the two branches,
\[
\begin{array}{c|c|c}
{\rm branch}&(g_\sigma)_3&(g_\sigma)_Y\\ \hline
+&0.374811401563\ldots&7.371454070999\ldots\\
-&0.572116571654\ldots&7.230005311453\ldots .
\end{array}
\]
Return to the local source definition
\[
  {\cal J}_A
  =
  \gamma_\Gamma^{-1}(\theta_\Gamma^A,\Theta_{\rm img}),
  \qquad A=3,Y .
\]
For a normalized image eigenmode \(\psi_\sigma\),
\[
  \int_{\Gamma_{\rm EW,int}}d\nu_\Gamma|\psi_\sigma|^2=1,
  \qquad
  {\cal M}_{\rm img}\psi_\sigma=m_\sigma\psi_\sigma ,
\]
the one-field packet is equivalent to
\[
  \int_{\Gamma_{\rm EW,int}}d\nu_\Gamma\,
  \psi_\sigma\,
  \gamma_\Gamma^{-1}(\theta_\Gamma^A,\Theta_{{\rm img},\sigma})
  =
  L_2\sqrt{m_\sigma}\,(g_\sigma)_A,
  \qquad A=3,Y .
\]
With \(\Xi_\sigma=\Theta_{{\rm img},\sigma}/\sqrt{m_\sigma}\), this reads
\[
  \frac{1}{L_2}
  \int_{\Gamma_{\rm EW,int}}d\nu_\Gamma\,
  \psi_\sigma\,
  \gamma_\Gamma^{-1}(\theta_\Gamma^A,\Xi_\sigma)
  =
  (g_\sigma)_A .
\]
A sufficient local profile is obtained from
\[
  {\cal G}_{AB}^{(\psi)}
  =
  \int_{\Gamma_{\rm EW,int}}d\nu_\Gamma\,
  \psi_\sigma^2\,
  \gamma_\Gamma^{-1}(\theta_\Gamma^A,\theta_\Gamma^B).
\]
When this matrix is invertible,
\[
  \Xi_\sigma
  =
  L_2\psi_\sigma\,
  c_\sigma^B\theta_{\Gamma B},
  \qquad
  c_\sigma^B
  =
  \left({\cal G}^{(\psi)-1}\right)^{BA}(g_\sigma)_A
\]
solves the two overlap equations.  The D8/O8 source row is therefore tested by
\({\cal G}^{(\psi)}\) and the Gram-dual current combination in the
sourced image quotient.
Writing
\[
  {\cal G}^{(\psi)}
  =
  \begin{pmatrix}
  A_\psi&B_\psi\\
  B_\psi&C_\psi
  \end{pmatrix},
  \qquad
  \Delta_\psi=A_\psi C_\psi-B_\psi^2,
\]
the coefficients are
\[
  c_\sigma^3
  =
  \frac{C_\psi(g_\sigma)_3-B_\psi(g_\sigma)_Y}{\Delta_\psi},
  \qquad
  c_\sigma^Y
  =
  \frac{-B_\psi(g_\sigma)_3+A_\psi(g_\sigma)_Y}{\Delta_\psi}.
\]
In an orthonormal weighted current frame,
\({\cal G}^{(\psi)}=\mathbf 1\), hence
\[
  \Xi_\sigma
  =
  L_2\psi_\sigma
  \left((g_\sigma)_3\theta_\Gamma^3
  +(g_\sigma)_Y\theta_\Gamma^Y\right).
\]
The general solution splits as
\[
  \Xi_\sigma
  =
  \Xi_{\sigma,\min}+\Xi_{\sigma,\perp},
  \qquad
  \int_{\Gamma_{\rm EW,int}}d\nu_\Gamma\,
  \psi_\sigma\,
  \gamma_\Gamma^{-1}(\theta_\Gamma^A,\Xi_{\sigma,\perp})=0 .
\]
The forced profile has norm
\[
  \frac{\|\Xi_{\sigma,\min}\|_\Gamma^2}{L_2^2}
  =
  g_\sigma^T{\cal G}^{(\psi)-1}g_\sigma .
\]
In an orthonormal weighted current frame this gives
\[
\begin{array}{c|c}
{\rm branch}&\|\Xi_{\sigma,\min}\|_\Gamma^2/L_2^2\\ \hline
+&54.478818707582\ldots\\
-&52.600294175197\ldots .
\end{array}
\]
Any nonzero \(\Xi_{\sigma,\perp}\) raises the image-profile norm without
changing the Schur matrix.
Define the profile residual
\[
  {\cal R}_{\Xi,\sigma}
  =
  \frac{\|\Xi_{\sigma,\perp}\|_\Gamma^2}
       {\|\Xi_{\sigma,\min}\|_\Gamma^2}.
\]
The minimal image-source realization has
\[
  {\cal R}_{\Xi,\sigma}=0 .
\]
A positive value records extra localized image content that leaves
\({\cal D}_\sigma\) invariant and belongs to the D8/O8 spectrum.

\begin{proposition}[One-mode image-source closure]
Fix \(\sigma\in\{+,-\}\) and fix the orientation of the real image field
\(r_\sigma\).  A positive one-mode image sector over \(\Gamma_{\rm EW}\)
realizes the Schur branch
\[
  {\cal D}_\sigma
  =
  \delta_\sigma
  \binom{1}{\rho_\sigma}
  \begin{pmatrix}1&\rho_\sigma\end{pmatrix}
\]
precisely when the mass-normalized image profile satisfies
\[
  \frac{1}{L_2}
  \int_{\Gamma_{\rm EW,int}}d\nu_\Gamma\,
  \psi_\sigma\,
  \gamma_\Gamma^{-1}(\theta_\Gamma^A,\Xi_\sigma)
  =
  (g_\sigma)_A,
  \qquad A=3,Y .
\]
\end{proposition}

\begin{proof}
The overlap vector is
\[
  \ell_A
  =
  \int_{\Gamma_{\rm EW,int}}d\nu_\Gamma\,
  \psi_\sigma\,
  \gamma_\Gamma^{-1}(\theta_\Gamma^A,\Theta_{{\rm img},\sigma})
  =
  L_2\sqrt{m_\sigma}(g_\sigma)_A .
\]
Integrating out \(r_\sigma\) gives
\[
  {\cal D}_{AB}
  =
  \frac{\ell_A\ell_B}{m_\sigma L_2^2}
  =
  (g_\sigma)_A(g_\sigma)_B .
\]
Since \(g_\sigma=\sqrt{\delta_\sigma}(1,\rho_\sigma)^T\), the displayed branch
matrix follows.  The reverse implication follows by choosing the orientation of
\(r_\sigma\) so that \(\ell_3>0\).
\end{proof}

The localized realization data are the packet
\[
  {\cal P}_{\Gamma,\sigma}
  =
  \left(
  d\nu_\Gamma,\,
  \psi_\sigma,\,
  m_\sigma,\,
  \Xi_\sigma,\,
  {\cal G}^{(\psi)},\,
  Z_{\Gamma,0},\,
  {\cal N}_\Gamma^\chi
  \right)
\]
with
\[
  \int_{\Gamma_{\rm EW,int}}d\nu_\Gamma|\psi_\sigma|^2=1,
  \qquad
  {\cal M}_{\rm img}\psi_\sigma=m_\sigma\psi_\sigma,
  \qquad
  m_\sigma>0,
\]
\[
  L_2^{-1}
  \int_{\Gamma_{\rm EW,int}}d\nu_\Gamma\,
  \psi_\sigma
  \gamma_\Gamma^{-1}(\theta_\Gamma^A,\Xi_\sigma)
  =
  (g_\sigma)_A,
  \qquad A=3,Y,
\]
\[
  {\cal R}_{\Xi,\sigma}=0,
  \qquad
  Z_{\Gamma,H}=\frac Q4Z_{\Gamma,0},
  \qquad
  {\cal N}_\Gamma^\chi=2\sigma_\chi d_\chi F(z).
\]

\subsection{Localized Quadratic Master Form}

Let
\[
  \Gamma_{\rm int}=\CP^2\times S^1_Q .
\]
Let \({\cal S}_\Gamma\) label the D8 image sectors retained by the O8
projection, and set
\[
  \kappa_8=T_8(2\pi\alpha')^2 .
\]
The string-frame DBI expansion on sector \(s\),
\[
  S_{{\rm DBI},s}
  =
  -T_8\int e^{-\phi_s}
  {\rm Str}_{E_s}
  \sqrt{-\det\bigl(P_s[G+B]+2\pi\alpha'F_s\bigr)},
\]
has quadratic Yang--Mills density
\[
  d\nu_{\Gamma,s}
  =
  \kappa_8\,\omega_s\,e^{-\phi_s}
  \sqrt{\gamma_{\Gamma,s}}\,
  d^5\sigma ,
\]
where \(\omega_s\) is the orientifold and endpoint multiplicity assigned to
the retained sector.  The common \(\Gamma\)-inner product is
\[
  \langle X,Y\rangle_\Gamma
  =
  \sum_{s\in{\cal S}_\Gamma}
  \int_{\Gamma_{\rm int}^{(s)}}d\nu_{\Gamma,s}\,
  \gamma_s^{mn}\,
  {\rm tr}_{E_s}
  \bigl(\Pi_sX_{m,s}\Pi_sY_{n,s}\bigr).
\]
The scalar density \(d\nu_\Gamma\) used below denotes this sector-sum measure
together with its finite trace projectors.  The immediate DBI readouts are
\[
  {\cal K}_{AB}^\Gamma
  =
  \langle\theta_\Gamma^A,\theta_\Gamma^B\rangle_\Gamma,
  \qquad
  {\cal G}_{AB}^{(\psi)}
  =
  \langle\psi_\sigma\theta_\Gamma^A,
  \psi_\sigma\theta_\Gamma^B\rangle_\Gamma,
\]
\[
  Z_{\Gamma,0}=\langle\eta_0,\eta_0\rangle_\Gamma,
  \qquad
  Z_{\Gamma,H}
  =
  \langle\omega_H n_H^aT_a\eta_0,
  \omega_H n_H^bT_b\eta_0\rangle_\Gamma .
\]
The image operator is an adjoint square,
\[
  {\cal M}_{\rm img}={\cal Q}_{\Gamma,\sigma}^\dagger
  {\cal Q}_{\Gamma,\sigma},
  \qquad
  \dagger \text{ taken with } \langle\,,\,\rangle_\Gamma .
\]
The weak-period Chan--Paton half-stack is
\[
  W_Q=H^0(\CP^1,\mathcal O(Q)),
  \qquad
  \dim W_Q=Q+1 .
\]
The local D8 endpoint space is the orientifold-real double
\[
  E_Q=(W_Q\oplus W_Q^\vee)^\Omega,
  \qquad
  {\rm rk}_{\mathbb C}(E_Q^{\mathbb C})=2(Q+1).
\]
At the active value \(Q=7\),
\[
  {\rm rk}_{\mathbb C}(E_Q^{\mathbb C})=16,
  \qquad
  q_8^{\rm net}=16-16=0,
  \qquad
  k=\frac{16}{2}-8=0 .
\]
Thus the retained endpoint is the locally neutral D8/O8 row in the doubled
type-I' convention.

Introduce row projectors
\[
  \Pi_Q=P_\Omega P_{E_Q},
  \qquad
  \Pi_\sigma=|e_\sigma\rangle\langle e_\sigma|,
  \qquad
  \Pi_\chi=P_1^\chi+P_2^\chi ,
\]
with
\[
  P_1^\chi={\rm proj}_{\rm span\{u,v\}},
  \qquad
  P_2^\chi={\rm proj}_{\rm span\{u^2,v^2\}} .
\]
For a row projector \(\Pi\), define
\[
  \langle X,Y\rangle_{\Gamma,\Pi}
  =
  \int_{\Gamma_{\rm int}}d\mu_\Gamma\,
  \gamma_\Gamma^{mn}\,
  {\rm tr}_{{\cal H}_\Gamma^{\rm loc}}
  \bigl(\Pi X_m\Pi Y_n\bigr),
  \qquad
  d\mu_\Gamma=\kappa_8e^{-\phi}\sqrt{\gamma_\Gamma}\,d^5\sigma .
\]
The row readouts become
\[
  {\cal K}_{AB}^\Gamma
  =
  \langle\theta_\Gamma^A,\theta_\Gamma^B\rangle_{\Gamma,\Pi_Q},
  \qquad
  {\cal G}_{AB}^{(\psi)}
  =
  \langle\psi_\sigma\theta_\Gamma^A,
  \psi_\sigma\theta_\Gamma^B\rangle_{\Gamma,\Pi_\sigma},
\]
\[
  Z_{\Gamma,0}=\langle\eta_0,\eta_0\rangle_{\Gamma,\Pi_Q},
  \qquad
  Z_{\Gamma,H}
  =
  \langle\omega_H n_H^aT_a\eta_0,
  \omega_H n_H^bT_b\eta_0\rangle_{\Gamma,\Pi_Q},
\]
and
\[
  {\cal N}_\Gamma^\chi
  =
  {\rm Str}_{\Pi_\chi\ker{\cal B}_\Gamma^\chi}
  (\tau_\alpha).
\]
Let
\[
  {\cal H}_n=H^0(\CP^1,\mathcal O(n)),
  \qquad
  e_k=u^{n-k}v^k,\quad k=0,\ldots,n .
\]
The \(\SU(2)\)-equivariant antiunitary
\[
  J_n e_k=(-1)^k e_{n-k}
\]
satisfies
\[
  J_n^2=(-1)^n .
\]
Thus \(n=1\) gives the pseudoreal doublet and \(n=2\) gives the real triplet.
Using the invariant pairing to identify \({\cal H}_n^\vee\) with
\({\cal H}_n\), the doubled endpoint involution is
\[
  \Omega_n(s,\lambda)=(J_n^{-1}\lambda,J_ns),
  \qquad
  \Omega_n^2=1,
  \qquad
  P_{\Omega,n}=\frac12(1+\Omega_n).
\]
On the charged bases,
\[
  J_1u=v,\qquad J_1v=-u,
\]
and
\[
  J_2u^2=v^2,\qquad J_2uv=-uv,\qquad J_2v^2=u^2 .
\]
Consequently \(P_{\Omega,1}\) preserves the \((u,v)\) charged doublet, and
\(P_{\Omega,2}\) preserves the charged pair \((u^2,v^2)\) after applying
\(P_2^\chi=1-|uv\rangle\langle uv|\).  The trace readout is
\[
  {\rm Tr}\,P_{\Omega,1}P_1^\chi=2,
  \qquad
  {\rm Tr}\,P_{\Omega,2}P_2^\chi=2 .
\]
The orientifold row therefore contributes
\[
  \alpha_{\rm parity}=0,
  \qquad
  d_1=d_2=d_\chi .
\]
The determinant sign \(\sigma_\chi\) and the common threshold \(M_\chi\) are
properties of the localized oscillator \({\cal B}_\Gamma^\chi\).
Let \(x=w_H\theta_H\).  On the charged Borel--Weil pair with charges
\[
  r=1,\qquad r=2 ,
\]
take the localized oscillator block to have eigenvalues
\[
  \lambda_{\ell,r}(x)
  =
  p_E^2+M_\chi^2+
  \frac{(2\pi\ell+r x)^2}{L_H^2},
  \qquad
  \ell\in\mathbb Z,
\]
on \(P_r^\chi\otimes E_\chi\).  The Euclidean complex-boson Gaussian gives
the positive one-loop sign
\[
  \sigma_\chi=+1.
\]
After Schwinger parameterization and Poisson resummation, the
angle-dependent term is
\[
  V_\chi(x)
  =
  \frac{3}{4\pi^2L_H^4}\,2d_\chi
  \sum_{m\ge1}
  \frac{F(mz_\chi)}{m^5}
  \bigl(\cos(mx)+\cos(2mx)\bigr),
  \qquad
  z_\chi=M_\chi L_H ,
\]
with
\[
  F(z)=e^{-z}\left(1+z+\frac{z^2}{3}\right).
\]
Thus the first-wrapping determinant density is
\[
  {\cal N}_\Gamma^\chi=2d_\chi F(z_\chi).
\]
For \(d_\chi=1\),
\[
  {\cal N}_\Gamma^\chi=2F(z_\chi),
  \qquad
  M_\chi=L_H^{-1}F^{-1}
  \left(\frac{{\cal N}_\Gamma^\chi}{2}\right),
\]
with \(F^{-1}\) taken on \(0<F\le1\).
Let \(\Pi_A\) denote projection of the pulled-back RR one-form to the retained
current generator \(A\in\{3,Y\}\).  Split the localized coframe as
\[
  \theta_\Gamma^A
  =
  \theta_{\Gamma,{\rm CP}}^A
  +
  q_A^{-1}\Pi_A\iota_\Gamma^*C_1,
  \qquad A=3,Y,
\]
with curvature check
\[
  d_\Gamma\theta_\Gamma^A
  =
  d_\Gamma\theta_{\Gamma,{\rm CP}}^A
  +
  q_A^{-1}\Pi_A\iota_\Gamma^*F_2 .
\]
The weak triplet is the equivariant completion of \(\theta_\Gamma^3\), and
the unit-current pass is
\[
  \langle\theta_\Gamma^a,\theta_\Gamma^b\rangle_{\Gamma,\Pi_Q}
  =
  L_{2,\Gamma}^2\delta^{ab},
  \qquad
  e_\Gamma^a=L_{2,\Gamma}^{-1}\theta_\Gamma^a .
\]
The neutral current matrix is
\[
  {\cal L}_\Gamma
  =
  \begin{pmatrix}
  K_{33}^\Gamma&K_{3Y}^\Gamma\\
  K_{3Y}^\Gamma&K_{YY}^\Gamma
  \end{pmatrix},
  \qquad
  K_{AB}^\Gamma
  =
  \langle\theta_\Gamma^A,\theta_\Gamma^B\rangle_{\Gamma,\Pi_Q}.
\]
The photon coframe and length are
\[
  \theta_\Gamma^Q=\theta_\Gamma^3+\theta_\Gamma^Y,
  \qquad
  L_Q^2=(1,1){\cal L}_\Gamma\binom{1}{1}.
\]
The neutral branch test becomes
\[
  \det{\cal L}_\Gamma
  =
  R_{\rm dV}\,K_{33}^\Gamma L_Q^2,
  \qquad
  R_{\rm dV}=\frac{\Xp(3/4)}{\Xp(2)}.
\]
In this row the vertical contribution is read directly as
\[
  L_{2,{\rm vert}}^2
  =
  q_3^{-2}
  \langle\Pi_3\iota_\Gamma^*C_1,
  \Pi_3\iota_\Gamma^*C_1\rangle_{\Gamma,\Pi_Q}.
\]
The Chern--Simons source for this one-form is
\[
  S_{\rm CS}^{\Gamma}
  =
  \mu_8\sum_{s\in{\cal S}_\Gamma}\epsilon_s
  \int_{\Gamma_{\rm EW}^{(s)}}
  \sum_r C_{9-2r}\wedge
  \left[{\cal Y}_{\Gamma,s}\right]_{2r},
\]
with
\[
  {\cal Y}_{\Gamma,s}
  =
  {\rm ch}({\cal F}_s)
  \sqrt{\frac{\widehat A(T\Gamma_s)}{\widehat A(N\Gamma_s)}} .
\]
On the spin\(^c\) \(\CP^2\) boundary row,
\[
  {\cal F}_Q/(2\pi)=f_mh_Q,
  \qquad
  f_m=m+\frac{13}{2},
  \qquad
  \int_{\CP^2}h_Q^2=1,
\]
and
\[
  {\cal Y}_{\Gamma}
  =
  N_8\left(
  1+f_mh_Q+
  \left(\frac{f_m^2}{2}-\frac{1}{16}\right)h_Q^2
  \right).
\]
Hence the \(C_1\) variation gives the boundary source coefficient
\[
  \mathfrak q_{C_1}
  =
  N_8\left(\frac{f_m^2}{2}-\frac{1}{16}\right),
  \qquad
  \delta_{C_1}S_{\rm CS}^{\Gamma}
  =
  \mu_8\int_{M_4\times S^1_Q}
  \delta C_{1,Q}\,\mathfrak q_{C_1}.
\]
For the minimal spin\(^c\) local stack \(m=0,\ N_8=16\),
\[
  \mathfrak q_{C_1}=337.
\]
Define the total projected one-form source by adding the orientifold and bulk
counterterms:
\[
  {\cal I}_{1,A}^{\Gamma}
  =
  \Pi_A(\mathfrak q_{C_1}\,ds_Q)
  +{\cal I}_{1,A}^{\rm O8}
  +{\cal I}_{1,A}^{\rm bulk}.
\]
Compact solvability imposes the zero-mode equation
\[
  \langle{\cal I}_{1,A}^{\Gamma},\omega_Q\rangle_\Gamma=0,
  \qquad
  \omega_Q=L_Q^{-1}ds_Q .
\]
The Green operator acts on
\[
  {\cal I}_{1,A}^{\perp}
  =
  {\cal I}_{1,A}^{\Gamma}
  -
  \frac{\langle{\cal I}_{1,A}^{\Gamma},\omega_Q\rangle_\Gamma}
       {\langle\omega_Q,\omega_Q\rangle_\Gamma}\omega_Q .
\]
The projected Green solve is
\[
  \Pi_A\iota_\Gamma^*C_1
  =
  \Pi_A C_{1,{\rm harm}}
  +
  \mu_8\,G_{1,\Gamma}{\cal I}_{1,A}^{\perp},
  \qquad A=3,Y,
\]
and therefore
\[
  L_{2,{\rm vert}}^2
  =
  q_3^{-2}
  \left\|
  \Pi_3 C_{1,{\rm harm}}
  +
  \mu_8\,G_{1,\Gamma}{\cal I}_{1,3}^{\perp}
  \right\|_{\Gamma,\Pi_Q}^2 .
\]
Choose a \(\Gamma\)-orthonormal eigenbasis
\(\{\alpha_i\}_{i\ge1}\) of the zero-mode-orthogonal retained one-form
sector:
\[
  \langle\alpha_i,\alpha_j\rangle_{\Gamma,\Pi_Q}=\delta_{ij},
  \qquad
  \Delta_{1,\Gamma}\alpha_i=\lambda_i\alpha_i,
  \qquad
  \lambda_i>0 .
\]
Set
\[
  B_{Ai}=\langle\alpha_i,{\cal I}_{1,A}^{\perp}\rangle_{\Gamma,\Pi_Q},
  \qquad
  P_{Ai}=\langle\alpha_i,\theta_{\Gamma,{\rm CP}}^A\rangle_{\Gamma,\Pi_Q},
\]
and
\[
  \Pi_A C_{1,{\rm harm}}=h_A\omega_Q,
  \qquad
  \Omega_Q=\langle\omega_Q,\omega_Q\rangle_{\Gamma,\Pi_Q},
  \qquad
  p_A=\langle\theta_{\Gamma,{\rm CP}}^A,\omega_Q\rangle_{\Gamma,\Pi_Q}.
\]
Then
\[
  G_{1,\Gamma}{\cal I}_{1,A}^{\perp}
  =
  \sum_{i\ge1}\frac{B_{Ai}}{\lambda_i}\alpha_i ,
\]
and the retained coframe is
\[
  \theta_\Gamma^A
  =
  \theta_{\Gamma,{\rm CP}}^A
  +q_A^{-1}h_A\omega_Q
  +q_A^{-1}\mu_8
  \sum_{i\ge1}\frac{B_{Ai}}{\lambda_i}\alpha_i .
\]
Thus the current matrix entries are
\[
\begin{aligned}
  K_{AB}^\Gamma
  &=
  K_{AB}^{\rm CP}
  +q_A^{-1}h_Ap_B
  +q_B^{-1}h_Bp_A
  +q_A^{-1}q_B^{-1}h_Ah_B\Omega_Q\\
  &\quad
  +\mu_8 q_A^{-1}\sum_i\frac{B_{Ai}P_{Bi}}{\lambda_i}
  +\mu_8 q_B^{-1}\sum_i\frac{B_{Bi}P_{Ai}}{\lambda_i}
  +\mu_8^2q_A^{-1}q_B^{-1}
  \sum_i\frac{B_{Ai}B_{Bi}}{\lambda_i^2},
\end{aligned}
\]
where
\[
  K_{AB}^{\rm CP}
  =
  \langle\theta_{\Gamma,{\rm CP}}^A,
  \theta_{\Gamma,{\rm CP}}^B\rangle_{\Gamma,\Pi_Q}.
\]
On the product zero-mode reduction of the boundary circle, let
\(\varphi\in[0,2\pi]\) parameterize \(S^1_Q\), with
\[
  ds_Q=R_Q\,d\varphi,
  \qquad
  L_Q=2\pi R_Q .
\]
After the \(\CP^2\) zero-mode projection, write
\[
  {\cal I}_{1,A}^{\Gamma}
  =
  \rho_A(\varphi)\,ds_Q .
\]
The compact condition removes the average
\[
  \bar\rho_A=\frac1{2\pi}\int_0^{2\pi}\rho_A(\varphi)\,d\varphi,
  \qquad
  \rho_A^\perp(\varphi)=\rho_A(\varphi)-\bar\rho_A .
\]
With
\[
  \rho_A^\perp(\varphi)
  =
  \sum_{n\ge1}
  \left(a_{A,n}\cos n\varphi+b_{A,n}\sin n\varphi\right),
\]
the retained one-form eigenvalues are
\[
  \lambda_n=\frac{n^2}{R_Q^2}
  =
  \left(\frac{2\pi n}{L_Q}\right)^2 .
\]
In an orthonormal Fourier basis, the Green part of the current matrix is
\[
  K_{AB}^{\rm Green}
  =
  \mu_8^2q_A^{-1}q_B^{-1}
  \sum_{n\ge1}
  \frac{a_{A,n}a_{B,n}+b_{A,n}b_{B,n}}{\lambda_n^2}.
\]
The constant spin\(^c\) stack source
\[
  \rho_A(\varphi)=\Pi_A\mathfrak q_{C_1}
\]
has \(\rho_A^\perp=0\) and contributes through the tadpole equation.  Its
Green-current contribution is
\[
  K_{AB}^{\rm Green}=0 .
\]
Thus a nonzero sourced vertical current from the CS row requires Fourier
content in
\[
  {\cal I}_{1,A}^{\rm O8}+{\cal I}_{1,A}^{\rm bulk}
\]
or a harmonic coefficient \(h_A\) in \(\Pi_A C_{1,{\rm harm}}\).
Let
\[
  v_A=q_A^{-1}h_A,\qquad A=3,Y,
\]
and collect the zero-mode-orthogonal part into
\[
  \widehat C_{AB}
  =
  K_{AB}^{\rm CP}
  +K_{AB}^{\rm Green}
  +K_{AB}^{\rm CP-Green}.
\]
With
\[
  p_A=\langle\theta_{\Gamma,{\rm CP}}^A,\omega_Q\rangle_{\Gamma,\Pi_Q},
  \qquad
  \Omega_Q=\langle\omega_Q,\omega_Q\rangle_{\Gamma,\Pi_Q},
\]
the neutral current matrix is
\[
  {\cal L}_{AB}^{\Gamma}
  =
  \widehat C_{AB}
  +v_Ap_B+v_Bp_A+\Omega_Qv_Av_B .
\]
Set
\[
  C_Q=(1,1)\widehat C\binom{1}{1},
  \qquad
  p_Q=p_3+p_Y,
  \qquad
  v_Q=v_3+v_Y .
\]
Then the branch-current closure condition becomes
\[
  \det{\cal L}^{\Gamma}
  =
  R_{\rm dV}
  \left(\widehat C_{33}+2v_3p_3+\Omega_Qv_3^2\right)
  \left(C_Q+2v_Qp_Q+\Omega_Qv_Q^2\right).
\]
For the product specialization \(p_A=0\) with vanishing perpendicular Green
contribution,
\[
  {\cal L}^{\Gamma}=C+\Omega_Qvv^T,
\]
and therefore
\[
\begin{aligned}
  \det C
  &+\Omega_Q
  \left(C_{YY}v_3^2-2C_{3Y}v_3v_Y+C_{33}v_Y^2\right)\\
  &=
  R_{\rm dV}
  \left(C_{33}+\Omega_Qv_3^2\right)
  \left(C_Q+\Omega_Q(v_3+v_Y)^2\right).
\end{aligned}
\]
This specialization gives
\[
  L_{2,{\rm vert}}^2=\Omega_Qv_3^2 .
\]
For a fixed harmonic direction
\[
  t=\frac{v_Y}{v_3},
  \qquad
  u=\Omega_Qv_3^2,
\]
define
\[
  \Delta_C=\det C,
  \qquad
  C_Q=C_{33}+2C_{3Y}+C_{YY},
  \qquad
  C_N=C_{33}-2C_{3Y}+C_{YY}.
\]
The same branch equation becomes
\[
  A_tu^2+B_tu+D_C=0,
\]
where
\[
  A_t=R_{\rm dV}(1+t)^2,
\]
\[
  B_t=
  R_{\rm dV}\left(C_{33}(1+t)^2+C_Q\right)
  -\left(C_{YY}-2tC_{3Y}+t^2C_{33}\right),
\]
and
\[
  D_C=R_{\rm dV}C_{33}C_Q-\Delta_C .
\]
For \(t\ne-1\),
\[
  u_\pm(t)=
  \frac{-B_t\pm\sqrt{B_t^2-4A_tD_C}}{2A_t}.
\]
The massive-neutral harmonic alignment \(v_Y=-v_3\) has \(v_Q=0\) and
\[
  L_Q^2=C_Q,
  \qquad
  u_N=\frac{D_C}{(1-R_{\rm dV})C_Q}.
\]
The \(Q\)-parallel alignment \(v_Y=v_3\) obeys
\[
  4R_{\rm dV}u^2+
  \left(R_{\rm dV}(4C_{33}+C_Q)-C_N\right)u+D_C=0 .
\]
Normalize the circle one-form by
\[
  \int_{S^1_Q}\omega_Q=1 .
\]
The harmonic RR potential is parametrized by Wilson-line variables
\[
  \Pi_A C_{1,{\rm harm}}=2\pi\vartheta_A\omega_Q,
  \qquad
  \vartheta_A\in\mathbb R/\mathbb Z,
  \qquad A=3,Y.
\]
Consequently
\[
  v_A=2\pi q_A^{-1}\vartheta_A,
\]
and, for \(v_3\ne0\),
\[
  t=\frac{q_3\vartheta_Y}{q_Y\vartheta_3},
  \qquad
  u=(2\pi)^2\Omega_Q q_3^{-2}\vartheta_3^2 .
\]
The massive-neutral harmonic row is
\[
  q_3^{-1}\vartheta_3+q_Y^{-1}\vartheta_Y=0,
\]
and the \(Q\)-parallel harmonic row is
\[
  q_3^{-1}\vartheta_3-q_Y^{-1}\vartheta_Y=0 .
\]
Retaining the CP--circle cross vector \(p=(p_3,p_Y)\) gives a closed
photon-preserving row.  Put
\[
  v=a\binom{1}{-1},
  \qquad
  C=
  \begin{pmatrix}
  A&B\\B&D
  \end{pmatrix},
  \qquad
  C_Q=A+2B+D,
  \qquad
  p_Q=p_3+p_Y,
\]
and
\[
  S_p=p_3(D+B)-p_Y(A+B).
\]
Then
\[
  \det(C+vp^T+pv^T+\Omega_Qvv^T)
  =
  \Delta_C+2aS_p
  +a^2(\Omega_QC_Q-p_Q^2).
\]
Since \(L_Q^2=C_Q\), branch closure becomes
\[
  \left((1-R_{\rm dV})\Omega_QC_Q-p_Q^2\right)a^2
  +2\left(S_p-R_{\rm dV}p_3C_Q\right)a
  -D_C=0 .
\]
Equivalently, with \(u=\Omega_Qa^2\), \(\epsilon={\rm sign}(a)\), and
\[
  A_N=(1-R_{\rm dV})C_Q-\frac{p_Q^2}{\Omega_Q},
  \qquad
  B_N=\frac{S_p-R_{\rm dV}p_3C_Q}{\sqrt{\Omega_Q}},
\]
the positive square-root variable \(y=\sqrt u\) obeys
\[
  A_Ny^2+2\epsilon B_Ny-D_C=0 .
\]
Thus the harmonic square is
\[
  u_{\epsilon}
  =
  \left(
  \frac{-\epsilon B_N+\sqrt{B_N^2+A_ND_C}}{A_N}
  \right)^2,
  \qquad A_N\ne0,
\]
with affine limit \(u_\epsilon=D_C^2/(4B_N^2)\) when
\(A_N=0\) and \(\epsilon B_N>0\).  For \(p_A=0\), the row gives
\[
  u_N=\frac{D_C}{(1-R_{\rm dV})C_Q}.
\]
The cross vector has a direct DBI readout.  Let \(s_Q\) be arclength on
the \(Q\)-circle, \(0\le s_Q<L_Q\), and take
\[
  \theta_{\Gamma,{\rm CP}}^A=\theta_i^A(y,\sigma)dy^i,
  \qquad
  \omega_Q=\frac{ds_Q}{L_Q}.
\]
With sector density
\[
  d\nu_{\Gamma,\sigma}
  =
  \kappa_8\omega_\sigma e^{-\phi_\sigma}
  \sqrt{\gamma_{\Gamma,\sigma}}\,
  d^4y\,ds_Q ,
\]
one obtains
\[
  p_A
  =
  \frac{1}{L_Q}
  \sum_{\sigma\in{\cal S}_\Gamma}
  \int_{\Gamma_{{\rm int},\sigma}}
  d\nu_{\Gamma,\sigma}\,
  \gamma_\sigma^{is_Q}\,
  {\rm tr}_{E_\sigma}
  \left(\Pi_Q\theta_i^A\Pi_Q\right),
\]
and
\[
  \Omega_Q
  =
  \frac{1}{L_Q^2}
  \sum_{\sigma\in{\cal S}_\Gamma}
  \int_{\Gamma_{{\rm int},\sigma}}
  d\nu_{\Gamma,\sigma}\,
  \gamma_\sigma^{s_Qs_Q}\,
  {\rm tr}_{E_\sigma}\Pi_Q .
\]
For a block-product internal metric,
\[
  \gamma_\sigma^{is_Q}=0,
  \qquad
  \gamma_\sigma^{s_Qs_Q}=1,
\]
hence
\[
  p_3=p_Y=0,
  \qquad
  \Omega_Q=\frac{Z_{\Gamma,Q}}{L_Q}.
\]
A nonzero \(p_A\) records vertical leakage of the retained CP coframe or an
off-diagonal frame contribution after choosing the adapted \(Q\)-circle
one-form.
Equivalently, for
\[
  \chi_Q=ds_Q+{\cal A}_i(y)dy^i,
  \qquad
  \omega_Q=\frac{\chi_Q}{L_Q},
\]
decompose
\[
  \theta_{\Gamma,{\rm CP}}^A
  =
  \theta_H^A+\ell_A\chi_Q,
  \qquad
  \gamma_\Gamma^{-1}(\theta_H^A,\chi_Q)=0 .
\]
Then
\[
  p_A
  =
  \frac{1}{L_Q}
  \sum_{\sigma\in{\cal S}_\Gamma}
  \int_{\Gamma_{{\rm int},\sigma}}
  d\nu_{\Gamma,\sigma}\,
  \ell_A\,
  {\rm tr}_{E_\sigma}\Pi_Q .
\]
Thus the horizontal adapted row has \(p_A=0\), and the leakage functions
\(\ell_3,\ell_Y\) are the localized sources for the cross-vector branch.
Pointwise,
\[
  \ell_A
  =
  \frac{\gamma_\Gamma^{-1}(\theta_{\Gamma,{\rm CP}}^A,\chi_Q)}
       {\gamma_\Gamma^{-1}(\chi_Q,\chi_Q)} .
\]
For the product representative
\[
  \pi_{\CP^2}:\CP^2\times S^1_Q\longrightarrow\CP^2,
  \qquad
  \theta_{\Gamma,{\rm CP}}^A=\pi_{\CP^2}^*\alpha_A,
\]
and adapted metric
\[
  ds_\Gamma^2=g_{ij}(y)dy^idy^j+L_Q^2\chi_Q^2,
\]
one obtains
\[
  \gamma_\Gamma^{-1}(\pi_{\CP^2}^*\alpha_A,\chi_Q)=0,
\]
hence
\[
  \ell_3=\ell_Y=0,
  \qquad
  p_3=p_Y=0,
  \qquad
  \Omega_Q=\frac{Z_{\Gamma,Q}}{L_Q}.
\]
A cross-vector branch with \(p_A\ne0\) is a separate leakage realization in
which the retained CP coframe has a \(\chi_Q\)-component after adapted
projection.
For the product \(\CP^2\times S^1_Q\) row, take
\[
  \omega_Q=\frac{ds_Q}{L_Q},
  \qquad
  \int_{S^1_Q}\omega_Q=1 .
\]
The CP2 zero-mode DBI density is
\[
  Z_{\Gamma,Q}
  =
  \sum_s\int_{\CP^2_s}
  \kappa_8\,\omega_s e^{-\phi_s}
  \sqrt{\gamma_{\CP^2,s}}\,
  {\rm tr}_{E_s}\Pi_Q .
\]
Hence
\[
  \Omega_Q
  =
  \langle\omega_Q,\omega_Q\rangle_{\Gamma,\Pi_Q}
  =
  \frac{Z_{\Gamma,Q}}{L_Q},
\]
and the harmonic weak vertical length is
\[
  L_{2,{\rm vert}}^2
  =
  (2\pi)^2\frac{Z_{\Gamma,Q}}{L_Q}\,
  q_3^{-2}\vartheta_3^2 .
\]
Define the finite trace ratio
\[
  \tau_Q=\frac{Z_{\Gamma,Q}}{Z_{\Gamma,0}} .
\]
Then
\[
  L_{2,{\rm vert}}^2
  =
  (2\pi)^2
  \frac{\tau_QZ_{\Gamma,0}}{L_Q}\,
  q_3^{-2}\vartheta_3^2 .
\]
For a shared CP2 density profile and normalized endpoint vector \(\eta_0\),
\[
  \tau_Q={\rm tr}_{E_Q}\Pi_Q .
\]
At \(Q=7\), the full retained \(E_Q\) trace gives
\[
  \tau_Q=16 .
\]
Using the integer residual convention
\[
  q_3=3Q,
\]
and the full retained trace \(\tau_Q=2(Q+1)\), this becomes
\[
  L_{2,{\rm vert}}^2
  =
  \frac{8\pi^2(Q+1)}{9Q^2}
  \frac{Z_{\Gamma,0}}{L_Q}\,\delta\vartheta_3^2 ,
\]
where \(\delta\vartheta_3\) is the shortest representative of
\(\vartheta_3\in\mathbb R/\mathbb Z\).  For \(Q=7\),
\[
  L_{2,{\rm vert}}^2
  =
  \frac{64\pi^2}{441}
  \frac{Z_{\Gamma,0}}{L_Q}\,\delta\vartheta_3^2 .
\]
Equivalently,
\[
  \delta\vartheta_3
  =
  \frac{21}{8\pi}
  \sqrt{\frac{L_QL_{2,{\rm vert}}^2}{Z_{\Gamma,0}}}.
\]
The Wilson chamber \(|\delta\vartheta_3|\le\half\) gives
\[
  L_{2,{\rm vert}}^2
  \le
  \frac{16\pi^2}{441}\frac{Z_{\Gamma,0}}{L_Q}.
\]
For a nonzero harmonic representative, define
\[
  \Xi_Q^{\rm harm}
  =
  \frac{Z_{\Gamma,0}}{L_QL_{2,{\rm vert}}^2}.
\]
Then the full-trace row gives
\[
  \Xi_Q^{\rm harm}
  =
  \frac{9Q^2}{8\pi^2(Q+1)}
  \frac{1}{\delta\vartheta_3^2},
\]
and, at \(Q=7\),
\[
  \Xi_7^{\rm harm}
  =
  \frac{441}{64\pi^2}
  \frac{1}{\delta\vartheta_3^2}.
\]
The fundamental Wilson chamber gives
\[
  \Xi_7^{\rm harm}
  \ge
  \frac{441}{16\pi^2}
  =
  2.792665124141\ldots .
\]
The D9 global-form descent fixes \(q_3=3Q\).  For the retained \(Q=7\)
stack, \(q_3=21\), so
\[
  \Theta_N(21,q_Y)
  =
  \half\min\left(1,\left|\frac{21}{q_Y}\right|\right).
\]
Equivalently,
\[
  \Theta_N(21,q_Y)=\half
  \qquad (|q_Y|\le21),
\]
and
\[
  \Theta_N(21,q_Y)=\frac{21}{2|q_Y|}
  \qquad (|q_Y|>21).
\]
For the endpoint character in the same retained \(Q=7\) row,
\[
  c_1(L_Y)=6h+Q\eta,
  \qquad
  c_1(L_{T_3})=\eta,
  \qquad
  c_1(L_Y)-Q\,c_1(L_{T_3})=6h .
\]
The quotient line operator
\[
  {\cal O}_Y(\gamma_Q)
  =
  {\rm Hol}_{L_Y}(\gamma_Q)\,
  {\rm Hol}_{L_{T_3}}(\gamma_Q)^{-Q}
\]
has Chern class \(6h\).  Since \(h\) pulls back from \(\CP^2\), this operator
has unit \(S^1_Q\) character.  The hypercharge boundary holonomy on
\(\gamma_Q\) is therefore the \(Q\)-fold weak endpoint holonomy:
\[
  {\rm Hol}_{L_Y}(\gamma_Q)
  =
  {\rm Hol}_{L_{T_3}}(\gamma_Q)^Q .
\]
For the retained endpoint vector \(\eta_Q=\eta_0^{\odot Q}\),
\[
  T_3\eta_Q=-\frac Q2\eta_Q,
  \qquad
  Y\eta_Q=\frac Q2\eta_Q .
\]
With integer hypercharge \(y=6Y\), the hypercharge Wilson period is
\[
  q_Y=3Q=21 .
\]
Consequently
\[
  \Theta_N(21,21)=\half,
  \qquad
  \delta\vartheta_Y=-\delta\vartheta_3 .
\]
For the photon-preserving massive-neutral branch,
\[
  v_Q=0,
  \qquad
  \delta\vartheta_Y=-\frac{q_Y}{q_3}\delta\vartheta_3 .
\]
Inside the identity chamber, both Wilson representatives obey
\[
  |\delta\vartheta_3|\le\Theta_N(q_3,q_Y),
  \qquad
  \Theta_N(q_3,q_Y)
  =
  \half\min\left(1,\left|\frac{q_3}{q_Y}\right|\right).
\]
Thus
\[
  \Xi_Q^{\rm harm}
  \ge
  \frac{9Q^2}{8\pi^2(Q+1)\Theta_N(q_3,q_Y)^2},
\]
with the \(Q=7\) specialization
\[
  \Xi_7^{\rm harm}
  \ge
  \frac{441}{64\pi^2\Theta_N(q_3,q_Y)^2}.
\]
For the massive-neutral amplitude
\[
  u_N=\frac{D_C}{(1-R_{\rm dV})C_Q},
\]
the density condition becomes
\[
  \frac{Z_{\Gamma,0}}{L_Q}
  \ge
  \frac{9Q^2}{8\pi^2(Q+1)\Theta_N(q_3,q_Y)^2}
  \frac{D_C}{(1-R_{\rm dV})C_Q}.
\]
At \(Q=7\),
\[
  \frac{Z_{\Gamma,0}}{L_Q}
  \ge
  \frac{441}{64\pi^2\Theta_N(q_3,q_Y)^2}
  \frac{D_C}{(1-R_{\rm dV})C_Q}.
\]
With \(q_3=21\), the same condition is
\[
  \frac{Z_{\Gamma,0}}{L_Q}
  \ge
  \frac{441}{16\pi^2}
  \frac{D_C}{(1-R_{\rm dV})C_Q}
  \qquad (|q_Y|\le21),
\]
and
\[
  \frac{Z_{\Gamma,0}}{L_Q}
  \ge
  \frac{q_Y^2}{16\pi^2}
  \frac{D_C}{(1-R_{\rm dV})C_Q}
  \qquad (|q_Y|>21).
\]
For the endpoint-character value \(q_Y=21\), the density condition reduces to
\[
  \frac{Z_{\Gamma,0}}{L_Q}
  \ge
  \frac{441}{16\pi^2}
  \frac{D_C}{(1-R_{\rm dV})C_Q}.
\]
Equivalently,
\[
  U_\Gamma^{(21)}
  =
  \frac{16\pi^2}{441}\frac{Z_{\Gamma,0}}{L_Q}
\]
obeys
\[
  U_\Gamma^{(21)}
  \ge
  \frac{D_C}{(1-R_{\rm dV})C_Q}.
\]
The corresponding shortest Wilson representatives satisfy
\[
  \delta\vartheta_3^2
  =
  \delta\vartheta_Y^2
  =
  \frac{1}{4U_\Gamma^{(21)}}
  \frac{D_C}{(1-R_{\rm dV})C_Q}.
\]
The product DBI density factors the endpoint norm.  Write
\[
  d\nu_{\Gamma,s}
  =
  d\mu_{\Gamma,s}^{\CP^2}\,ds_Q,
  \qquad
  d\mu_{\Gamma,s}^{\CP^2}
  =
  \kappa_8\omega_se^{-\phi_s}
  \sqrt{\gamma_{\CP^2,s}}\,d^4y .
\]
For an \(S^1_Q\)-invariant neutral endpoint section,
\[
  Z_{\Gamma,0}
  =
  L_Q\,{\cal Z}_{\Gamma,0}^{\CP^2},
\]
where
\[
  {\cal Z}_{\Gamma,0}^{\CP^2}
  =
  \sum_s
  \int_{\CP^2_s}d\mu_{\Gamma,s}^{\CP^2}\,
  {\rm tr}_{E_s}\left(\Pi_Q|\eta_Q|^2\Pi_Q\right).
\]
Hence
\[
  U_\Gamma^{(21)}
  =
  \frac{16\pi^2}{441}{\cal Z}_{\Gamma,0}^{\CP^2}.
\]
The localized density condition becomes
\[
  {\cal Z}_{\Gamma,0}^{\CP^2}
  \ge
  \frac{441}{16\pi^2}
  \frac{D_C}{(1-R_{\rm dV})C_Q}.
\]
Let \({\mathfrak z}_{\Gamma}^{\CP^2}\) denote the common CP2 density per
finite endpoint state,
\[
  {\mathfrak z}_{\Gamma}^{\CP^2}
  =
  \sum_s
  \int_{\CP^2_s}d\mu_{\Gamma,s}^{\CP^2}.
\]
The unit rank-one endpoint normalization gives
\[
  {\rm tr}_{E_Q}
  \left(\Pi_Q|\eta_Q\rangle\langle\eta_Q|\Pi_Q\right)=1,
\]
and the full retained trace gives
\[
  {\rm tr}_{E_Q}\Pi_Q=2(Q+1)=16.
\]
Thus
\[
  {\cal Z}_{\Gamma,0}^{\CP^2}
  =
  {\mathfrak z}_{\Gamma}^{\CP^2},
  \qquad
  {\cal Z}_{\Gamma,Q}^{\CP^2}
  =
  16\,{\mathfrak z}_{\Gamma}^{\CP^2},
  \qquad
  \tau_Q=16 .
\]
The D8/O8 density datum is the scalar \({\mathfrak z}_{\Gamma}^{\CP^2}\).
In a constant product profile, \(\omega_s\), \(\phi_s\), and
\(\gamma_{\CP^2,s}\) are constant on each retained sector, giving
\[
  {\mathfrak z}_{\Gamma}^{\CP^2}
  =
  \kappa_8
  \sum_s\omega_s e^{-\phi_s}
  {\rm Vol}(\CP^2_s,\gamma_{\CP^2,s}) .
\]
The density-dominated pass is therefore the effective volume inequality
\[
  \kappa_8
  \sum_s\omega_s e^{-\phi_s}
  {\rm Vol}(\CP^2_s,\gamma_{\CP^2,s})
  \ge
  9.724809426610\ldots .
\]
For a common CP2 metric and dilaton, this becomes
\[
  \kappa_8 e^{-\phi_0}\Omega_\Gamma
  {\rm Vol}(\CP^2,\gamma_{\CP^2})
  \ge
  9.724809426610\ldots,
  \qquad
  \Omega_\Gamma=\sum_s\omega_s .
\]
Use the same WZ normalization as the CS source row,
\[
  \int_{\CP^2}h_Q^2=1.
\]
Writing
\[
  J_{\Gamma,s}=r_{\Gamma,s}h_Q
\]
gives
\[
  {\rm Vol}(\CP^2_s,\gamma_{\CP^2,s})
  =
  \frac12\int_{\CP^2_s}J_{\Gamma,s}^2
  =
  \frac12 r_{\Gamma,s}^2 .
\]
Thus the density condition is
\[
  \frac{\kappa_8}{2}
  \sum_s\omega_s e^{-\phi_s}r_{\Gamma,s}^2
  \ge
  \frac{441}{16\pi^2}
  \frac{D_C}{(1-R_{\rm dV})C_Q}.
\]
The density-dominated metric pass is
\[
  \sum_s\omega_s e^{-\phi_s}r_{\Gamma,s}^2
  \ge
  \frac{19.449618853220\ldots}{\kappa_8}.
\]
For a common metric and dilaton,
\[
  r_\Gamma^2
  \ge
  \frac{19.449618853220\ldots}
       {\kappa_8e^{-\phi_0}\Omega_\Gamma}.
\]
The density scalar fixes \(U_\Gamma^{\CP^2}\).  The point-row shape parameter
comes from a separate normalized Hodge readout.  Split the retained weak
coframe into a unit endpoint component and a horizontal excess,
\[
  \theta_\Gamma^3=\theta_{3,0}+\theta_{3,h},
  \qquad
  \langle\theta_{3,0},\theta_{3,h}\rangle_\Gamma=0 .
\]
Define
\[
  C_0^\Gamma
  =
  \sum_s\int_{\CP^2_s}d\mu_{\Gamma,s}^{\CP^2}\,
  {\rm tr}_{E_s}
  \left(\Pi_Q|\theta_{3,0}|^2\Pi_Q\right),
\]
and
\[
  C_h^\Gamma
  =
  \sum_s\int_{\CP^2_s}d\mu_{\Gamma,s}^{\CP^2}\,
  {\rm tr}_{E_s}
  \left(\Pi_Q|\theta_{3,h}|^2\Pi_Q\right).
\]
The point-evaluation row reads
\[
  \widehat C_{33}^\Gamma
  =
  \frac{C_0^\Gamma+C_h^\Gamma}{C_0^\Gamma}
  =
  1+\frac{Q^2+1}{(Q+1)^2}q_\Gamma .
\]
Hence the localized shape readout is
\[
  q_\Gamma
  =
  \frac{(Q+1)^2}{Q^2+1}
  \frac{C_h^\Gamma}{C_0^\Gamma}.
\]
For \(Q=7\),
\[
  q_\Gamma=\frac{32}{25}\frac{C_h^\Gamma}{C_0^\Gamma}.
\]
Define the dimensionless shape ratio
\[
  \varrho_\Gamma=\frac{C_h^\Gamma}{C_0^\Gamma}.
\]
Then
\[
  q_\Gamma=\frac{32}{25}\varrho_\Gamma,
  \qquad
  0<\varrho_\Gamma<\varrho_c,
  \qquad
  \varrho_c=\frac{25}{32}q_c=2.743448440896\ldots .
\]
The point-row amplitude as a function of the localized shape ratio is
\[
  u_\Gamma^{\rm sh}(\varrho_\Gamma)
  =
  \frac{
  625R_{\rm dV}(1+\varrho_\Gamma)^2
  -576\varrho_\Gamma^2-900\varrho_\Gamma}
  {625(1-R_{\rm dV})(1+\varrho_\Gamma)} .
\]
The charge chamber becomes the pair of localized inequalities
\[
  U_\Gamma^{\CP^2}\ge u_\Gamma^{\rm sh}(\varrho_\Gamma),
  \qquad
  0<\varrho_\Gamma<\varrho_c .
\]
Equivalently, the localized density satisfies
\[
  {\cal Z}_{\Gamma,0}^{\CP^2}
  \ge
  \frac{441}{16\pi^2}
  u_\Gamma^{\rm sh}(\varrho_\Gamma).
\]
In the constant-profile metric row this becomes
\[
  \sum_s\omega_s e^{-\phi_s}r_{\Gamma,s}^2
  \ge
  \frac{441}{8\pi^2\kappa_8}
  u_\Gamma^{\rm sh}(\varrho_\Gamma).
\]
For a common metric and dilaton,
\[
  r_\Gamma^2
  \ge
  \frac{441}{8\pi^2\kappa_8e^{-\phi_0}\Omega_\Gamma}
  u_\Gamma^{\rm sh}(\varrho_\Gamma).
\]
The point-evaluation row gives a concrete base matrix.  Put
\[
  q=\lambda^2,
  \qquad
  t_Q=-\frac{Q}{Q+1},
\]
and
\[
  C_{33}^{\rm pt}
  =
  1+\frac{Q^2+1}{(Q+1)^2}q .
\]
Then
\[
  C_Q^{\rm pt}=(Q+1)^2C_{33}^{\rm pt},
  \qquad
  \Delta_C^{\rm pt}=(Q-1)^2q(q+2),
\]
so
\[
  u_N^{\rm pt}(q)
  =
  \frac{
  R_{\rm dV}(Q+1)^2(C_{33}^{\rm pt})^2
  -(Q-1)^2q(q+2)}
  {(1-R_{\rm dV})(Q+1)^2C_{33}^{\rm pt}}.
\]
For \(Q=7\),
\[
  u_N^{\rm pt}(q)
  =
  \frac{R_{\rm dV}(25q+32)^2-576q(q+2)}
       {32(1-R_{\rm dV})(25q+32)}.
\]
The zero of this amplitude is
\[
  q_c=3.511614004347\ldots .
\]
The derivative is
\[
  \frac{d u_N^{\rm pt}}{dq}
  =
  -\frac{
  (14400-15625R_{\rm dV})q^2
  +(36864-40000R_{\rm dV})q
  +(36864-25600R_{\rm dV})}
  {32(1-R_{\rm dV})(25q+32)^2}.
\]
Since \(0<R_{\rm dV}<1\), the amplitude is strictly decreasing for
\(q>0\), with
\[
  u_N^{\rm pt}(0)=\frac{R_{\rm dV}}{1-R_{\rm dV}}
  =
  3.482268368857\ldots,
  \qquad
  u_N^{\rm pt}(q_c)=0.
\]
Together with the product DBI readout, this gives a density-dominated pass:
\[
  {\cal Z}_{\Gamma,0}^{\CP^2}
  \ge
  \frac{441}{16\pi^2}
  \frac{R_{\rm dV}}{1-R_{\rm dV}}
  =
  9.724809426610\ldots .
\]
In that regime the full interval \(0<q<q_c\) satisfies the
endpoint-character chamber condition.  For lower positive density, set
\[
  U_\Gamma^{\CP^2}
  =
  \frac{16\pi^2}{441}{\cal Z}_{\Gamma,0}^{\CP^2}.
\]
Then the active interval is
\[
  q_{\rm sat}(U_\Gamma^{\CP^2})\le q<q_c .
\]
For \(0<q<q_c\), the photon-preserving harmonic lift has positive
amplitude, and the \(Q=7\) density condition becomes
\[
  \frac{Z_{\Gamma,0}}{L_Q}
  \ge
  \frac{
  441\left(R_{\rm dV}(25q+32)^2-576q(q+2)\right)}
  {2048\pi^2\Theta_N(q_3,q_Y)^2(1-R_{\rm dV})(25q+32)} .
\]
Define the normalized density
\[
  U_\Gamma
  =
  \frac{64\pi^2\Theta_N(q_3,q_Y)^2}{441}
  \frac{Z_{\Gamma,0}}{L_Q}.
\]
The chamber condition is \(U_\Gamma\ge u_N^{\rm pt}(q)\).  On the saturation
curve \(U_\Gamma=u_N^{\rm pt}(q)\), \(q\) solves
\[
  A_q q^2+B_q(U_\Gamma)q+C_q(U_\Gamma)=0,
\]
where
\[
  A_q=576-625R_{\rm dV},
\]
\[
  B_q(U_\Gamma)
  =
  1152-1600R_{\rm dV}+800(1-R_{\rm dV})U_\Gamma,
\]
and
\[
  C_q(U_\Gamma)
  =
  1024\left((1-R_{\rm dV})U_\Gamma-R_{\rm dV}\right).
\]
For
\[
  0\le U_\Gamma\le \frac{R_{\rm dV}}{1-R_{\rm dV}},
\]
the positive branch is
\[
  q_{\rm sat}(U_\Gamma)
  =
  \frac{
  -B_q(U_\Gamma)
  +\sqrt{B_q(U_\Gamma)^2-4A_qC_q(U_\Gamma)}}
  {2A_q}.
\]
The chamber pass is \(q_{\rm sat}(U_\Gamma)\le q<q_c\).  Given \(q\) and
\(U_\Gamma\), the Wilson representative is
\[
  \delta\vartheta_3^2
  =
  \Theta_N(q_3,q_Y)^2
  \frac{u_N^{\rm pt}(q)}{U_\Gamma}.
\]
A quadratic local expansion has the form
\[
\begin{aligned}
S_\Gamma^{(2)}
=\int_{M_4}\Bigg[
&-\frac14{\cal K}_{AB}^\Gamma F_{\mu\nu}^AF^{B\mu\nu}
+\frac12(f_H^2)_{ab}
  \partial_\mu\theta_H^a\partial^\mu\theta_H^b\\
&+\frac12\langle r,{\cal M}_{\rm img}r\rangle_\Gamma
+L_2\,{\mathfrak j}_A
  \langle r,{\cal J}_A\rangle_\Gamma
+\frac12{\rm Str}_\chi\log{\cal B}_\Gamma^\chi(\theta_H)
\Bigg].
\end{aligned}
\]
The current matrix is
\[
  {\cal K}_{AB}^\Gamma
  =
  \int_{\Gamma_{\rm int}}d\nu_\Gamma\,
  \gamma_\Gamma^{-1}(\theta_\Gamma^A,\theta_\Gamma^B),
  \qquad A,B=3,Y,Q .
\]
The image source functions are
\[
  {\cal J}_A
  =
  \gamma_\Gamma^{-1}(\theta_\Gamma^A,\Theta_{\rm img}),
  \qquad A=3,Y .
\]
Eliminating \(r\) gives
\[
  \Delta S_{\rm eff}
  =
  -\frac12L_2^2{\mathfrak j}_A
  {\cal D}_{AB}{\mathfrak j}_B,
  \qquad
  {\cal D}_{AB}
  =
  L_2^{-2}
  \langle{\cal J}_A,
  {\cal M}_{\rm img}^{-1}{\cal J}_B\rangle_\Gamma .
\]
Thus the same local expression yields
\[
  L_2^2={\cal K}_{33}^\Gamma,
  \qquad
  L_Q^2=(1,1){\cal K}_{(3,Y)}^\Gamma\binom{1}{1},
\]
\[
  \Delta_{\rm rank}
  =
  {\cal D}_{33}{\cal D}_{YY}-{\cal D}_{3Y}^2,
  \qquad
  \rho_D=\frac{{\cal D}_{3Y}}{{\cal D}_{33}},
\]
and the Wilson kinetic entry
\[
  (f_H^2)_{ab}
  =
  \int_{\Gamma_{\rm int}}d\nu_\Gamma\,
  g^{ss}\omega_{H,s}\omega_{H,s}\,
  {\rm tr}_{\rm end}(n_H^cT_c\,n_H^dT_d)\Pi_{ab}^{\rm Schur}.
\]
For the point-projector endpoint metric this row gives
\[
  Z_{\Gamma,H}=\frac Q4Z_{\Gamma,0}.
\]
The charged determinant readout is
\[
  {\cal N}_\Gamma^\chi
  =
  {\rm Str}_{\ker{\cal B}_\Gamma^\chi}
  (\tau_\alpha)
  =
  2\sigma_\chi d_\chi F(z)
\]
on the equal-threshold Borel--Weil charged pair.  The single local coefficient
test is therefore
\[
  \left(
  {\cal D}_{AB},\,
  \frac{Z_{\Gamma,H}}{Z_{\Gamma,0}},\,
  {\cal N}_\Gamma^\chi
  \right)
  =
  \left(
  {\cal D}_\sigma,\,
  \frac Q4,\,
  2\sigma_\chi d_\chi F(z)
  \right).
\]

\subsection{Finite-Matrix Extraction}

Choose a \(\Gamma\)-orthonormal basis
\(\{e_i\}_{i=1}^N\) for the sourced image subspace and define
\[
  M_{ij}=\langle e_i,{\cal M}_{\rm img}e_j\rangle_\Gamma,
  \qquad
  C_{Ai}=\langle e_i,{\cal J}_A\rangle_\Gamma ,
  \qquad A=3,Y .
\]
The image correction is
\[
  {\cal D}_{AB}
  =
  L_2^{-2}\,C_{Ai}(M^{-1})^{ij}C_{Bj}.
\]
With
\[
  w_A^i=L_2^{-1}(M^{-1/2})^{ij}C_{Aj},
\]
one has
\[
  {\cal D}_{AB}=w_A\cdot w_B .
\]
The Schur branch condition is the finite-vector equation
\[
  w_Y=\rho_\sigma w_3,
  \qquad
  \|w_3\|^2=\delta_\sigma .
\]
Thus any localized image spectrum produces the two Schur numbers by a direct
matrix inversion.  Sourced modes away from this line contribute to
\(\Delta_{\rm rank}\); unsourced modes drop out of \({\cal D}_{AB}\).

For the charged endpoint determinant, let
\[
  ({\cal B}_{r,\Gamma}^\chi)^\dagger
  {\cal B}_{r,\Gamma}^\chi u_{\alpha r}
  =
  M_{\alpha r}^2u_{\alpha r},
  \qquad r=1,2 .
\]
The equal-amplitude Borel--Weil row is the spectral equality
\[
  \{(\sigma_{\alpha1},d_{\alpha1},M_{\alpha1})\}_\alpha
  =
  \{(\sigma_{\alpha2},d_{\alpha2},M_{\alpha2})\}_\alpha .
\]
It gives
\[
  {\cal N}_\Gamma^\chi
  =
  2\sum_\alpha
  \sigma_\alpha d_\alpha
  F(M_\alpha L_H).
\]
The minimal endpoint row is
\[
  M_\alpha=M,\qquad
  {\cal N}_\Gamma^\chi=2\sigma_\chi d_\chi F(ML_H).
\]

\subsection{Minimal Finite Spectrum Package}

The finite extraction has a minimal image realization on
\[
  {\cal H}_{{\rm img},\sigma}^{\rm src}
  =
  {\rm span}\{e_\sigma\},
  \qquad
  \langle e_\sigma,e_\sigma\rangle_\Gamma=1,
  \qquad
  {\cal M}_{\rm img}e_\sigma=m_\sigma e_\sigma .
\]
The source couplings are
\[
  C_{3\sigma}
  =
  L_2\sqrt{m_\sigma}\sqrt{\delta_\sigma},
  \qquad
  C_{Y\sigma}
  =
  L_2\sqrt{m_\sigma}\sqrt{\delta_\sigma}\rho_\sigma .
\]
Equivalently,
\[
  \frac{1}{L_2\sqrt{m_\sigma}}
  \binom{C_{3\sigma}}{C_{Y\sigma}}
  =
  g_\sigma .
\]
For the two branches,
\[
\begin{array}{c|c|c}
{\rm branch}&(g_\sigma)_3&(g_\sigma)_Y\\ \hline
+&0.374811401563\ldots&7.371454070999\ldots\\
-&0.572116571654\ldots&7.230005311453\ldots .
\end{array}
\]
Substitution gives
\[
  {\cal D}_{AB}
  =
  \frac{C_{A\sigma}C_{B\sigma}}{m_\sigma L_2^2}
  =
  (g_\sigma)_A(g_\sigma)_B
  =
  ({\cal D}_\sigma)_{AB}.
\]
In the constant-profile product row, the same density gives the normalized
source mode
\[
  \psi_{\Gamma,0}
  =
  \left(L_Q{\mathfrak z}_{\Gamma}^{\CP^2}\right)^{-1/2}\eta_Q,
  \qquad
  \int_{\Gamma_{\rm EW,int}}d\nu_\Gamma\,
  {\rm tr}_{E_Q}\left(\Pi_Q|\psi_{\Gamma,0}|^2\Pi_Q\right)=1 .
\]
Taking \(e_\sigma=\psi_{\Gamma,0}\), the branch source current is
\[
  {\cal J}_A^{(\sigma)}
  =
  L_2\sqrt{m_\sigma}\,(g_\sigma)_A\,e_\sigma^\flat,
  \qquad A=3,Y.
\]
For the current Gram matrix
\[
  G_{AB}^{(0)}
  =
  \int_{\Gamma_{\rm EW,int}}d\nu_\Gamma\,
  \psi_{\Gamma,0}^2
  \gamma_\Gamma^{-1}(\theta_\Gamma^A,\theta_\Gamma^B),
\]
the Gram-dual image one-form is
\[
  \Xi_\sigma
  =
  L_2\psi_{\Gamma,0}\,
  (G^{(0)})^{-1\,BC}(g_\sigma)_C\theta_\Gamma^B .
\]
It satisfies
\[
  L_2^{-1}
  \int d\nu_\Gamma\,\psi_{\Gamma,0}\,
  \gamma_\Gamma^{-1}(\theta_\Gamma^A,\Xi_\sigma)
  =
  (g_\sigma)_A,
  \qquad
  R_{\Xi,\sigma}=0 .
\]
The constant-profile row therefore forms a single-pass data packet
\[
  {\cal P}_{\Gamma}^{\rm cp}
  =
  \left(
  \kappa_8,\,
  e^{-\phi_0}\Omega_\Gamma,\,
  r_\Gamma,\,
  L_Q,\,
  \varrho_\Gamma,\,
  m_\sigma,\,
  z_\chi,\,
  \sigma_\chi d_\chi
  \right).
\]
It returns
\[
  e^{-\phi_0}\Omega_\Gamma
  =
  \sum_{s\in{\cal S}_\Gamma}\omega_s e^{-\phi_s},
  \qquad
  \Omega_\Gamma
  =
  \sum_{s\in{\cal S}_\Gamma}\omega_s .
\]
The Chan--Paton trace factor remains the finite endpoint multiplier
\(\tau_Q=16\) in \(Z_{\Gamma,Q}\).  In the common-metric sector row,
\[
  e^{-\phi_0}
  =
  \frac{\sum_{s\in{\cal S}_\Gamma}\omega_s e^{-\phi_s}}
  {\sum_{s\in{\cal S}_\Gamma}\omega_s}.
\]
The measure and CP2 density are
\[
  d\nu_\Gamma
  =
  \kappa_8e^{-\phi_0}\Omega_\Gamma
  \sqrt{\gamma_{\CP^2}}\,d^4y\,ds_Q,
  \qquad
  {\mathfrak z}_{\Gamma}^{\CP^2}
  =
  \frac{\kappa_8}{2}e^{-\phi_0}\Omega_\Gamma r_\Gamma^2 ,
\]
\[
  Z_{\Gamma,0}=L_Q{\mathfrak z}_{\Gamma}^{\CP^2},
  \qquad
  Z_{\Gamma,Q}=16L_Q{\mathfrak z}_{\Gamma}^{\CP^2},
  \qquad
  Z_{\Gamma,H}=\frac Q4Z_{\Gamma,0},
\]
and
\[
  q_\Gamma=\frac{32}{25}\varrho_\Gamma,
  \qquad
  U_\Gamma^{\CP^2}
  =
  \frac{16\pi^2}{441}{\mathfrak z}_{\Gamma}^{\CP^2}.
\]
The shape parameter is the current-frame Rayleigh quotient
\[
  \varrho_\Gamma
  =
  \frac{C_h^\Gamma}{C_0^\Gamma},
  \qquad
  C_0^\Gamma
  =
  \langle P_0^\Gamma\theta_\Gamma^3,
  P_0^\Gamma\theta_\Gamma^3\rangle_{\Gamma,\Pi_Q},
\]
\[
  C_h^\Gamma
  =
  \langle P_h^\Gamma\theta_\Gamma^3,
  P_h^\Gamma\theta_\Gamma^3\rangle_{\Gamma,\Pi_Q},
  \qquad
  C_0^\Gamma>0 .
\]
Here \(P_0^\Gamma\) projects to the endpoint component and \(P_h^\Gamma\)
projects to the horizontal \(\CP^2\) component in the localized current
frame.
In sector form,
\[
  C_{0,s}^\Gamma
  =
  \omega_s e^{-\phi_s}
  \langle P_0^\Gamma\theta_{\Gamma,s}^3,
  P_0^\Gamma\theta_{\Gamma,s}^3\rangle_{s,\Pi_Q},
  \qquad
  C_{h,s}^\Gamma
  =
  \omega_s e^{-\phi_s}
  \langle P_h^\Gamma\theta_{\Gamma,s}^3,
  P_h^\Gamma\theta_{\Gamma,s}^3\rangle_{s,\Pi_Q}.
\]
Therefore
\[
  C_0^\Gamma=\sum_sC_{0,s}^\Gamma,
  \qquad
  C_h^\Gamma=\sum_sC_{h,s}^\Gamma,
  \qquad
  \varrho_\Gamma
  =
  \frac{\sum_sC_{h,s}^\Gamma}{\sum_sC_{0,s}^\Gamma}.
  \tag{5.1}
\]
On \({\cal S}_{\Gamma,0}=\{s:C_{0,s}^\Gamma>0\}\), set
\[
  b_{\Gamma,s}=\frac{C_{0,s}^\Gamma}{C_0^\Gamma},
  \qquad
  \rho_{\Gamma,s}=\frac{C_{h,s}^\Gamma}{C_{0,s}^\Gamma}.
\]
Then
\[
  b_{\Gamma,s}\ge0,\qquad
  \sum_{s\in{\cal S}_{\Gamma,0}}b_{\Gamma,s}=1,
  \qquad
  \varrho_\Gamma
  =
  \sum_{s\in{\cal S}_{\Gamma,0}}b_{\Gamma,s}\rho_{\Gamma,s}.
  \tag{5.2}
\]
The density-shape pass is
\[
  0<\varrho_\Gamma<2.743448440896\ldots,
  \qquad
  U_\Gamma^{\CP^2}\ge u_\Gamma^{\rm sh}(\varrho_\Gamma).
\]
The image and determinant entries are
\[
  \psi_\sigma=\psi_{\Gamma,0},
  \qquad
  {\cal M}_{\rm img}\psi_\sigma=m_\sigma\psi_\sigma,
  \qquad
  N_\Gamma^\chi=2\sigma_\chi d_\chi F(z_\chi).
\]
An endpoint-matrix realization makes the image mass a spectral readout:
\[
  {\cal M}_{\rm img}
  =
  {\cal Q}_{\rm img}^\dagger{\cal Q}_{\rm img},
  \qquad
  {\cal Q}_{\rm img}
  =
  {\bf 1}_{\CP^2\times S^1_Q}\otimes B_\sigma .
\]
The product density cancels from the normalized image quotient:
\[
  m_\sigma
  =
  \frac{
  \langle\eta_Q,B_\sigma^\dagger B_\sigma\eta_Q\rangle_{E_Q}}
  {\langle\eta_Q,\eta_Q\rangle_{E_Q}} .
\]
The one-line image sector is the eigenline condition
\[
  B_\sigma^\dagger B_\sigma\eta_Q=m_\sigma\eta_Q,
  \qquad
  \Pi_\sigma^\perp B_\sigma^\dagger B_\sigma\eta_Q=0,
  \qquad
  m_\sigma>0 .
\]
Let
\[
  \Pi_{\eta_Q}
  =
  \frac{|\eta_Q\rangle\langle\eta_Q|}
  {\langle\eta_Q,\eta_Q\rangle_{E_Q}},
  \qquad
  \Pi_{\eta_Q}^\perp={\bf 1}_{E_Q}-\Pi_{\eta_Q}.
\]
The finite endpoint normal form is
\[
  B_\sigma^\dagger B_\sigma
  =
  m_\sigma\Pi_{\eta_Q}
  +
  {\cal B}_{\sigma,\perp}^\dagger{\cal B}_{\sigma,\perp},
  \qquad
  {\cal B}_{\sigma,\perp}
  =
  \Pi_{\eta_Q}^\perp{\cal B}_{\sigma,\perp}\Pi_{\eta_Q}^\perp .
\]
The one-mode reduction has a separated complement whenever
\[
  {\cal B}_{\sigma,\perp}^\dagger{\cal B}_{\sigma,\perp}
  \ge
  M_{\sigma,\perp}^2\Pi_{\eta_Q}^\perp,
  \qquad
  M_{\sigma,\perp}^2>m_\sigma .
\]
Equivalently, define
\[
  g_{\sigma,\perp}
  =
  \frac{M_{\sigma,\perp}^2}{m_\sigma}.
\]
The finite endpoint image row passes exactly for
\[
  g_{\sigma,\perp}>1 .
  \tag{5.3}
\]
For a diagonal endpoint spectrum on \(\Pi_{\eta_Q}^\perp E_Q\),
\[
  B_\sigma^\dagger B_\sigma e_{\sigma,\alpha}
  =
  \lambda_{\sigma,\alpha}e_{\sigma,\alpha},
  \qquad
  e_{\sigma,\alpha}\in\Pi_{\eta_Q}^\perp E_Q,
\]
the ratio is
\[
  g_{\sigma,\perp}
  =
  \frac{\min_\alpha\lambda_{\sigma,\alpha}}{m_\sigma}.
  \tag{5.4}
\]
The finite eigenline residual is
\[
  \varepsilon_{\sigma,Q}^2
  =
  \frac{
  \left\|
  \Pi_{\eta_Q}^\perp B_\sigma^\dagger B_\sigma\eta_Q
  \right\|_{E_Q}^2}
  {m_\sigma^2\|\eta_Q\|_{E_Q}^2}.
\]
The one-line image condition is \(\varepsilon_{\sigma,Q}=0\).
For the minimal positive endpoint block,
\[
  {\cal B}_\chi^\dagger{\cal B}_\chi
  =
  M_\chi^2{\bf 1}_{E_\chi},
  \qquad
  E_\chi=\mathbb C,
  \qquad
  \sigma_\chi d_\chi=+1 .
\]
The determinant threshold coordinate is \(z_\chi=M_\chi L_H\).  Combining
with the density target gives
\[
  F(z_\chi)
  =
  \Xi_\Gamma^{\rm cp},
  \qquad
  \Xi_\Gamma^{\rm cp}
  =
  \frac{21375.716460702053\ldots\,\lambda_F}{2w_H^8}
  \left(
  \frac{L_Q{\mathfrak z}_{\Gamma}^{\CP^2}}{v_{\rm EW}}
  \right)^4 .
\]
Since
\[
  F'(z)=-\frac13e^{-z}z(1+z),
\]
the minimal positive endpoint row passes exactly for
\[
  0<\Xi_\Gamma^{\rm cp}\le1,
  \qquad
  z_\chi=F^{-1}(\Xi_\Gamma^{\rm cp}).
\]
The same determinant equation gives the neutral-density upper constraint
\[
  0<Z_{\Gamma,0}
  =
  L_Q{\mathfrak z}_{\Gamma}^{\CP^2}
  \le
  Z_{\chi,\max}.
\]
Here
\[
  \begin{aligned}
  Z_{\chi,\max}
  &=
  v_{\rm EW}
  \left(
  \frac{2w_H^8}{21375.716460702053\ldots\,\lambda_F}
  \right)^{1/4}  \\
  &=
  0.0983506713707568\ldots\,
  w_H^2\lambda_F^{-1/4}v_{\rm EW}.
  \end{aligned}
\]
Together with the density-shape lower row, this gives the loop-length
envelope
\[
  0<L_Q
  \le
  \frac{Z_{\chi,\max}}{{\mathfrak z}_{\Gamma}^{\CP^2}}
  \le
  \frac{16\pi^2 Z_{\chi,\max}}
  {441\,u_\Gamma^{\rm sh}(\varrho_\Gamma)} .
\]
Equivalently, the CP2 radius must lie in the band
\[
  \frac{441\,u_\Gamma^{\rm sh}(\varrho_\Gamma)}
  {8\pi^2\kappa_8e^{-\phi_0}\Omega_\Gamma}
  \le
  r_\Gamma^2
  \le
  \frac{2Z_{\chi,\max}}
  {\kappa_8e^{-\phi_0}\Omega_\Gamma L_Q}.
\]
For the full sector-sum row define the weighted radius moment
\[
  {\cal R}_\Gamma^2
  =
  \sum_{s\in{\cal S}_\Gamma}
  \omega_s e^{-\phi_s}r_{\Gamma,s}^2 .
\]
Equivalently, introduce the sector label space
\[
  {\cal H}_{\cal S}
  =
  \bigoplus_{s\in{\cal S}_\Gamma}\mathbb C|s\rangle
\]
and diagonal sector operators
\[
  {\sf W}_\Gamma|s\rangle
  =
  \omega_s e^{-\phi_s}|s\rangle,
  \qquad
  {\sf R}_\Gamma|s\rangle
  =
  r_{\Gamma,s}^2|s\rangle .
\]
Then
\[
  e^{-\phi_0}\Omega_\Gamma={\rm Tr}_{\cal S}{\sf W}_\Gamma,
  \qquad
  {\cal R}_\Gamma^2
  =
  {\rm Tr}_{\cal S}({\sf W}_\Gamma{\sf R}_\Gamma).
\]
Then
\[
  {\mathfrak z}_{\Gamma}^{\CP^2}
  =
  \frac{\kappa_8}{2}{\cal R}_\Gamma^2,
  \qquad
  e^{-\phi_0}\Omega_\Gamma r_\Gamma^2
  =
  {\cal R}_\Gamma^2
  \quad
  \hbox{in the common-metric row}.
\]
The sector-sum compactification band is
\[
  \frac{441\,u_\Gamma^{\rm sh}(\varrho_\Gamma)}
  {8\pi^2\kappa_8}
  \le
  {\cal R}_\Gamma^2
  \le
  \frac{2Z_{\chi,\max}}{\kappa_8L_Q}.
\]
The determinant threshold depends on the weighted product
\[
  {\cal V}_{\Gamma,Q}
  =
  L_Q{\cal R}_\Gamma^2,
  \qquad
  Z_{\Gamma,0}
  =
  \frac{\kappa_8}{2}{\cal V}_{\Gamma,Q}.
\]
In trace form the local density outputs read
\[
  Z_{\Gamma,0}
  =
  \frac{\kappa_8L_Q}{2}
  {\rm Tr}_{\cal S}({\sf W}_\Gamma{\sf R}_\Gamma),
  \qquad
  Z_{\Gamma,Q}
  =
  8\kappa_8L_Q
  {\rm Tr}_{\cal S}({\sf W}_\Gamma{\sf R}_\Gamma),
\]
\[
  Z_{\Gamma,H}
  =
  \frac{Q\kappa_8L_Q}{8}
  {\rm Tr}_{\cal S}({\sf W}_\Gamma{\sf R}_\Gamma).
\]
Set
\[
  {\cal T}_\Gamma
  =
  {\rm Tr}_{\cal S}({\sf W}_\Gamma{\sf R}_\Gamma).
\]
The trace pass surface is
\[
  {\cal T}_\Gamma
  \ge
  {\cal T}_{\Gamma,\min}(\varrho_\Gamma)
  =
  \frac{441\,u_\Gamma^{\rm sh}(\varrho_\Gamma)}
  {8\pi^2\kappa_8},
  \qquad
  0<L_Q\le
  \frac{2Z_{\chi,\max}}{\kappa_8{\cal T}_\Gamma}.
\]
Equivalently, with \(0<x_{\Gamma,Q}\le1\),
\[
  L_Q
  =
  \frac{2Z_{\chi,\max}}{\kappa_8{\cal T}_\Gamma}
  x_{\Gamma,Q},
  \qquad
  {\cal T}_\Gamma\ge
  {\cal T}_{\Gamma,\min}(\varrho_\Gamma).
\]
At lower-density saturation,
\[
  {\cal T}_\Gamma={\cal T}_{\Gamma,\min}(\varrho_\Gamma),
  \qquad
  L_Q=x_{\Gamma,Q}L_{Q,\max}(\varrho_\Gamma).
\]
Define
\[
  \chi_\Gamma
  =
  \frac{{\cal T}_\Gamma}
  {{\cal T}_{\Gamma,\min}(\varrho_\Gamma)}
  \ge1 .
\]
Then
\[
  L_Q
  =
  \frac{x_{\Gamma,Q}}{\chi_\Gamma}
  L_{Q,\max}(\varrho_\Gamma),
  \qquad
  x_{\Gamma,Q}
  =
  \chi_\Gamma
  \frac{L_Q}{L_{Q,\max}(\varrho_\Gamma)} .
\]
Therefore
\[
  0<L_Q\le
  \frac{L_{Q,\max}(\varrho_\Gamma)}{\chi_\Gamma}.
\]
In sector variables set
\[
  \zeta_{\Gamma,s}
  =
  \frac{\omega_s e^{-\phi_s}r_{\Gamma,s}^2}
  {{\cal T}_{\Gamma,\min}(\varrho_\Gamma)}
  \ge0 .
\]
Then
\[
  \chi_\Gamma
  =
  \sum_{s\in{\cal S}_\Gamma}\zeta_{\Gamma,s},
  \qquad
  L_Q
  =
  \frac{x_{\Gamma,Q}}
  {\sum_{s\in{\cal S}_\Gamma}\zeta_{\Gamma,s}}
  L_{Q,\max}(\varrho_\Gamma).
  \tag{6.1}
\]
For a fixed loop length,
\[
  x_{\Gamma,Q}
  =
  \frac{L_Q}{L_{Q,\max}(\varrho_\Gamma)}
  \sum_{s\in{\cal S}_\Gamma}\zeta_{\Gamma,s},
  \qquad
  1\le
  \sum_{s\in{\cal S}_\Gamma}\zeta_{\Gamma,s}
  \le
  \frac{L_{Q,\max}(\varrho_\Gamma)}{L_Q}.
  \tag{6.2}
\]
Define sector fractions
\[
  p_{\Gamma,s}
  =
  \frac{\zeta_{\Gamma,s}}{\chi_\Gamma},
  \qquad
  p_{\Gamma,s}\ge0,
  \qquad
  \sum_{s\in{\cal S}_\Gamma}p_{\Gamma,s}=1 .
\]
Then each retained sector obeys
\[
  \omega_s e^{-\phi_s}r_{\Gamma,s}^2
  =
  p_{\Gamma,s}\chi_\Gamma
  {\cal T}_{\Gamma,\min}(\varrho_\Gamma)
  =
  p_{\Gamma,s}\chi_\Gamma
  \frac{441\,u_\Gamma^{\rm sh}(\varrho_\Gamma)}
  {8\pi^2\kappa_8}.
  \tag{6.3}
\]
Equivalently,
\[
  r_{\Gamma,s}^2
  =
  \frac{441\,p_{\Gamma,s}\chi_\Gamma
  u_\Gamma^{\rm sh}(\varrho_\Gamma)}
  {8\pi^2\kappa_8\omega_s e^{-\phi_s}} .
  \tag{6.4}
\]
Collecting the local rows, define the finite D8/O8 data record
\[
  {\cal D}_\Gamma
  =
  \left(
  {\cal S}_\Gamma;\,
  \omega_s,\phi_s,p_{\Gamma,s},b_{\Gamma,s},\rho_{\Gamma,s};\,
  \chi_\Gamma,x_{\Gamma,Q},L_Q;\,
  m_\sigma,\varepsilon_{\sigma,Q},\lambda_{\sigma,\alpha}
  \right).
\]
It gives
\[
  \varrho_\Gamma
  =
  \sum_{s\in{\cal S}_{\Gamma,0}}b_{\Gamma,s}\rho_{\Gamma,s},
  \qquad
  {\cal T}_\Gamma
  =
  \chi_\Gamma{\cal T}_{\Gamma,\min}(\varrho_\Gamma),
  \qquad
  L_Q
  =
  \frac{x_{\Gamma,Q}}{\chi_\Gamma}
  L_{Q,\max}(\varrho_\Gamma),
  \tag{6.5}
\]
\[
  z_\chi=F^{-1}(x_{\Gamma,Q}^4),
  \qquad
  \varepsilon_{\sigma,Q}^2
  =
  \frac{
  \left\|
  \Pi_{\eta_Q}^\perp B_\sigma^\dagger B_\sigma\eta_Q
  \right\|_{E_Q}^2}
  {m_\sigma^2\|\eta_Q\|_{E_Q}^2},
  \qquad
  g_{\sigma,\perp}
  =
  \frac{\min_\alpha\lambda_{\sigma,\alpha}}{m_\sigma}.
  \tag{6.6}
\]
The admissible finite row satisfies
\[
  \chi_\Gamma\ge1,\quad
  0<x_{\Gamma,Q}\le1,\quad
  0<\varrho_\Gamma<\varrho_c,\quad
  \varepsilon_{\sigma,Q}=0,\quad
  g_{\sigma,\perp}>1,
  \tag{6.7}
\]
\[
  \sum_{s\in{\cal S}_\Gamma}p_{\Gamma,s}=1,
  \qquad
  \sum_{s\in{\cal S}_{\Gamma,0}}b_{\Gamma,s}=1,
  \qquad
  p_{\Gamma,s},b_{\Gamma,s},\rho_{\Gamma,s}\ge0 .
  \tag{6.8}
\]
The corresponding localized-action map starts from the quadratic data
\[
  {\cal Q}_\Gamma^{(2)}
  =
  \left(
  {\cal S}_\Gamma;\,
  \omega_s,\phi_s,\gamma_{\Gamma,s}^{\CP^2};\,
  \theta_{\Gamma,s}^3,P_0^\Gamma,P_h^\Gamma,\Pi_Q;\,
  B_\sigma,\eta_Q,L_Q
  \right).
\]
Let \(\widehat\gamma_s\) be the unit \(\CP^2\) reference metric in sector
\(s\), and define
\[
  {\rm Av}_{s,\Pi_Q}(f)
  =
  \frac{
  \int_{\CP^2_s}\sqrt{\widehat\gamma_s}\,
  {\rm tr}_{E_s}(\Pi_Qf\Pi_Q)}
  {\int_{\CP^2_s}\sqrt{\widehat\gamma_s}\,
  {\rm tr}_{E_s}(\Pi_Q)} .
  \tag{6.9}
\]
For each retained sector, set
\[
  r_{\Gamma,s}^2
  =
  {\rm Av}_{s,\Pi_Q}\!\left(
  \frac14{\rm tr}_{\widehat\gamma_s}
  \gamma_{\Gamma,s}^{\CP^2}
  \right),
  \qquad
  {\cal T}_\Gamma
  =
  \sum_s\omega_s e^{-\phi_s}r_{\Gamma,s}^2 .
  \tag{6.10}
\]
Then
\[
  \varrho_\Gamma
  =
  \frac{\sum_sC_{h,s}^\Gamma}{\sum_sC_{0,s}^\Gamma},
  \qquad
  \chi_\Gamma
  =
  \frac{{\cal T}_\Gamma}
  {{\cal T}_{\Gamma,\min}(\varrho_\Gamma)},
  \qquad
  x_{\Gamma,Q}
  =
  \frac{\kappa_8L_Q{\cal T}_\Gamma}{2Z_{\chi,\max}} .
  \tag{6.11}
\]
The endpoint spectrum supplies
\[
  m_\sigma
  =
  \frac{\langle\eta_Q,B_\sigma^\dagger B_\sigma\eta_Q\rangle_{E_Q}}
  {\langle\eta_Q,\eta_Q\rangle_{E_Q}},
  \qquad
  \{\lambda_{\sigma,\alpha}\}
  =
  {\rm Spec}\!
  \left(
  \Pi_{\eta_Q}^{\perp}B_\sigma^\dagger B_\sigma
  \Pi_{\eta_Q}^{\perp}
  \right).
  \tag{6.12}
\]
The same endpoint operator gives
\[
  \varepsilon_{\sigma,Q}^2
  =
  \frac{
  \left\|
  \Pi_{\eta_Q}^\perp B_\sigma^\dagger B_\sigma\eta_Q
  \right\|_{E_Q}^2}
  {m_\sigma^2\|\eta_Q\|_{E_Q}^2}.
\]
Together with
\[
  p_{\Gamma,s}
  =
  \frac{\omega_s e^{-\phi_s}r_{\Gamma,s}^2}
  {{\cal T}_\Gamma},
  \qquad
  b_{\Gamma,s}
  =
  \frac{C_{0,s}^\Gamma}{\sum_tC_{0,t}^\Gamma},
  \qquad
  \rho_{\Gamma,s}
  =
  \frac{C_{h,s}^\Gamma}{C_{0,s}^\Gamma},
  \tag{6.13}
\]
these equations define
\[
  {\mathfrak E}_\Gamma({\cal Q}_\Gamma^{(2)})
  =
  {\cal D}_\Gamma .
  \tag{6.14}
\]
The algebraic inverse local packet is explicit on the support
\[
  {\cal S}_{\Gamma,+}=\{s:p_{\Gamma,s}>0\}.
\]
Choose \({\cal C}_\Gamma>0\), endpoint-normalized one-forms
\(\widehat e_{0,s}\) and \(\widehat e_{h,s}\), and set
\[
  \gamma_{\Gamma,s}^{\CP^2}=r_{\Gamma,s}^2\widehat\gamma_s,
  \qquad
  r_{\Gamma,s}^2
  =
  \frac{p_{\Gamma,s}\chi_\Gamma
  {\cal T}_{\Gamma,\min}(\varrho_\Gamma)}
  {\omega_s e^{-\phi_s}} .
  \tag{6.15}
\]
The coframe pieces may be chosen as
\[
  P_0^\Gamma\theta_{\Gamma,s}^3
  =
  \left(
  \frac{b_{\Gamma,s}{\cal C}_\Gamma}
  {\omega_s e^{-\phi_s}}
  \right)^{1/2}\widehat e_{0,s},
  \qquad
  P_h^\Gamma\theta_{\Gamma,s}^3
  =
  \left(
  \frac{b_{\Gamma,s}\rho_{\Gamma,s}{\cal C}_\Gamma}
  {\omega_s e^{-\phi_s}}
  \right)^{1/2}\widehat e_{h,s}.
  \tag{6.16}
\]
With \(e_0=\eta_Q/\|\eta_Q\|_{E_Q}\) and an orthonormal complement
\(\{e_\alpha\}\subset\Pi_{\eta_Q}^\perp E_Q\), take
\[
  B_\sigma^\dagger B_\sigma
  =
  m_\sigma |e_0\rangle\langle e_0|
  +
  \sum_\alpha
  \lambda_{\sigma,\alpha}
  |e_\alpha\rangle\langle e_\alpha|.
  \tag{6.17}
\]
Then
\[
  {\mathfrak I}_\Gamma({\cal D}_\Gamma;{\cal C}_\Gamma,\widehat e)
  =
  {\cal Q}_\Gamma^{(2)},
  \qquad
  {\mathfrak E}_\Gamma
  \!\left(
  {\mathfrak I}_\Gamma({\cal D}_\Gamma;{\cal C}_\Gamma,\widehat e)
  \right)
  =
  {\cal D}_\Gamma .
  \tag{6.18}
\]
The charged determinant extension uses the same threshold coordinate.  Let
\[
  z_\chi=F^{-1}(x_{\Gamma,Q}^4),
  \qquad
  M_\chi=\frac{z_\chi}{L_H},
  \qquad
  E_\chi=\mathbb C,
  \qquad
  \sigma_\chi d_\chi=+1 .
  \tag{6.19}
\]
Set
\[
  P_1^\chi={\rm span}\{u,v\},
  \qquad
  P_2^\chi={\rm span}\{u^2,v^2\},
  \tag{6.20}
\]
and
\[
  B_\chi^\dagger B_\chi=M_\chi^2{\bf 1}_{E_\chi},
  \qquad
  {\cal B}_{r,\Gamma}^\chi
  =
  {\bf 1}_{P_r^\chi}\otimes B_\chi,
  \qquad r=1,2 .
  \tag{6.21}
\]
Then
\[
  {\cal N}_\Gamma^\chi
  =
  2F(M_\chi L_H)
  =
  2x_{\Gamma,Q}^4 .
  \tag{6.22}
\]
The extended local packet is
\[
  \widetilde{\cal Q}_\Gamma^{(2)}
  =
  \left(
  {\cal Q}_\Gamma^{(2)};\,
  P_1^\chi,P_2^\chi,E_\chi,B_\chi,L_H,\sigma_\chi d_\chi
  \right).
  \tag{6.23}
\]
The finite package carries the density, coframe, image, and determinant rows.
For concrete localized data
\[
  {\cal A}_\Gamma
  =
  \left(
  P_s[G+B],F_s,\phi_s,\omega_s,\Pi_s;\,
  \iota_\Gamma^*C_1;\,
  B_\sigma,B_\chi,L_Q,L_H
  \right)_{s\in{\cal S}_\Gamma},
\]
the action-matching residual is the finite vector
\[
  {\cal R}_\Gamma({\cal A}_\Gamma;{\cal D}_\Gamma)
  =
  \left(
  R_T,\ R_x,\ R_\varrho,\ R_\epsilon,\ R_g,\ R_\chi
  \right),
  \tag{6.24}
\]
where
\[
  R_T
  =
  \frac{{\cal T}_\Gamma({\cal A}_\Gamma)}
  {{\cal T}_{\Gamma,\min}(\varrho_\Gamma)}
  -\chi_\Gamma,
  \qquad
  R_x
  =
  \frac{\kappa_8L_Q{\cal T}_\Gamma({\cal A}_\Gamma)}
  {2Z_{\chi,\max}}
  -x_{\Gamma,Q},
  \tag{6.25}
\]
\[
  R_\varrho
  =
  \frac{C_h^\Gamma({\cal A}_\Gamma)}
  {C_0^\Gamma({\cal A}_\Gamma)}
  -\varrho_\Gamma,
  \qquad
  R_\epsilon=\varepsilon_{\sigma,Q},
  \tag{6.26}
\]
\[
  R_g
  =
  \max\!\left\{0,\,
  1-\frac{\min_\alpha\lambda_{\sigma,\alpha}}{m_\sigma}
  \right\},
  \qquad
  R_\chi
  =
  \frac{M_\chi L_H-z_\chi}{1+z_\chi}.
  \tag{6.27}
\]
The localized action produces the finite packet exactly when
\[
  {\cal R}_\Gamma({\cal A}_\Gamma;{\cal D}_\Gamma)=0 .
  \tag{6.28}
\]
A localized finite normal form realizing an admissible row is explicit.  Write
\[
  W_s=\omega_s e^{-\phi_s},
  \qquad
  {\cal S}_{\Gamma,+}=\{s:p_{\Gamma,s}>0\}.
\]
On \({\cal S}_{\Gamma,+}\), take endpoint-normalized one-forms
\(\widehat e_{0,s},\widehat e_{h,s}\) and set
\[
  P_s[G+B]_{\CP^2}
  =
  r_{\Gamma,s}^2\widehat\gamma_s,
  \qquad
  r_{\Gamma,s}^2
  =
  \frac{p_{\Gamma,s}\chi_\Gamma
  {\cal T}_{\Gamma,\min}(\varrho_\Gamma)}
  {W_s}.
\]
The retained current pullback is
\[
  \Pi_{\rm EW}\iota_\Gamma^*C_1|_s
  =
  q_\Gamma\theta_{\Gamma,s}^3,
  \qquad q_\Gamma>0,
\]
with
\[
  P_0^\Gamma\theta_{\Gamma,s}^3
  =
  \left(\frac{b_{\Gamma,s}{\cal C}_\Gamma}{W_s}\right)^{1/2}
  \widehat e_{0,s},
  \qquad
  P_h^\Gamma\theta_{\Gamma,s}^3
  =
  \left(\frac{b_{\Gamma,s}\rho_{\Gamma,s}{\cal C}_\Gamma}{W_s}\right)^{1/2}
  \widehat e_{h,s}.
\]
The endpoint image and determinant blocks are
\[
  B_\sigma^\dagger B_\sigma
  =
  m_\sigma |e_0\rangle\langle e_0|
  +
  \sum_\alpha
  \lambda_{\sigma,\alpha}|e_\alpha\rangle\langle e_\alpha|,
  \qquad
  \frac{\min_\alpha\lambda_{\sigma,\alpha}}{m_\sigma}>1,
\]
\[
  B_\chi^\dagger B_\chi
  =
  \left(\frac{z_\chi}{L_H}\right)^2{\bf 1}_{E_\chi},
  \qquad
  L_Q
  =
  \frac{x_{\Gamma,Q}}{\chi_\Gamma}L_{Q,\max}(\varrho_\Gamma).
\]
Then
\[
  {\cal T}_\Gamma({\cal A}_\Gamma^{\rm nf})
  =
  \chi_\Gamma{\cal T}_{\Gamma,\min}(\varrho_\Gamma),
  \qquad
  \frac{C_h^\Gamma({\cal A}_\Gamma^{\rm nf})}
  {C_0^\Gamma({\cal A}_\Gamma^{\rm nf})}
  =
  \varrho_\Gamma,
\]
\[
  \varepsilon_{\sigma,Q}=0,
  \qquad
  M_\chi L_H=z_\chi,
  \qquad
  {\cal R}_\Gamma({\cal A}_\Gamma^{\rm nf};{\cal D}_\Gamma)=0 .
\]
A global Type IIA source datum can be written as
\[
  {\mathfrak G}_\Gamma
  =
  \left(
  X_{10},{\cal I}_{O8},
  \iota_{\Gamma,s},E_s,F_s,\phi_s,G,B,C_1,F_0,\Omega
  \right)_{s\in{\cal S}_\Gamma},
\]
with localized DBI/CS/Payen readout
\[
  {\mathfrak L}_\Gamma({\mathfrak G}_\Gamma)={\cal A}_\Gamma .
\]
The finite normal form is the local image of this global datum when
\[
  {\mathfrak L}_\Gamma({\mathfrak G}_\Gamma)={\cal A}_\Gamma^{\rm nf}.
\]
Component form:
\[
  P_s[G+B]_{\CP^2}=r_{\Gamma,s}^2\widehat\gamma_s,
  \qquad
  \Pi_{\rm EW}\iota_{\Gamma,s}^*C_1=q_\Gamma\theta_{\Gamma,s}^3,
\]
\[
  \Pi_{\rm EW}\iota_{\Gamma,s}^*F_2
  =
  q_\Gamma d_\Gamma\theta_{\Gamma,s}^3,
  \qquad
  {\rm End}_\Omega(E_s,F_s)
  =
  (\Pi_Q,e_0,E_\chi,B_\sigma,B_\chi).
\]
The source and lower-charge rows are
\[
  dF_0
  =
  \sum_sN_{8,s}\delta_{\Gamma_s}
  +
  \sum_aQ_{O8,a}\delta_{O8,a},
  \qquad
  \sum_sN_{8,s}+\sum_aQ_{O8,a}=0,
\]
\[
  Q_{\rm lower}^{\rm tot}
  =
  \left[
  \sum_s{\rm ch}(E_s)e^{F_s/2\pi}
  \sqrt{\frac{\widehat A(T\Gamma_s)}
  {\widehat A(N\Gamma_s)}}
  +
  \sum_aQ_{O8,a}^{\rm curv}
  \right]_{<8}
  =
  0 .
\]
Equivalently, define
\[
  {\cal R}_\Gamma^{\rm glob}({\mathfrak G}_\Gamma)
  =
  \left(
  {\mathfrak L}_\Gamma({\mathfrak G}_\Gamma)-{\cal A}_\Gamma^{\rm nf},
  dF_0-\sum_sN_{8,s}\delta_{\Gamma_s}
  -\sum_aQ_{O8,a}\delta_{O8,a},
  Q_{\rm lower}^{\rm tot}
  \right).
\]
Then
\[
  {\cal R}_\Gamma^{\rm glob}({\mathfrak G}_\Gamma)=0
  \quad\Longrightarrow\quad
  {\cal R}_\Gamma({\cal A}_\Gamma;{\cal D}_\Gamma)=0 .
\]
The minimal endpoint-real type-I' datum has
\[
  Q_{O8,L}=Q_{O8,R}=-16,
  \qquad
  N_{8,L}=N_{8,R}=16.
\]
Thus
\[
  \sum_{\epsilon=L,R}(N_{8,\epsilon}+Q_{O8,\epsilon})=0,
  \qquad
  k_\epsilon=\frac{N_{8,\epsilon}}{2}-8=0.
\]
At \(Q=7\),
\[
  W_7=H^0(\CP^1,\mathcal O(7)),
  \qquad
  \dim W_7=8,
\]
\[
  E_7=(W_7\oplus W_7^\vee)^\Omega,
  \qquad
  {\rm rk}_{\mathbb C}(E_7^{\mathbb C})=16.
\]
The minimal spin\(^c\) \(\CP^2\) source has \(m=0\), \(f_0=13/2\), and
\[
  {\mathfrak q}_{C_1}
  =
  16\left(\frac{f_0^2}{2}-\frac1{16}\right)
  =
  337.
\]
For the two endpoints,
\[
  Q_{D8/O8}^{<8}
  =
  208h_Q+673h_Q^2.
\]
A bulk lower-source completion
\[
  Q_{\rm bulk}^{<8}
  =
  -208h_Q-673h_Q^2
\]
gives
\[
  Q_{\rm lower}^{\rm tot}=0.
\]
Consequently, a datum \({\mathfrak G}_\Gamma^{(7)}\) with
\[
  {\mathfrak L}_\Gamma({\mathfrak G}_\Gamma^{(7)})
  =
  {\cal A}_\Gamma^{\rm nf}
\]
satisfies
\[
  {\cal R}_\Gamma^{\rm glob}({\mathfrak G}_\Gamma^{(7)})=0.
\]
The localized type-I' readout realizing the finite normal form is explicit.
Set
\[
  W_s=\omega_s e^{-\phi_s},
  \qquad
  P_sB_{\CP^2}=0,
\]
and choose
\[
  P_sG_{\CP^2}
  =
  r_{\Gamma,s}^2\widehat\gamma_s,
  \qquad
  W_s r_{\Gamma,s}^2
  =
  p_{\Gamma,s}\chi_\Gamma
  {\cal T}_{\Gamma,\min}(\varrho_\Gamma).
\]
Then
\[
  {\cal T}_\Gamma^{\rm rd}
  =
  \sum_sW_s r_{\Gamma,s}^2
  =
  \chi_\Gamma{\cal T}_{\Gamma,\min}(\varrho_\Gamma).
\]
The RR/current row is
\[
  \Pi_{\rm EW}\iota_{\Gamma,s}^*C_1
  =
  q_\Gamma\theta_{\Gamma,s}^3,
  \qquad
  \Pi_{\rm EW}\iota_{\Gamma,s}^*F_2
  =
  q_\Gamma d_\Gamma\theta_{\Gamma,s}^3.
\]
Choose endpoint-normalized coframes by
\[
  P_0^\Gamma\theta_{\Gamma,s}^3
  =
  \left(\frac{b_{\Gamma,s}{\cal C}_\Gamma}{W_s}\right)^{1/2}
  \widehat e_{0,s},
  \qquad
  P_h^\Gamma\theta_{\Gamma,s}^3
  =
  \left(\frac{b_{\Gamma,s}\rho_{\Gamma,s}{\cal C}_\Gamma}{W_s}\right)^{1/2}
  \widehat e_{h,s}.
\]
The readout gives
\[
  C_{0,s}^{\Gamma,{\rm rd}}=b_{\Gamma,s}{\cal C}_\Gamma,
  \qquad
  C_{h,s}^{\Gamma,{\rm rd}}
  =
  b_{\Gamma,s}\rho_{\Gamma,s}{\cal C}_\Gamma,
\]
and hence
\[
  \frac{\sum_sC_{h,s}^{\Gamma,{\rm rd}}}
  {\sum_sC_{0,s}^{\Gamma,{\rm rd}}}
  =
  \sum_s b_{\Gamma,s}\rho_{\Gamma,s}
  =
  \varrho_\Gamma.
\]
The endpoint operators are
\[
  B_\sigma^\dagger B_\sigma
  =
  m_\sigma |e_0\rangle\langle e_0|
  +
  \sum_\alpha
  \lambda_{\sigma,\alpha}|e_\alpha\rangle\langle e_\alpha|,
  \qquad
  \min_\alpha\lambda_{\sigma,\alpha}>m_\sigma,
\]
\[
  B_\chi^\dagger B_\chi
  =
  \left(\frac{z_\chi}{L_H}\right)^2{\bf 1}_{E_\chi},
  \qquad
  L_Q=
  \frac{x_{\Gamma,Q}}{\chi_\Gamma}L_{Q,\max}(\varrho_\Gamma).
\]
The residual coordinates satisfy
\[
  R_T=R_\varrho=R_\epsilon=R_g=R_\chi=0,
\]
and
\[
  R_x
  =
  \frac{\kappa_8L_Q\chi_\Gamma
  {\cal T}_{\Gamma,\min}(\varrho_\Gamma)}
  {2Z_{\chi,\max}}
  -x_{\Gamma,Q}
  =
  0.
\]
Therefore
\[
  {\mathfrak L}_\Gamma({\mathfrak G}_{\Gamma,{\rm loc}}^{(7)})
  =
  {\cal A}_\Gamma^{\rm nf}.
\]
The extension from the local readout to a global type-I' representative
separates into a period condition and an endpoint spectral condition.  Let
\(\omega_Q\) be the integral generator on \(S_Q^1\), normalized by
\(\int_{S_Q^1}\omega_Q=1\), and set
\[
  \ell_{\Gamma,s}
  =
  \int_{S_Q^1}\theta_{\Gamma,s}^3 .
\]
The RR one-form periods obey
\[
  R_{{\rm per},s}
  =
  \frac{q_\Gamma\ell_{\Gamma,s}}{2\pi}
  -n_{\Gamma,s}-\nu_{\Gamma,s},
  \qquad
  n_{\Gamma,s}\in{\mathbb Z},
  \qquad
  \nu_{\Gamma,s}\in{\mathbb Q}/{\mathbb Z}.
\]
A common \(q_\Gamma\) lies on the rational-period locus
\[
  q_\Gamma\ell_{\Gamma,s}
  =
  2\pi(n_{\Gamma,s}+\nu_{\Gamma,s})
  \qquad
  \hbox{for every retained }s.
\]
The metric, dilaton, and \(B\)-field rows extend through disjoint endpoint
collars by partition of unity with preservation of the local readout values.
For the endpoint operator, let \(D_{\partial,\Gamma}^{(7)}\) be the real
boundary Dirac--Dolbeault operator on
\[
  E_7\oplus E_\sigma\oplus
  (P_1^\chi\oplus P_2^\chi)\otimes E_\chi .
\]
The endpoint spectral condition is
\[
  R_{\rm DD}
  =
  \left\|
  \Pi_{\eta_Q}D_{\partial,\Gamma}^{(7)\,2}\Pi_{\eta_Q}
  -m_\sigma\Pi_{\eta_Q}
  \right\|
  +
  \left\|
  \Pi_{\eta_Q}^{\perp}D_{\partial,\Gamma}^{(7)\,2}\Pi_{\eta_Q}
  \right\|
\]
\[
  \quad+
  \left\|
  \Pi_\chi D_{\partial,\Gamma}^{(7)\,2}\Pi_\chi
  -
  \left(\frac{z_\chi}{L_H}\right)^2\Pi_\chi
  \right\|,
  \qquad
  \Pi_\chi=P_1^\chi\oplus P_2^\chi .
\]
The global representative satisfies
\[
  R_{{\rm per},s}=0\quad(\forall s),
  \qquad
  R_{\rm DD}=0,
  \qquad
  Q_{\rm lower}^{\rm tot}=0 .
\]
Thus
\[
  \Xi_\Gamma^{\rm cp}
  =
  \left(
  \frac{Z_{\Gamma,0}}{Z_{\chi,\max}}
  \right)^4
  =
  \left(
  \frac{\kappa_8{\cal V}_{\Gamma,Q}}
  {2Z_{\chi,\max}}
  \right)^4 .
\]
The determinant pass is equivalently
\[
  0<{\cal V}_{\Gamma,Q}\le
  \frac{2Z_{\chi,\max}}{\kappa_8},
  \qquad
  z_\chi
  =
  F^{-1}\!\left[
  \left(
  \frac{\kappa_8{\cal V}_{\Gamma,Q}}
  {2Z_{\chi,\max}}
  \right)^4
  \right].
\]
The normalized endpoint coordinate is
\[
  x_{\Gamma,Q}
  =
  \frac{Z_{\Gamma,0}}{Z_{\chi,\max}}
  =
  \frac{\kappa_8{\cal V}_{\Gamma,Q}}{2Z_{\chi,\max}}
  =
  \frac{\kappa_8L_Q{\cal T}_\Gamma}{2Z_{\chi,\max}} .
\]
Then
\[
  \Xi_\Gamma^{\rm cp}=x_{\Gamma,Q}^4,
  \qquad
  z_\chi=F^{-1}(x_{\Gamma,Q}^4),
  \qquad
  0<x_{\Gamma,Q}\le1 .
\]
Combining the sector-sum lower row with the determinant upper row gives
\[
  \frac{L_Q}{L_{Q,\max}(\varrho_\Gamma)}
  \le
  x_{\Gamma,Q}
  \le
  1 .
\]
Define
\[
  r_{\min}^2(\varrho_\Gamma)
  =
  \frac{441\,u_\Gamma^{\rm sh}(\varrho_\Gamma)}
  {8\pi^2\kappa_8e^{-\phi_0}\Omega_\Gamma},
  \qquad
  r_{\max}^2(L_Q)
  =
  \frac{2Z_{\chi,\max}}
  {\kappa_8e^{-\phi_0}\Omega_\Gamma L_Q}.
\]
The compactification interval
\[
  {\cal I}_{r^2}(\varrho_\Gamma,L_Q)
  =
  [r_{\min}^2(\varrho_\Gamma),r_{\max}^2(L_Q)]
\]
is populated exactly when
\[
  0<L_Q\le L_{Q,\max}(\varrho_\Gamma),
  \qquad
  L_{Q,\max}(\varrho_\Gamma)
  =
  \frac{16\pi^2Z_{\chi,\max}}
  {441\,u_\Gamma^{\rm sh}(\varrho_\Gamma)} .
\]

The matching endpoint determinant spectrum is a paired charged oscillator:
\[
  {\cal E}_{1}^{\chi}
  =
  {\rm span}\{u,v\}\otimes E_\chi,
  \qquad
  {\cal E}_{2}^{\chi}
  =
  {\rm span}\{u^2,v^2\}\otimes E_\chi,
\]
with
\[
  ({\cal B}_{r,\Gamma}^\chi)^\dagger{\cal B}_{r,\Gamma}^\chi
  =
  M_\chi^2{\bf 1}_{2d_\chi},
  \qquad r=1,2 .
\]
For positive determinant sign this gives
\[
  {\cal N}_\Gamma^\chi=2d_\chi F(M_\chi L_H).
\]
The minimal row has \(d_\chi=1\), hence
\[
  {\cal N}_\Gamma^\chi=2F(z_\chi),
  \qquad z_\chi=M_\chi L_H .
\]
The local spectrum package identifies \(e_\sigma\) and the paired oscillator
block \(E_\chi\) inside the same D8/O8 boundary spectrum.

\subsection{Single-Block Boundary Complex}

A local complex that generates the spectrum package has Hilbert space
\[
  {\cal H}_\Gamma^{\rm loc}
  =
  {\rm span}\{e_\sigma\}
  \oplus
  \left(P_1^\chi\oplus P_2^\chi\right)\otimes E_\chi
  \oplus{\cal H}_\perp ,
\]
where
\[
  P_1^\chi={\rm span}\{u,v\},
  \qquad
  P_2^\chi={\rm span}\{u^2,v^2\}.
\]
Let
\[
  {\cal Q}_{\Gamma,\sigma}
  =
  \sqrt{m_\sigma}\,\Pi_{e_\sigma}
  \oplus
  \left({\bf 1}_{P_1^\chi}\oplus{\bf 1}_{P_2^\chi}\right)
  \otimes B_\chi
  \oplus Q_\perp .
\]
Then
\[
  {\cal M}_{\rm img}
  =
  {\cal Q}_{\Gamma,\sigma}^\dagger
  {\cal Q}_{\Gamma,\sigma}\big|_{{\rm span}\{e_\sigma\}},
  \qquad
  {\cal B}_{r,\Gamma}^\chi
  =
  {\bf 1}_{P_r^\chi}\otimes B_\chi .
\]
The current-source row is
\[
  {\cal J}_A^{\rm src}
  =
  C_{A\sigma}e_\sigma^\flat,
  \qquad A=3,Y,
\]
with
\[
  C_{3\sigma}=L_2\sqrt{m_\sigma\delta_\sigma},
  \qquad
  C_{Y\sigma}=\rho_\sigma C_{3\sigma}.
\]
The charged block obeys
\[
  B_\chi^\dagger B_\chi=M_\chi^2{\bf 1}_{E_\chi},
  \qquad
  {\rm Str}_{E_\chi}{\bf 1}=d_\chi,
  \qquad
  \sigma_\chi=+1 .
\]
It follows that
\[
  F(M_\chi L_H)=\frac{{\cal N}_\Gamma^\chi}{2d_\chi}.
\]
Using the density target,
\[
  F(M_\chi L_H)
  =
  \frac{21375.716460702053\ldots\,\lambda_F}
       {2d_\chi w_H^8}
  \left(\frac{Z_{\Gamma,0}}{v_{\rm EW}}\right)^4 .
\]
For \(d_\chi=1\), the right hand side is \(\Xi_\Gamma^{\rm cp}\).  Values in
\((0,1]\) determine
\[
  M_\chi=\frac{z_\chi}{L_H},
  \qquad
  F(z_\chi)=\Xi_\Gamma^{\rm cp}.
\]
The image mass \(m_\sigma\) enters the source coupling \(C_{A\sigma}\) and
cancels from \({\cal D}_{AB}\).  The local action supplies \(m_\sigma\) and
the profile \(e_\sigma\).

\subsection{Profile-To-Matrix Closure}

Identify the image eigenline with
\[
  e_\sigma=\psi_\sigma,
  \qquad
  \int_{\Gamma_{\rm int}}d\nu_\Gamma|\psi_\sigma|^2=1 .
\]
For the weighted current Gram matrix
\[
  {\cal G}_{AB}^{(\psi)}
  =
  \int_{\Gamma_{\rm int}}d\nu_\Gamma\,
  \psi_\sigma^2
  \gamma_\Gamma^{-1}(\theta_\Gamma^A,\theta_\Gamma^B),
  \qquad A,B=3,Y ,
\]
define
\[
  c_\sigma^B
  =
  \left({\cal G}^{(\psi)-1}\right)^{BA}(g_\sigma)_A .
\]
The minimal source profile is
\[
  \Xi_{\sigma,\min}
  =
  L_2\psi_\sigma c_\sigma^B\theta_{\Gamma B},
  \qquad
  \Theta_{{\rm img},\sigma}
  =
  \sqrt{m_\sigma}\,\Xi_{\sigma,\min}.
\]
Then
\[
\begin{aligned}
  C_{A\sigma}
  &=
  \int_{\Gamma_{\rm int}}d\nu_\Gamma\,
  \psi_\sigma
  \gamma_\Gamma^{-1}
  (\theta_\Gamma^A,\Theta_{{\rm img},\sigma})\\
  &=
  L_2\sqrt{m_\sigma}\,
  {\cal G}_{AB}^{(\psi)}c_\sigma^B
  =
  L_2\sqrt{m_\sigma}(g_\sigma)_A .
\end{aligned}
\]
Thus
\[
  M_{11}=m_\sigma,
  \qquad
  {\cal D}_{AB}
  =
  \frac{C_{A\sigma}C_{B\sigma}}{m_\sigma L_2^2}
  =
  (g_\sigma)_A(g_\sigma)_B .
\]
The source-profile norm is
\[
  \frac{\|\Xi_{\sigma,\min}\|_\Gamma^2}{L_2^2}
  =
  g_\sigma^T{\cal G}^{(\psi)-1}g_\sigma .
\]
In an orthonormal weighted current frame this gives
\[
\begin{array}{c|c}
{\rm branch}&\|\Xi_{\sigma,\min}\|_\Gamma^2/L_2^2\\ \hline
+&54.478818707582\ldots\\
-&52.600294175197\ldots .
\end{array}
\]
The image row is therefore reduced to the local production of
\(\psi_\sigma\), \(m_\sigma\), and \({\cal G}^{(\psi)}\).

If a Schur slope came from an effective weak period,
\[
  \rho_\sigma=2Q_{\rm eff}^{(\sigma)}+1 .
\]
The two branches give
\[
  Q_{\rm eff}^{(+)}
  =
  9.333551007611\ldots,
  \qquad
  Q_{\rm eff}^{(-)}
  =
  5.818646994052\ldots .
\]
Both values lie outside the integral weak-period lattice.  A period relabel is
therefore excluded; the viable realization rows are Green-source coframe
deformation or projection inside the localized image sector.
% CONTROL: The scale row uses C_W to separate the Weinberg length constant
% from the Hodge entry C in the local A,B,C notation.

\subsection{Boundary Operator Package}

Collect the local realization data as
\[
  {\cal B}_\Gamma
  =
  \left(
  d\nu_\Gamma,\,
  \theta_\Gamma^a,\,
  \theta_\Gamma^Y,\,
  \eta_0,\,
  \Delta_{\Gamma,j},\,
  \mu_\Gamma^2
  \right).
\]
The branch-operator realization is then governed by the boundary system
\[
\begin{array}{ll}
{\rm current:}&
\displaystyle
\int_{\Gamma_{\rm EW,int}}d\nu_\Gamma\,
\gamma_\Gamma^{-1}(\theta_\Gamma^a,\theta_\Gamma^b)
=L_{2,\Gamma}^2\delta^{ab},
\\[3mm]
{\rm neutral:}&
\displaystyle
L_Q^2=(1,1){\cal L}_\Gamma\binom{1}{1},
\qquad
Q\eta_0=0,
\\[3mm]
{\rm endpoint:}&
\displaystyle
Z_{\mathcal E,1}
=
\int_{\Gamma_{\rm EW,int}}d\nu_\Gamma\,|\eta_0|^2
\|\widehat\eta_0\odot\widehat\eta_0\|^2
=Z_0I_0,
\\[3mm]
{\rm complement:}&
\displaystyle
\ker\Delta_{\Gamma,j}={\rm Im}\,A_j,
\qquad
M_{\perp,j}^2>0,
\\[3mm]
{\rm image:}&
\displaystyle
{\cal D}_{AB}
=
L_2^{-2}
\langle{\cal J}_A,{\cal M}_{\rm img}^{-1}{\cal J}_B\rangle_\Gamma,
\qquad A,B=3,Y,
\\[1mm]
&
\displaystyle
\Delta_{\rm rank}=0,\qquad
({\cal D}_{33},{\cal D}_{3Y}/{\cal D}_{33})
\in
\{(\delta_+,\rho_+),(\delta_-,\rho_-)\},
\\[3mm]
{\rm auxiliary:}&
\displaystyle
S_{\rm aux}^{\Gamma}
\leadsto
\mu_\Gamma^2
\begin{pmatrix}
0&A_j^\dagger\\
A_j&-A_jA_j^\dagger
\end{pmatrix}.
\end{array}
\]
A D10 realization that hosts the branch operator supplies the displayed
boundary system.  The absolute electromagnetic normalization, the neutral
mass, and the Schur coefficient are outputs once \({\cal B}_\Gamma\) is
computed.
% CONTROL: Candidate-filter audits remain in codexVersion-proof-gates.md,
% project-status-gates-audit.md, and progressReport10.tex.

\begin{proposition}[Boundary-current closure constraint]
\label{prop:d10currentunit}
A D10 IIA realization hosting the branch operator with unit current metric
supplies a localized weak-current coframe satisfying
\[
  \int_{\Gamma_{\rm EW,int}}d\nu_\Gamma\,
  \gamma_\Gamma^{-1}(\theta_\Gamma^a,\theta_\Gamma^b)
  =
  L_{2,\Gamma}^2\delta^{ab}.
\]
When this closure equation holds, define
\[
  e_\Gamma^a=\frac{\theta_\Gamma^a}{L_{2,\Gamma}}.
\]
Then
\[
  \langle e_\Gamma^a,e_\Gamma^b\rangle_{\rm cur}=\delta^{ab},
  \qquad
  \kappa_{\rm cur}=1,
\]
and
\[
  A_j=\sum_aT_a^{(j)}\otimes e_\Gamma^a,
  \qquad
  A_j^\dagger A_j=j(j+1)\bm1_{V_j}.
\]
\end{proposition}

\begin{proof}
The displayed norm equation makes \(e_\Gamma^a\) orthonormal in the retained
current Hilbert space.  Substitution into the current operator leaves the
finite \(\SU(2)\) Casimir \(\sum_aT_a^{(j)}T_a^{(j)}=j(j+1)\bm1\).
\end{proof}

% CONTROL: Derive theta_Gamma^a, dnu_Gamma, and L_2,Gamma from the localized
% D10 boundary action in codexVersion-proof-gates.md and Report 10.

\section{Product-To-Flag Interface}

The source-backed D10 base is \(\CP^2\times\CP^1\).  The active refinement is
the flag
\[
  F_{1,2}(\C^3)=\SU(3)/T^2=P(Q)\longrightarrow\CP^2,
\]
where \(Q\) is the tautological quotient bundle.  With
\[
  h=c_1(\cO_{\CP^2}(1)),
  \qquad
  \eta=c_1(\cO_{P(Q)}(1)),
\]
the projective-bundle relation is
\[
  \eta^2-h\eta+h^2=0.
\]
For a retained primitive class
\[
  \zeta=P h+Q\eta,
\]
the cubic is
\[
  \zeta^3=3PQ(P+Q)h^2\eta.
\]
This gives a D10 flag-ring normalization candidate for the boundary/current
sector.

The KLT flag metric uses
\[
  J=-ae^{12}+be^{34}-ce^{56},
  \qquad a,b,c>0,
\]
with
\[
  d^2=abc.
\]
The \(\CP^2\)-base-symmetric line is
\[
  a=b=A,\qquad c=\lambda A.
\]
On this line,
\[
  W_2^-
  =
  -\frac{2iA^2(1-\lambda)}{3d}
  \left(e^{12}-e^{34}-2\lambda e^{56}\right),
  \qquad d^2=\lambda A^3.
\]
The KLT admissibility expression reduces to
\[
  \frac{-5\lambda^2+12\lambda-4}{8A\lambda},
  \qquad
  \frac25\leq\lambda\leq2.
\]
The D10 stationary value \(\lambda_*\) must lie in this interval.

% CONTROL: The product-to-flag transition is kept as a proof gate in the
% auxiliary file. The rendered text records the algebraic interface.

\section{\texorpdfstring{Finite Direction, Branched Base, And \(S^3\) Fiber}{Finite Direction, Branched Base, And S3 Fiber}}

The finite-direction reading places the neutral \(O(1)\) field as a
two-sheet connection.  For a finite field \(\Phi_0\),
\[
  \ell_H=\frac{1}{\|\Phi_0\|}=\frac{\sqrt2}{v_{\rm EW}}.
\]
This supplies a language for the scalar Hessian and the electroweak scale.
The same \(Q\)-neutral \(O(1)\) component protects the photon and provides the
branch \(j=\half\) sector.

The branched-projective clue supplies the four-dimensional base replacement
\[
  \CP^2/\hbox{complex conjugation}=S^4.
\]
The quotient map \(\CP^2\to S^4\) is a two-fold branched covering with branch
locus \(\mathbb{RP}^2\).  Away from \(\mathbb{RP}^2\) it has two local sheets;
along \(\mathbb{RP}^2\) the sheets meet.  This gives a geometric setting for
double-cover memory, branch/interface data, and spin or twisted sectors.

The corresponding seven-dimensional Hopf model is
\[
  S^3\longrightarrow S^7\longrightarrow S^4,
  \qquad 3+4=7.
\]
Replacing the spherical base by the projective base asks for an
\(S^3\)-type parent over \(\CP^2\):
\[
  S^3\hbox{-type fiber}\longrightarrow K_7^{\rm br}\longrightarrow\CP^2.
\]
A precise geometric test object is the pullback
\[
  K_7^{\rm br}=\CP^2\times_{S^4}S^7.
\]
Over the smooth locus of the branched map, \(K_7^{\rm br}\) gives an
\(S^3\)-fibered seven-dimensional space over \(\CP^2\); over
\(\mathbb{RP}^2\) it carries the special branch/interface data.

The D10 compactification keeps the six-dimensional projective object
\[
  K_6=\CP^2\times\CP^1
\]
with retained \(C_1\), \(F_2\), current, boundary, and branch data.  Local
\(\SU(2)\) double-cover memory can be carried by current, flat, spin, branch,
or boundary sectors inside this D10 object.  The flag target
\[
  \CP^1\longrightarrow F_{1,2}(\C^3)=\SU(3)/T^2\longrightarrow\CP^2
\]
ties the weak projective fiber to the projective base.

The D9 projection keeps the \(\CP^2\) base and the electromagnetic stabilizer:
\[
  K_{5,{\rm eff}}=\CP^2\times S^1_Q.
\]
The weak projective direction is projected by the neutral \(O(1)\) section,
\[
  \eta_0\in H^0(\CP^1,\cO(1)),
  \qquad Q\eta_0=0.
\]
The ladder is therefore
\[
\begin{array}{ccccc}
S^3 &\longrightarrow& S^7 &\longrightarrow& S^4\\[1mm]
S^3\hbox{-type} &\longrightarrow& K_7^{\rm br} &\longrightarrow& \CP^2\\[1mm]
\CP^1 &\longrightarrow& K_6 &\longrightarrow& \CP^2\\
&& \downarrow Q &&\\[-1mm]
S^1_Q &\longrightarrow& K_5 && \CP^2 .
\end{array}
\]
The branched \(\CP^2\to S^4\) map is the comparison between the spherical
Hopf base and the projective base; the \(S^3\)-type fiber supplies the
seven-dimensional parent.

The pure Hopf branch gives a useful control endpoint.  Keeping \(S^7\) as the
parent and projecting to the electromagnetic stabilizer leads to
\[
  S^4\times S^1_Q,
  \qquad
  \mathrm{Spin}(5)\times\U(1)_Q\simeq \mathrm{Sp}(2)\times\U(1)_Q
\]
at the double-cover level, or \(\SO(5)\times\U(1)_Q\) at the adjoint level.
The projective lift
\[
  \CP^2\simeq \SU(3)/S(\U(2)\times\U(1))
\]
supplies the eight-generator \(SU(3)\) action and the \(\Z_3\) center entering
\[
  (\SU(3)_c\times\U(1)_Q)/\Z_3.
\]
Thus the \(S^7/S^4\) chain is the quaternionic control branch.  The \(\CP^2\)
lift is the color branch.

% CONTROL: The pullback K7^br=CP2 x_{S4} S7 and the branch-sector
% representation test are recorded in branched-cp2-s4-s3-ladder.md.

\section{Numerical Consequences}

At \(c=1\),
\[
  \Xp(3/4)=\frac18(-3+\sqrt{57})=0.568729\ldots,
  \qquad
  \Xp(2)=-1+\sqrt3=0.732051\ldots .
\]
Hence
\[
  \sin^2\theta_W
  =
  1-\frac{\Xp(3/4)}{\Xp(2)}
  =
  0.223101\ldots,
\]
and
\[
  \tan^2\theta_{\rm pole}=0.2871691363429836.
\]
The same definitions give the custodial tree-level identity.  Since
\[
  \cos^2\theta_W=\frac{\Xp(3/4)}{\Xp(2)},
  \qquad
  \frac{M_W^2}{M_Z^2}=\frac{\Xp(3/4)}{\Xp(2)},
\]
one has
\[
  \rho_{\rm cust}
  =
  \frac{M_W^2}{M_Z^2\cos^2\theta_W}
  =1.
\]
The absolute electromagnetic normalization is assigned to the D10 boundary
length matrix.  With
\[
  L_Q^2=(1,1){\cal L}_\Gamma\binom{1}{1},
  \qquad
  \alpha_\Gamma=\frac{C^2}{4\pi L_Q^2},
\]
alpha is obtained from the same localized current calculation that gives
\(M_Z\) and the photon null direction.
% CONTROL: Tuple-family scans are density diagnostics and rejected as
% construction evidence.  The active calculation is the D10 Gamma_EW action.

\section{Conclusion}

The D10 boundary-current route has a compact operator core.  The
Borel--Weil weak sectors supply \(j=0,\half,1\).  The finite current zero modes
give \(A_j^\dagger A_j=J\).  A single Stueckelberg boundary Schur completion
gives the block \(\cA(J)\).  The \(Q\)-boundary sector protects the photon and
descends the global form to \((\SU(3)_c\times\U(1)_Q)/\Z_3\).  The D8-like
boundary \(\Gamma_{\rm EW}\) also supplies a natural location for the Romans
jump, the RR/current vertical length, and the auxiliary field producing the
negative diagonal.

The finite electroweak Connes layer supplies the scalar-vacuum dictionary:
\(\ell_H=\sqrt2/v_{\rm EW}\), \(Q\eta_0=0\), the standard \(W/Z/\gamma\)
Hessian, and the local \((h,\sigma)\) extension.  This fixes the finite sheet
scale.  A stationary D10 realization is specified by the common compactification
potential
\[
  V_{\rm eff}(\rho,\lambda,v_F,\theta_H,\sigma,\ldots)
\]
with a stable or metastable D10 threshold window above the D9 endpoint.  The
KO6 layer supplies the finite-triple tests; fermion data are assigned to the
full Standard-Model lift.

The resulting conditional manuscript is concentrated on one boundary/current
sector.  Its realization data are the localized vertical current, the unit
current metric, the Schur auxiliary coefficient and complement gap, and the
finite/KK scalar Hessian on the D10 interface.

\begingroup
\raggedright
\begin{thebibliography}{99}

\bibitem{Weinberg1983}
S. Weinberg,
``Charges from extra dimensions,''
Phys. Lett. B 125 (1983) 265.

\bibitem{KoerberLustTsimpis2008}
P. Koerber, D. Lust, and D. Tsimpis,
``Type IIA \(\AdS_4\) compactifications on cosets, interpolations and domain
walls,''
JHEP 07 (2008) 017, arXiv:0804.0614.

\bibitem{HertzbergKachruTaylorTegmark2007}
M. P. Hertzberg, S. Kachru, W. Taylor, and M. Tegmark,
``Inflationary constraints on type IIA string theory,''
arXiv:0711.2512.

\bibitem{BergshoeffKalloshOrtinRoestVanProeyen2001}
E. Bergshoeff, R. Kallosh, T. Ortin, D. Roest, and A. Van Proeyen,
``New formulations of D=10 supersymmetry and D8-O8 domain walls,''
arXiv:hep-th/0103233.

\bibitem{GreenHarveyMoore1996}
M. B. Green, J. A. Harvey, and G. Moore,
``I-brane inflow and anomalous couplings on D-branes,''
arXiv:hep-th/9605033.

\bibitem{MinasianMoore1997}
R. Minasian and G. Moore,
``K-theory and Ramond-Ramond charge,''
arXiv:hep-th/9710230.

\bibitem{FreedWitten1999}
D. S. Freed and E. Witten,
``Anomalies in string theory with D-branes,''
arXiv:hep-th/9907189.

\bibitem{HenryLabordereJulia2001}
P. Henry-Labordere and B. Julia,
``Gravitational couplings of orientifold-planes,''
arXiv:hep-th/0112065.

\bibitem{Payen2007}
R. Payen,
``An action for Chan-Paton factors,''
arXiv:0708.0888.

\bibitem{BouwknegtSchoutensLudwig1995}
P. Bouwknegt, A. W. W. Ludwig, and K. Schoutens,
``Spinon basis for higher level SU(2) WZW models,''
arXiv:hep-th/9412108.

\bibitem{Manton1979}
N. S. Manton,
``A new six-dimensional approach to the Weinberg--Salam model,''
Nucl. Phys. B 158 (1979) 141.

\bibitem{Hosotani1983}
Y. Hosotani,
``Dynamical mass generation by compact extra dimensions,''
Phys. Lett. B 126 (1983) 309.

\bibitem{ConnesLott1990}
A. Connes and J. Lott,
``Particle models and noncommutative geometry,''
Nucl. Phys. B Proc. Suppl. 18B (1990) 29.

\bibitem{Connes1989}
A. Connes,
``Essay on physics and non-commutative geometry,''
IHES preprint M/89/69.

\bibitem{ChamseddineConnes2012}
A. H. Chamseddine and A. Connes,
``Resilience of the spectral standard model,''
JHEP 09 (2012) 104, arXiv:1208.1030.

\bibitem{ChamseddineConnesVanSuijlekom2013a}
A. H. Chamseddine, A. Connes, and W. D. van Suijlekom,
``Inner fluctuations in noncommutative geometry without the first order
condition,''
J. Geom. Phys. 73 (2013) 222, arXiv:1304.7583.

\bibitem{ChamseddineConnesVanSuijlekom2013b}
A. H. Chamseddine, A. Connes, and W. D. van Suijlekom,
``Beyond the spectral standard model: emergence of Pati--Salam unification,''
JHEP 11 (2013) 132, arXiv:1304.8050.

\bibitem{Martinetti2015}
P. Martinetti,
``Twisted spectral geometry for the standard model,''
arXiv:1503.07548.

\bibitem{FilaciMartinettiPesco2020}
M. Filaci, P. Martinetti, and S. Pesco,
``Twisted Standard Model in noncommutative geometry I: the field content,''
arXiv:2008.01629.

\bibitem{MartinettiWulkenhaar2002}
P. Martinetti and R. Wulkenhaar,
``Discrete Kaluza--Klein from scalar fluctuations in noncommutative
geometry,''
J. Math. Phys. 43 (2002) 182, arXiv:hep-th/0104108.

\bibitem{ChamseddineConnesMarcolli2007}
A. H. Chamseddine, A. Connes, and M. Marcolli,
``Gravity and the standard model with neutrino mixing,''
Adv. Theor. Math. Phys. 11 (2007) 991, arXiv:hep-th/0610241.

\bibitem{Stephan2009}
C. A. Stephan,
``Krajewski diagrams and the standard model,''
J. Math. Phys. 50 (2009) 043515, arXiv:0809.5137.

\bibitem{DAndreaDabrowski2016}
F. D'Andrea and L. Dabrowski,
``The standard model in noncommutative geometry and Morita equivalence,''
J. Noncommut. Geom. 10 (2016) 551, arXiv:1501.00156.

\bibitem{DabrowskiSitarz2019}
L. Dabrowski and A. Sitarz,
``Fermion masses, mass-mixing and the almost commutative geometry of the
standard model,''
JHEP 02 (2019) 068, arXiv:1806.07282.

\bibitem{DabrowskiMukhopadhyayPozar2025}
L. Dabrowski, S. Mukhopadhyay, and F. Pozar,
``Spectral torsion of the internal noncommutative geometry of the standard
model,''
arXiv:2511.08159.

\bibitem{Hillman2017}
J. A. Hillman,
``An explicit formula for a branched covering from \(CP^2\) to \(S^4\),''
arXiv:1705.05038.

\bibitem{AtiyahMantonSchroers2011}
M. F. Atiyah, N. S. Manton, and B. J. Schroers,
``Geometric models of matter,''
Proc. R. Soc. A 468 (2012) 1252, arXiv:1108.5151.

\end{thebibliography}
\endgroup

\end{document}
